% =====================================================================
%  LIBRO DE ANÁLISIS MATEMÁTICO — ESTRUCTURA MODULAR
%  Compilar con: xelatex -> biber -> xelatex -> xelatex
%  (makeindex se ejecuta automáticamente)
% =====================================================================
\documentclass[11pt,twoside,openright]{book}

% --- Preámbulo modular ---
% ---------------------------------------------------------------------
%  1. GEOMETRÍA DE PÁGINA — márgenes clásicos de libro impreso
% ---------------------------------------------------------------------
\usepackage[
  paperwidth=21cm,
  paperheight=29.7cm,
  inner=3cm,
  outer=2.5cm,
  top=2.8cm,
  bottom=3cm,
  headsep=18pt,
  footskip=32pt
]{geometry}

% ---------------------------------------------------------------------
%  2. TIPOGRAFÍA — Latin Modern (texto) + Latin Modern Math
% ---------------------------------------------------------------------
\usepackage{fontspec}

\usepackage{polyglossia}
\setdefaultlanguage{spanish}

\usepackage{microtype}          % justificación y espaciado óptico finos
\frenchspacing                  % un solo espacio tras punto, uso editorial

% Separación vertical cómoda en ecuaciones desplegadas
\usepackage{amsmath}
\usepackage{mathtools}

\usepackage{unicode-math}
\setmainfont{Latin Modern Roman}
\setsansfont{Latin Modern Sans}
\setmonofont{Latin Modern Mono}[Scale=0.92]
\setmathfont{Latin Modern Math}

\setlength{\abovedisplayskip}{12pt plus 3pt minus 6pt}
\setlength{\belowdisplayskip}{12pt plus 3pt minus 6pt}
\setlength{\abovedisplayshortskip}{6pt plus 3pt}
\setlength{\belowdisplayshortskip}{6pt plus 3pt}
\allowdisplaybreaks

% ---------------------------------------------------------------------
%  3. COLOR — únicamente negro
% ---------------------------------------------------------------------
\usepackage{xcolor}

% ---------------------------------------------------------------------
%  4. CAPÍTULOS Y SECCIONES — diseño sobrio, sin adornos
% ---------------------------------------------------------------------
\usepackage{titlesec}
\usepackage{titletoc}

% Capítulo: centrado, tipografía sobria con número destacado
\titleformat{\chapter}[display]
  {\centering\normalfont}
  {\large\MakeUppercase{\chaptertitlename}\\[6pt]\fontsize{28}{34}\selectfont\thechapter}
  {20pt}
  {\LARGE\itshape\centering}
\titlespacing*{\chapter}{0pt}{1.5in}{0.8in}

\titleformat{\section}
  {\normalfont\bfseries}{\thesection}{0.8em}{}
\titlespacing*{\section}{0pt}{22pt plus 4pt minus 2pt}{10pt}

\titleformat{\subsection}
  {\normalfont\bfseries\itshape}{\thesubsection}{0.8em}{}
\titlespacing*{\subsection}{0pt}{16pt plus 3pt minus 2pt}{8pt}

\titleformat{\subsubsection}
  {\normalfont\itshape}{\thesubsubsection}{0.8em}{}
\titlespacing*{\subsubsection}{0pt}{12pt plus 2pt}{6pt}

% Espaciado de párrafo clásico: sangría, sin salto entre párrafos
\setlength{\parindent}{1.4em}
\setlength{\parskip}{0pt}

% ---------------------------------------------------------------------
%  5. ENCABEZADOS Y PIES DE PÁGINA — sobrios, sin reglas horizontales
% ---------------------------------------------------------------------
\usepackage{fancyhdr}
\setlength{\headheight}{13pt}
\pagestyle{fancy}
\fancyhf{}
\renewcommand{\headrulewidth}{0pt}
\renewcommand{\footrulewidth}{0pt}
\fancyhead[LE]{\small\itshape\nouppercase{\leftmark}}
\fancyhead[RO]{\small\itshape\nouppercase{\rightmark}}
\fancyhead[RE,LO]{}
\fancyfoot[C]{\small\thepage}

\renewcommand{\chaptermark}[1]{\markboth{\thechapter.\ #1}{}}
\renewcommand{\sectionmark}[1]{\markright{\thesection.\ #1}}

\fancypagestyle{plain}{%
  \fancyhf{}
  \fancyfoot[C]{\small\thepage}
  \renewcommand{\headrulewidth}{0pt}
}

% ---------------------------------------------------------------------
%  PÁGINAS EN BLANCO ENTRE CAPÍTULOS — totalmente vacías
% ---------------------------------------------------------------------
\makeatletter
\def\cleardoublepage{\clearpage\if@twoside \ifodd\c@page\else
\hbox{}\thispagestyle{empty}\newpage\if@twocolumn\hbox{}\newpage\fi\fi\fi}
\makeatother

% ---------------------------------------------------------------------
%  6. ÍNDICE GENERAL — sobrio, sin puntos decorativos gruesos
% ---------------------------------------------------------------------
\usepackage{tocloft}
\renewcommand{\cftchapfont}{\normalfont}
\renewcommand{\cftchappagefont}{\normalfont}
\renewcommand{\cftsecfont}{\normalfont}
\renewcommand{\cftchapleader}{\cftdotfill{\cftdotsep}}
\setlength{\cftbeforechapskip}{10pt}

% ---------------------------------------------------------------------
%  7. TIKZ — dibujo de diagramas y figuras matemáticas
% ---------------------------------------------------------------------
\usepackage{tikz}
\usetikzlibrary{arrows.meta, calc, positioning, decorations.pathmorphing, decorations.pathreplacing}

% ---------------------------------------------------------------------
%  HERRAMIENTAS MATEMÁTICAS Y ENTORNOS
% ---------------------------------------------------------------------
\usepackage{bm}
\usepackage{amsthm}

% --- estilo "plain": encabezado en negrita, cuerpo en cursiva ---
\newtheoremstyle{plainbook}
  {12pt plus 2pt minus 2pt}   % espacio antes
  {12pt plus 2pt minus 2pt}   % espacio después
  {\itshape}                  % cuerpo
  {}                          % sangría
  {\bfseries}                 % encabezado
  {.}                         % puntuación tras el encabezado
  {0.6em}                     % espacio tras el encabezado
  {}

% --- estilo "definición": encabezado en negrita, cuerpo en romana ---
\newtheoremstyle{defbook}
  {12pt plus 2pt minus 2pt}
  {12pt plus 2pt minus 2pt}
  {\normalfont}
  {}
  {\bfseries}
  {.}
  {0.6em}
  {}

% --- estilo "observación/nota": encabezado en cursiva, cuerpo romano ---
\newtheoremstyle{notebook}
  {10pt plus 2pt minus 2pt}
  {10pt plus 2pt minus 2pt}
  {\normalfont}
  {}
  {\itshape}
  {.}
  {0.6em}
  {}

% Cada entorno recibe su propio contador (para que cleveref use el nombre
% correcto de cada tipo), pero todos se alían al contador de "theorem".
\theoremstyle{plainbook}
\newtheorem{theorem}{Teorema}[chapter]

\theoremstyle{plainbook}
\newtheorem{lemma}{Lema}[chapter]
\newtheorem{proposition}{Proposición}[chapter]
\newtheorem{corollary}{Corolario}[chapter]
\newtheorem{conjecture}{Conjetura}[chapter]

\theoremstyle{defbook}
\newtheorem{axiom}{Axioma}[chapter]
\newtheorem{definition}{Definición}[chapter]
\newtheorem{example}{Ejemplo}[chapter]
\newtheorem{exercise}{Ejercicio}[chapter]
\newtheorem{problem}{Problema}[chapter]

\theoremstyle{notebook}
\newtheorem{remark}{Observación}[chapter]
\newtheorem{notation}{Notación}[chapter]

\makeatletter
\let\c@lemma\c@theorem
\let\c@proposition\c@theorem
\let\c@corollary\c@theorem
\let\c@conjecture\c@theorem
\let\c@axiom\c@theorem
\let\c@definition\c@theorem
\let\c@example\c@theorem
\let\c@exercise\c@theorem
\let\c@problem\c@theorem
\let\c@remark\c@theorem
\let\c@notation\c@theorem
\makeatother

% Demostración con marca cuadrada discreta
\renewcommand{\qedsymbol}{$\mdlgblksquare$}
\renewcommand{\proofname}{Demostración}

% ---------------------------------------------------------------------
%  HIPERVÍNCULOS — discretos, sin recuadros
% ---------------------------------------------------------------------
\usepackage[
  colorlinks=true,
  linkcolor=blue,
  citecolor=blue,
  urlcolor=blue,
  pdfusetitle,
  bookmarksnumbered=true,
  bookmarksopen=true
]{hyperref}
\usepackage{bookmark}

% ---------------------------------------------------------------------
%  REFERENCIAS CRUZADAS Y BIBLIOGRAFÍA
% ---------------------------------------------------------------------
\usepackage[capitalise,spanish]{cleveref}
\crefname{theorem}{teorema}{teoremas}
\Crefname{theorem}{Teorema}{Teoremas}
\crefname{lemma}{lema}{lemas}
\Crefname{lemma}{Lema}{Lemas}
\crefname{proposition}{proposición}{proposiciones}
\Crefname{proposition}{Proposición}{Proposiciones}
\crefname{corollary}{corolario}{corolarios}
\Crefname{corollary}{Corolario}{Corolarios}
\crefname{definition}{definición}{definiciones}
\Crefname{definition}{Definición}{Definiciones}
\crefname{example}{ejemplo}{ejemplos}
\Crefname{example}{Ejemplo}{Ejemplos}
\crefname{remark}{observación}{observaciones}
\Crefname{remark}{Observación}{Observaciones}
\crefname{equation}{ecuación}{ecuaciones}
\Crefname{equation}{Ecuación}{Ecuaciones}
\crefname{chapter}{capítulo}{capítulos}
\Crefname{chapter}{Capítulo}{Capítulos}
\crefname{section}{sección}{secciones}
\Crefname{section}{Sección}{Secciones}

% Bibliografía con biblatex (motor moderno, mejor para matemáticas)
\usepackage[backend=biber, style=numeric, sorting=ynt, hyperref=true]{biblatex}


% --- Índice de palabras ---
\usepackage{imakeidx}
\makeindex

% --- Metadatos del libro ---
\title{Análisis Matemático\\[-4pt] \large Del rigor foundational a las estructuras del análisis moderno}
\author{}
\date{}

% --- Bibliografía ---
\addbibresource{bibliografia.bib}

% --- Metadatos PDF ---
\hypersetup{
  pdfinfo={
    Title={Análisis Matemático},
    Author={},
    Subject={Análisis matemático, espacios métricos, topología, convergencia, completitud},
    Keywords={análisis matemático, espacios métricos, topología, convergencia, completitud, compacidad, teoremas estructurales},
    Creator={XeLaTeX con biblatex},
    Producer={MiKTeX}
  }
}

% =====================================================================
\begin{document}

% --- Front matter: portada, prólogo, índices (páginas en números romanos) ---
% Cada elemento ocupa siempre una página propia.
\frontmatter

% Portada
% -------------------- PORTADA INTERIOR (sobria) ---------------------
\begingroup
\thispagestyle{empty}
\begin{center}
\vspace*{1.6in}
{\Large\itshape Fundamentos de\\[2pt] Análisis Matemático\par}
\vspace{0.6in}
{\normalsize M.H. Alberto\par}
\vfill
{\small Editorial Universitaria\par}
\vspace{0.3in}
\end{center}
\endgroup
\clearpage

\thispagestyle{empty}\mbox{}\clearpage

\clearpage

% Dedicatoria
\thispagestyle{empty}
\vspace*{2in}
\begin{flushright}
\itshape
A quienes encuentran belleza\\
en la precisión de una demostración.
\end{flushright}
\clearpage

\clearpage

% Prólogo
\chapter*{Prólogo}
\markboth{Prólogo}{Prólogo}
\addcontentsline{toc}{chapter}{Prólogo}

Las matemáticas nacieron como una herramienta para describir el mundo. Durante siglos, la geometría griega, el álgebra y posteriormente el cálculo permitieron resolver problemas de enorme complejidad con resultados sorprendentemente precisos. Sin embargo, ese éxito ocultaba una dificultad fundamental: muchos de los argumentos utilizados carecían de una justificación rigurosa.

El cálculo infinitesimal desarrollado por Newton y Leibniz revolucionó la ciencia, pero se apoyaba en conceptos como los \emph{infinitésimos}, los \emph{límites} y las \emph{series infinitas}, cuya naturaleza no estaba claramente definida. Los matemáticos obtenían resultados correctos, aunque con frecuencia los demostraban mediante argumentos intuitivos o manipulaciones que hoy consideraríamos insuficientes.

A lo largo del siglo~XIX comenzaron a aparecer contraejemplos que mostraban que la intuición no siempre era un guía confiable. Existían funciones continuas no derivables, series que cambiaban de valor al reordenar sus términos, sucesiones cuyo comportamiento desafiaba las ideas tradicionales y paradojas derivadas de un uso poco preciso del infinito. Estos fenómenos revelaron que las matemáticas necesitaban una base mucho más sólida.

La respuesta fue el nacimiento del \emph{análisis matemático moderno}. Matemáticos como Cauchy, Bolzano, Weierstrass, Dedekind, Cantor, Borel, Lebesgue, Fréchet, Hausdorff, Banach y muchos otros desarrollaron un nuevo lenguaje basado en definiciones precisas y demostraciones rigurosas. Conceptos como la \emph{completitud}, la \emph{continuidad}, la \emph{compacidad}, la \emph{convergencia} o la \emph{topología} dejaron de ser ideas intuitivas para convertirse en estructuras matemáticas perfectamente definidas.

A partir de ese momento, los grandes teoremas del análisis adquirieron un papel fundamental. Cada uno de ellos establece las condiciones exactas bajo las cuales un procedimiento matemático es válido: cuándo una sucesión converge, cuándo una función alcanza un máximo, cuándo una familia de funciones posee una subsucesión convergente o cuándo una función puede aproximarse arbitrariamente bien por polinomios. Los teoremas no son únicamente resultados aislados; son las garantías que permiten trabajar con el infinito de manera segura.

Este libro sigue ese recorrido histórico y conceptual. Comienza construyendo los fundamentos lógicos y axiomáticos sobre los que descansa el análisis, desarrolla la teoría de espacios métricos y las propiedades esenciales de las funciones, y culmina con algunos de los resultados estructurales más importantes del análisis moderno. Finalmente, muestra cómo estas ideas constituyen la base de disciplinas como la teoría de la medida, el análisis funcional, la teoría espectral, las ecuaciones diferenciales, la topología y la probabilidad moderna.

Más que aprender técnicas de cálculo, el objetivo de esta obra es comprender \emph{por qué las matemáticas funcionan}, cuáles son los principios que garantizan la validez de sus argumentos y cómo el rigor permitió transformar el cálculo intuitivo en una de las teorías más profundas y sólidas de toda la matemática.

\vspace{1em}
\begin{flushright}
\textit{M.H. Alberto}
\end{flushright}

\clearpage

\tableofcontents
\clearpage

\listoffigures
\clearpage

\listoftables

\mainmatter

% --- Capítulos ---
\chapter{Espacios métricos}
\label{ch:espacios-metricos}
\index{Espacio métrico}\index{Métrica}\index{Distancia}\index{Desigualdad de Cauchy--Schwarz}\index{Desigualdad de Hölder}\index{Desigualdad de Minkowski}\index{Métricas equivalentes}\index{Subespacio métrico}\index{Bola abierta}\index{Diámetro}\index{Entorno}\index{Clausura}\index{Interior}\index{Frontera}\index{Punto de acumulación}\index{Punto aislado}\index{Conjunto denso}\index{Separabilidad}\index{Axiomas de separación}\index{Isometría}\index{Conjunto de Cantor}

\section{Introducción}

Con los números reales ya construidos, el siguiente paso natural en la evolución del rigor es abstraer la noción de distancia. El análisis matemático no estudia únicamente funciones sobre $\mathbb{R}$; sus objetos fundamentales son los espacios métricos, introducidos por Maurice Fréchet en su tesis de 1906 como una generalización que unifica el tratamiento de sucesiones, funciones, conjuntos y transformaciones.

\subsection{Definición de espacio métrico}

\begin{definition}
\label{def:metrica}
Sea $X$ un conjunto no vacío. Una \emph{métrica} sobre $X$ es una función $d\colon X\times X\to\mathbb{R}$ que satisface, para todos $x,y,z\in X$:
\begin{enumerate}
  \item[(i)] \emph{No negatividad y separación}: $d(x,y)\geq 0$, y $d(x,y)=0$ si y solo si $x=y$;
  \item[(ii)] \emph{Simetría}: $d(x,y)=d(y,x)$;
  \item[(iii)] \emph{Desigualdad triangular}: $d(x,z)\leq d(x,y)+d(y,z)$.
\end{enumerate}
Al par $(X,d)$ se le llama \emph{espacio métrico}.
\end{definition}

La condición (iii) es la que confiere a la métrica su carácter de medida de proximidad compatible con una estructura de límite. Intuitivamente, dice que ir directamente de $x$ a $z$ nunca puede ser más largo que pasar por un intermediario $y$.

\subsection{Ejemplos de métricas}

\begin{example}
\label{ej:euclidea}
El espacio $\mathbb{R}^n$ dotado de la distancia euclidiana
\begin{equation}\label{eq:euclidea}
  d_2(x,y)=\biggl(\sum_{i=1}^{n}(x_i-y_i)^2\biggr)^{1/2}
\end{equation}
es un espacio métrico. La desigualdad triangular se sigue de la desigualdad de Cauchy--Schwarz.
\end{example}

\begin{example}
\label{ej:discreta}
Sobre cualquier conjunto no vacío $X$, la \emph{métrica discreta}
\[
  d(x,y)=\begin{cases}0&\text{si }x=y,\\1&\text{si }x\neq y\end{cases}
\]
es una métrica. La desigualdad triangular es inmediata: si $x\neq z$, entonces al menos uno de $d(x,y)$ o $d(y,z)$ es $1$.
\end{example}

\begin{example}
\label{ej:supremo}
Sea $C([a,b])$ el conjunto de las funciones continuas $f\colon[a,b]\to\mathbb{R}$, definidas sobre el intervalo cerrado $[a,b]$. La fórmula
\begin{equation}\label{eq:supremo}
  d(f,g)=\sup_{t\in[a,b]}|f(t)-g(t)|
\end{equation}
define la \emph{métrica del supremo} (también llamada \emph{métrica infinita}, por analogía con $d_\infty$) sobre el espacio $C([a,b])$ de las funciones continuas sobre el intervalo cerrado $[a,b]$.
\end{example}

\subsection{La desigualdad de Cauchy--Schwarz}

Antes de continuar con más ejemplos, demostramos la herramienta básica que justifica la desigualdad triangular de la métrica euclidiana.

\begin{theorem}[Cauchy--Schwarz en $\mathbb{R}^n$]
\label{teo:cauchy-schwarz-cap5}
Para cualesquiera $x,y\in\mathbb{R}^n$,
\[
  \biggl|\sum_{i=1}^{n}x_i y_i\biggr|\leq\biggl(\sum_{i=1}^{n}x_i^2\biggr)^{1/2}\biggl(\sum_{i=1}^{n}y_i^2\biggr)^{1/2}.
\]
La igualdad se da si y solo si $x$ e $y$ son proporcionales.
\end{theorem}
\begin{proof}
Si $y=0$ el resultado es trivial. Supongamos $y\neq 0$ y sea $\lambda=\langle x,y\rangle/\langle y,y\rangle$. Entonces
\[
  0\leq\langle x-\lambda y,x-\lambda y\rangle=\|x\|^2-\frac{\langle x,y\rangle^2}{\|y\|^2},
\]
de donde se deduce el teorema.
\end{proof}

\begin{corollary}
\label{cor:euclidiana-metrica}
La función $d_2(x,y)=(\sum|x_i-y_i|^2)^{1/2}$ satisface la desigualdad triangular, luego es una métrica sobre $\mathbb{R}^n$.
\end{corollary}
\begin{proof}
Escribiendo $z_i=x_i-y_i$ y $w_i=y_i-z_i$ (ajustando índices), se aplica Cauchy--Schwarz para acotar $\|x-z\|^2$.
\end{proof}

\subsection{La desigualdad de Hölder}

La desigualdad de Cauchy--Schwarz admite una generalización natural a toda pareja de exponentes conjugados.

\begin{definition}
\label{def:exponentes-conjugados}
Sean $p,q\in(1,\infty)$ tales que $\frac{1}{p}+\frac{1}{q}=1$. Diremos que $p$ y $q$ son \emph{exponentes conjugados}.
\end{definition}

\begin{theorem}[Desigualdad de Hölder en $\mathbb{R}^n$]
\label{teo:holder-rn}
Para cualesquiera $x,y\in\mathbb{R}^n$ y exponentes conjugados $p,q$,
\[
  \sum_{i=1}^{n}|x_i y_i|\leq\biggl(\sum_{i=1}^{n}|x_i|^p\biggr)^{1/p}\biggl(\sum_{i=1}^{n}|y_i|^q\biggr)^{1/q}.
\]
El caso $p=q=2$ no es otro que la desigualdad de Cauchy--Schwarz (\cref{teo:cauchy-schwarz-cap5}).
\end{theorem}
\begin{proof}
Recordemos la desigualdad de Young: para $a,b\geq 0$ y exponentes conjugados $p,q$,
\[
  ab\leq\frac{a^p}{p}+\frac{b^q}{q}.
\]
Si $A=\bigl(\sum|x_i|^p\bigr)^{1/p}$ y $B=\bigl(\sum|y_i|^q\bigr)^{1/q}$ se anulan, el resultado es trivial. En caso contrario aplicamos Young a $a=|x_i|/A$ y $b=|y_i|/B$:
\[
  \frac{|x_i y_i|}{AB}\leq\frac{|x_i|^p}{pA^p}+\frac{|y_i|^q}{qB^q},
\]
y sumando sobre $i$ obtenemos
\[
  \frac{\sum_{i=1}^{n}|x_i y_i|}{AB}\leq\frac{1}{p}+\frac{1}{q}=1.
\]
\end{proof}

La desigualdad de Hölder es la pieza que permite transferir al caso $p\neq 2$ los argumentos que la métrica euclidiana recibe de Cauchy--Schwarz: con ella se demuestra la desigualdad de Minkowski y, en última instancia, la desigualdad triangular de todas las métricas $\ell^p$.

\subsection{La desigualdad de Minkowski}

\begin{theorem}[Desigualdad de Minkowski en $\mathbb{R}^n$]
\label{teo:minkowski-rn}
Para $1\leq p<\infty$ y cualesquiera $x,y\in\mathbb{R}^n$,
\[
  \biggl(\sum_{i=1}^{n}|x_i+y_i|^p\biggr)^{1/p}\leq\biggl(\sum_{i=1}^{n}|x_i|^p\biggr)^{1/p}+\biggl(\sum_{i=1}^{n}|y_i|^p\biggr)^{1/p}.
\]
\end{theorem}
\begin{proof}
Para $p=1$ es la desigualdad triangular usual. Supongamos $1<p<\infty$ y sea $q$ el conjugado de $p$. Escribamos $S=\bigl(\sum_{i=1}^{n}|x_i+y_i|^p\bigr)^{1/p}$. Entonces
\[
  S^p=\sum_{i=1}^{n}|x_i+y_i|^p=\sum_{i=1}^{n}|x_i+y_i||x_i+y_i|^{p-1}\leq\sum_{i=1}^{n}\bigl(|x_i|+|y_i|\bigr)|x_i+y_i|^{p-1}.
\]
Por la desigualdad de Hölder (\cref{teo:holder-rn}),
\[
  \sum_{i=1}^{n}|x_i||x_i+y_i|^{p-1}\leq\biggl(\sum_{i=1}^{n}|x_i|^p\biggr)^{1/p}\biggl(\sum_{i=1}^{n}|x_i+y_i|^{q(p-1)}\biggr)^{1/q},
\]
y como $q(p-1)=p$, el segundo factor es $S^{p/q}=S^{p-1}$. La estimación análoga vale para el término con $|y_i|$, de modo que
\[
  S^p\leq S^{p-1}\Bigl(\bigl(\sum|x_i|^p\bigr)^{1/p}+\bigl(\sum|y_i|^p\bigr)^{1/p}\Bigr),
\]
y basta dividir entre $S^{p-1}$ (si $S\neq 0$; el caso $S=0$ es trivial).
\end{proof}

\begin{corollary}
\label{cor:lp-metrica-rn}
Para $1\leq p<\infty$,
\[
  d_p(x,y)=\biggl(\sum_{i=1}^{n}|x_i-y_i|^p\biggr)^{1/p}
\]
es una métrica sobre $\mathbb{R}^n$, pues la desigualdad triangular se sigue de aplicar Minkowski a los vectores $x-y$ e $y-z$, cuya suma es $x-z$.
\end{corollary}

\subsection{Las desigualdades en forma integral}

Los espacios de funciones $C([a,b])$ heredan las mismas desigualdades en versión integral, con las sumas sustituidas por la integral de Riemann.

\begin{theorem}[Hölder y Minkowski en forma integral]
\label{teo:desigualdades-integrales}
Sean $f,g\in C([a,b])$ y $1<p<\infty$.
\begin{enumerate}
  \item[(i)] \emph{Hölder integral}: si $q$ es el conjugado de $p$,
  \[
    \int_a^b|f(t)g(t)|\,dt\leq\biggl(\int_a^b|f(t)|^p\,dt\biggr)^{1/p}\biggl(\int_a^b|g(t)|^q\,dt\biggr)^{1/q};
  \]
  \item[(ii)] \emph{Minkowski integral}: para $1\leq p<\infty$,
  \[
    \biggl(\int_a^b|f(t)+g(t)|^p\,dt\biggr)^{1/p}\leq\biggl(\int_a^b|f(t)|^p\,dt\biggr)^{1/p}+\biggl(\int_a^b|g(t)|^p\,dt\biggr)^{1/p}.
  \]
\end{enumerate}
\end{theorem}
\begin{proof}
Para (i) se replica la demostración de la \cref{teo:holder-rn} sustituyendo sumas por integrales y aplicando la desigualdad de Young puntualmente: con
\[
A=\biggl(\int_a^b|f(t)|^p\,dt\biggr)^{1/p},\qquad
B=\biggl(\int_a^b|g(t)|^q\,dt\biggr)^{1/q},
\]
se tiene, para cada $t\in[a,b]$,
\[
\frac{|f(t)g(t)|}{AB}\leq\frac{|f(t)|^p}{pA^p}+\frac{|g(t)|^q}{qB^q},
\]
e integrando se obtiene (i). Para (ii) se sigue, palabra por palabra, la prueba del \cref{teo:minkowski-rn} usando (i) en lugar de Hölder en sumas.
\end{proof}

\begin{corollary}
\label{cor:lp-metrica}
Para $1\leq p<\infty$,
\[
  d_p(f,g)=\biggl(\int_a^b|f(t)-g(t)|^p\,dt\biggr)^{1/p}
\]
es una métrica sobre $C([a,b])$, pues la desigualdad triangular se sigue de aplicar (ii) a $f-g=(f-h)+(h-g)$.
\end{corollary}

\begin{example}[Métrica del taxi en funciones]
\label{ej:l1}
Sobre $C([a,b])$, la integral
\[
  d_1(f,g)=\int_a^b|f(t)-g(t)|\,dt
\]
define otra métrica, distinta de la del supremo. La convergencia en $d_1$ se llama \emph{convergencia en media}. (La integral es la de Riemann.)
\end{example}

\subsection{Métricas en los espacios de sucesiones}

Los espacios de sucesiones constituyen los ejemplos más importantes de espacios métricos infinito-dimensionales.

\begin{definition}[Espacios $\ell^p$ de sucesiones]
\label{def:spaces-ellp}
Para $1\leq p<\infty$ se define
\[
  \ell^p=\Bigl\{x=(x_n)\in\mathbb{R}^{\mathbb{N}}:\sum_{n=1}^{\infty}|x_n|^p<\infty\Bigr\},
\]
y para $p=\infty$,
\[
  \ell^{\infty}=\Bigl\{x=(x_n)\in\mathbb{R}^{\mathbb{N}}:\sup_{n\in\mathbb{N}}|x_n|<\infty\Bigr\}.
\]
\end{definition}

\begin{proposition}
\label{prop:metricas-sucesiones}
Para $1\leq p<\infty$, la aplicación
\[
  d_p(x,y)=\Bigl(\sum_{n=1}^{\infty}|x_n-y_n|^p\Bigr)^{1/p}
\]
define una métrica sobre $\ell^p$; y $d_\infty(x,y)=\sup_{n\in\mathbb{N}}|x_n-y_n|$ define una métrica sobre $\ell^{\infty}$.
\end{proposition}
\begin{proof}
La convergencia de las series está garantizada por la estimación $|x_n-y_n|^p\leq 2^p(|x_n|^p+|y_n|^p)$, y la desigualdad triangular es la versión en series de las desigualdades de Minkowski (\cref{teo:minkowski-rn}, \cref{teo:desigualdades-integrales}), que se extienden por paso al límite a partir de los truncamientos finitos. El caso $p=\infty$ es inmediato.
\end{proof}

\begin{example}[Métricas $\ell^p$ sobre $\mathbb{R}^n$]
\label{ej:lp}
Para $1\leq p<\infty$, la fórmula
\[
  d_p(x,y)=\biggl(\sum_{i=1}^{n}|x_i-y_i|^p\biggr)^{1/p}
\]
define una métrica. El caso $p=2$ es la euclidiana (véase el \cref{cor:euclidiana-metrica}); $p=1$ se llama \emph{taxicab} o \emph{Manhattan}; y el límite $p\to\infty$ produce la métrica del máximo $d_\infty(x,y)=\max_i|x_i-y_i|$.
\end{example}

\begin{example}[Espacio de Hilbert $\ell^2$]
\label{ej:hilbert}
El espacio de las sucesiones cuadrado-sumables,
\[
  \ell^2=\Bigl\{x=(x_n)\in\mathbb{R}^{\mathbb{N}}:\sum_{n=1}^{\infty}x_n^2<\infty\Bigr\},
\]
con la métrica $d(x,y)=(\sum|x_n-y_n|^2)^{1/2}$, es un espacio métrico. La convergencia de la serie que define la métrica se garantiza por la desigualdad de Cauchy--Schwarz en cada truncamiento finito y un paso al límite.
\end{example}

El espacio $\ell^2$ es el prototipo de \emph{espacio de Hilbert}; los espacios $\ell^p$ con $p\neq 2$, en cambio, no admiten un producto interior que induzca su métrica.

\subsection{Productos de espacios métricos}

\begin{definition}
\label{def:producto-metrico}
Sean $(X,d_X)$ y $(Y,d_Y)$ espacios métricos. Sobre $X\times Y$ se pueden definir varias métricas; la más usuales son:
\begin{itemize}
  \item $d_1((x_1,y_1),(x_2,y_2))=d_X(x_1,x_2)+d_Y(y_1,y_2)$;
  \item $d_2((x_1,y_1),(x_2,y_2))=(d_X(x_1,x_2)^2+d_Y(y_1,y_2)^2)^{1/2}$;
  \item $d_\infty((x_1,y_1),(x_2,y_2))=\max\{d_X(x_1,x_2),d_Y(y_1,y_2)\}$.
\end{itemize}
Todas son equivalentes.
\end{definition}

\subsection{Ejercicios}

\begin{exercise}
Verifique que la métrica del supremo sobre $C([a,b])$ satisface la desigualdad triangular.
\end{exercise}

\section{Bolas abiertas y diámetro}

\subsection{Bolas abiertas}

\begin{definition}
\label{def:bola-abierta}
Sea $(X,d)$ un espacio métrico, $x\in X$ y $r>0$. La \emph{bola abierta} de centro $x$ y radio $r$ es
\[
  B(x,r)=\{y\in X:\,d(x,y)<r\}.
\]
\end{definition}

Las bolas abiertas constituyen los ``ladrillos'' con los que se construye toda la topología métrica. Nótese que el radio $r$ es un número real positivo arbitrario; esta arbitrariedad es la que permite capturar la noción de ``proximidad local'' sin cuantificarla de manera única.

\begin{notation}
\label{not:bolas}
A lo largo del libro, $B(x,r)$ denota la bola abierta, $\overline{B}(x,r)$ la bola cerrada $\{y:d(x,y)\leq r\}$, y $S(x,r)$ la esfera $\{y:d(x,y)=r\}$.
\end{notation}

\subsection{El diámetro de un conjunto}

\begin{definition}
\label{def:diametro}
Sea $(X,d)$ un espacio métrico y $A\subseteq X$ no vacío. El \emph{diámetro} de $A$ es
\[
  \operatorname{diam}(A)=\sup\{d(x,y):\,x,y\in A\}\in[0,+\infty].
\]
Diremos que $A$ es \emph{acotado} si $\operatorname{diam}(A)<+\infty$; en caso contrario, $A$ es \emph{no acotado}.
\end{definition}

Todo conjunto finito es acotado. En la métrica discreta, $\operatorname{diam}(A)=0$ si $A$ tiene un solo punto y $1$ en cualquier otro caso. Para la bola abierta de la métrica usual de $\mathbb{R}$, $\operatorname{diam}\bigl((a-\delta,a+\delta)\bigr)=2\delta$. En general, si $A\subseteq B(x,r)$, entonces $\operatorname{diam}(A)\leq 2r$. La noción de diámetro será esencial en el \cref{ch:grandes-teoremas}, donde aparecerán familias de cerrados encajados cuyos diámetros tienden a cero.

\subsection{Ejercicios}

\begin{exercise}
Demuestre que en un espacio métrico discreto, toda bola abierta $B(x,r)$ con $r\leq 1$ coincide con $\{x\}$.
\end{exercise}

\section{Funciones continuas}
\label{sec:funciones-continuas}

\subsection{Definición de continuidad}

La definición $\epsilon$-$\delta$ de continuidad que se utiliza habitualmente en cálculo de una variable supone, aunque no siempre se haga explícito, que estamos trabajando con la métrica euclidiana. En efecto, una función $f\colon\mathbb{R}\to\mathbb{R}$ es continua en un punto $a\in\mathbb{R}$ si, para todo $\epsilon>0$, existe $\delta>0$ tal que, para todo $x\in\mathbb{R}$,
\[
  |x-a|<\delta\implies|f(x)-f(a)|<\epsilon.
\]
Esta definición utiliza directamente la distancia euclidiana en $\mathbb{R}$. La condición $|x-a|<\delta$ describe el intervalo abierto $(a-\delta,a+\delta)$, mientras que la condición $|f(x)-f(a)|<\epsilon$ describe el intervalo abierto $(f(a)-\epsilon,f(a)+\epsilon)$. Desde el punto de vista de los espacios métricos, estos intervalos son precisamente bolas abiertas. Por lo tanto, la idea fundamental de la continuidad no depende de la estructura particular de $\mathbb{R}$ ni de la métrica euclidiana, sino que puede formularse en cualquier espacio métrico mediante sus respectivas bolas abiertas.

\begin{definition}
\label{def:continuidad-metrica}
Sean $(X,d_X)$ y $(Y,d_Y)$ espacios métricos. Decimos que $f\colon X\to Y$ es \emph{continua} en el punto $x_0\in X$ si, para todo $\epsilon>0$, existe $\delta>0$ tal que, para todo $x\in X$,
\[
  d_X(x,x_0)<\delta\implies d_Y\bigl(f(x),f(x_0)\bigr)<\epsilon.
\]
\end{definition}
Si una función $f\colon X\to Y$ es continua en todo punto de $X$, decimos simplemente que $f$ es \emph{continua en $X$} o que es una \emph{función continua}. La definición anterior puede interpretarse mediante bolas abiertas. En efecto, para todo $\epsilon>0$ debe existir $\delta>0$ tal que
\[
  f\bigl(B_{d_X}(x_0,\delta)\bigr)\subseteq B_{d_Y}\bigl(f(x_0),\epsilon\bigr).
\]
Es decir, los puntos suficientemente cercanos a $x_0$ son enviados por $f$ a puntos suficientemente cercanos a $f(x_0)$.

\subsection{Ejemplos básicos}

\begin{example}[La función afín en la métrica usual]
Sean $X=Y=\mathbb{R}$ con la métrica usual $d(x,y)=|x-y|$ y sea $f\colon\mathbb{R}\to\mathbb{R}$ definida por $f(x)=2x+1$. Demostraremos que $f$ es continua en cualquier punto $x_0\in\mathbb{R}$.

Sea $\epsilon>0$. Buscamos $\delta>0$ tal que $|x-x_0|<\delta$ implique $|f(x)-f(x_0)|<\epsilon$. Observamos que
\[
  |f(x)-f(x_0)|=|(2x+1)-(2x_0+1)|=2|x-x_0|.
\]
Basta entonces elegir $\delta=\epsilon/2$: si $|x-x_0|<\delta$,
\[
  |f(x)-f(x_0)|=2|x-x_0|<2\delta=\epsilon.
\]
Así, $f$ es continua en $x_0$, y como $x_0$ fue arbitrario, $f$ es continua en todo $\mathbb{R}$.
\end{example}

\begin{example}
Sean $X=\mathbb{R}$ y $Y=\mathbb{R}$, ambos dotados de la métrica discreta
\[
d(x,y)=
\begin{cases}
0, & x=y,\\
1, & x\neq y.
\end{cases}
\]
Consideremos la función
\[
f\colon X\to Y,
\qquad
f(x)=x^2.
\]
Demostraremos que $f$ es continua en cualquier punto $x_0\in X$.

Sea $\epsilon>0$. Tomemos $\delta=1$. Si
\[
d_X(x,x_0)<\delta,
\]
entonces, por la definición de la métrica discreta, necesariamente $x=x_0$. Por consiguiente,
\[
d_Y\bigl(f(x),f(x_0)\bigr)
=
d_Y\bigl(x^2,x_0^2\bigr)
=
0
<
\epsilon.
\]
Por lo tanto, $f$ es continua en $x_0$. Como $x_0$ fue arbitrario, $f$ es continua en todo $X$.

Este ejemplo muestra que la continuidad no depende únicamente de la expresión algebraica de una función, sino también de las métricas que determinan la estructura de los espacios involucrados.
\end{example}

\subsection{Isometrías}

Una clase especial de funciones continuas está formada por aquellas que conservan de manera exacta la distancia entre los puntos.

\begin{definition}
\label{def:isometria}
Sean $(X,d_X)$ y $(Y,d_Y)$ espacios métricos. Una función $f\colon X\to Y$ es una \emph{isometría} si
\[
  d_Y\bigl(f(x),f(x')\bigr)=d_X(x,x')\quad\text{para todo }x,x'\in X.
\]
Si, además, $f$ es sobreyectiva, se dice que los espacios $X$ e $Y$ son \emph{isométricos}.
\end{definition}

Toda isometría es inyectiva: si $f(x)=f(x')$, entonces $d_X(x,x')=d_Y\bigl(f(x),f(x')\bigr)=0$, de donde $x=x'$. Además, una isometría es siempre continua: dado $\epsilon>0$ basta tomar $\delta=\epsilon$, pues
\[
  d_X(x,x_0)<\delta\implies d_Y\bigl(f(x),f(x_0)\bigr)=d_X(x,x_0)<\epsilon.
\]
Si la isometría es sobreyectiva, su inversa también lo es, luego también es continua, y por tanto toda isometría sobreyectiva es en particular un homeomorfismo. La noción de isometría es, sin embargo, más fina que la de homeomorfismo: exige la conservación exacta de las distancias, no solo de las propiedades topológicas.

\begin{example}
La aplicación $f\colon\mathbb{R}\to\mathbb{R}$ definida por $f(x)=-x$ es una isometría de $\mathbb{R}$ sobre sí mismo con la métrica usual, pues $|f(x)-f(y)|=|-x+y|=|x-y|$. En cambio, la inclusión $i\colon\mathbb{Q}\to\mathbb{R}$ es una isometría (restricción de la distancia usual) que no es sobreyectiva.
\end{example}

\subsection{Homeomorfismos}

Existen situaciones en las que ciertas propiedades de un espacio
métrico son sencillas de identificar directamente, mientras que las
mismas propiedades pueden resultar más difíciles de estudiar en otro
espacio. Esto motiva la búsqueda de funciones que permitan relacionar
dos espacios métricos de manera que su estructura esencial se conserve.

La idea consiste en establecer una correspondencia entre los puntos de
ambos espacios que permita trasladar información de uno a otro. Si
dicha correspondencia y su inversa son continuas, entonces ambos
espacios pueden considerarse equivalentes desde el punto de vista
topológico. Esta noción conduce al concepto de homeomorfismo.

\begin{definition}
Sean $(X,d_X)$ y $(Y,d_Y)$ espacios métricos. Decimos que una función
\[
f\colon X\to Y
\]
es un \emph{homeomorfismo} si $f$ es biyectiva, continua y su función
inversa
\[
f^{-1}\colon Y\to X
\]
también es continua.

Cuando existe un homeomorfismo entre $X$ y $Y$, decimos que $X$ y $Y$
son \emph{homeomorfos} y escribimos
\[
X\cong Y.
\]
\end{definition}

Uno de los ejemplos más importantes de homeomorfismo consiste en
relacionar el intervalo abierto $(-1,1)$ con toda la recta real
\[
\mathbb{R}=(-\infty,\infty).
\]
Aunque estos conjuntos son diferentes, desde el punto de vista
topológico poseen la misma estructura. El siguiente ejemplo muestra
explícitamente una función que establece dicha correspondencia.

\begin{example}
Sean
\[
X=\mathbb{R}
\qquad\text{y}\qquad
Y=(-1,1),
\]
ambos dotados de la métrica euclidiana. Consideremos la función
\[
f\colon X\to Y,
\qquad
f(x)=\frac{2}{\pi}\arctan(x).
\]

Como $\arctan$ es estrictamente creciente y continua en $\mathbb{R}$,
y además
\[
\lim_{x\to-\infty}\frac{2}{\pi}\arctan(x)=-1,
\qquad
\lim_{x\to\infty}\frac{2}{\pi}\arctan(x)=1,
\]
la función $f$ es una biyección de $\mathbb{R}$ sobre $(-1,1)$.

Para determinar su inversa, sea $y=f(x)$. Entonces
\[
y=\frac{2}{\pi}\arctan(x),
\]
de donde
\[
\arctan(x)=\frac{\pi}{2}y.
\]
Aplicando la función tangente obtenemos
\[
x=\tan\left(\frac{\pi}{2}y\right).
\]
Por consiguiente,
\[
f^{-1}\colon (-1,1)\to\mathbb{R},
\qquad
f^{-1}(y)=\tan\left(\frac{\pi}{2}y\right).
\]

La función $f$ es continua porque $\arctan$ es continua en
$\mathbb{R}$. Asimismo, $f^{-1}$ es continua porque la función
tangente es continua en
$\left(-\frac{\pi}{2},\frac{\pi}{2}\right)$. Por lo tanto, $f$ es un
homeomorfismo y concluimos que
\[
\mathbb{R}\cong(-1,1).
\]

Así, aunque $\mathbb{R}$ y $(-1,1)$ son conjuntos diferentes, son
equivalentes desde el punto de vista de la estructura topológica.
\end{example}

\section{Subespacios}

\subsection{Definición y ejemplos}

\begin{definition}
\label{def:subespacio}
Sea $(X,d)$ un espacio métrico y $Y\subseteq X$ no vacío. La restricción $d|_{Y\times Y}$ define una métrica sobre $Y$, llamada la \emph{métrica inducida}. Al par $(Y,d|_{Y\times Y})$ se le llama \emph{subespacio métrico}.
\end{definition}

\begin{example}
$\mathbb{Q}$, como subconjunto de $\mathbb{R}$ con la métrica usual, es un subespacio métrico. Como veremos en el \cref{ch:completitud}, su comportamiento respecto a sucesiones de Cauchy difiere radicalmente del de $\mathbb{R}$.
\end{example}

\subsection{Continuidad y subespacios}

La construcción de subespacios interacciona de manera natural con las funciones continuas de la \cref{sec:funciones-continuas}.

\begin{proposition}
\label{prop:inclusion-isometria}
Sea $(X,d)$ un espacio métrico y $Y\subseteq X$ no vacío con la métrica inducida. La inclusión $i\colon Y\to X$, $i(y)=y$, es una isometría. En particular, es continua.
\end{proposition}
\begin{proof}
Para cualesquiera $y_1,y_2\in Y$, la métrica inducida está definida por $d_Y(y_1,y_2)=d(y_1,y_2)$, de modo que $d\bigl(i(y_1),i(y_2)\bigr)=d_Y(y_1,y_2)$.
\end{proof}

\begin{proposition}
\label{prop:restriccion-continua}
Sean $(X,d_X)$, $(Y,d_Y)$ y $(Z,d_Z)$ espacios métricos y $Y\subseteq X$ un subespacio. Si $f\colon X\to Z$ es continua en $x_0\in Y$, entonces la restricción $f|_Y\colon Y\to Z$ es continua en $x_0$.
\end{proposition}
\begin{proof}
Dado $\epsilon>0$, por la continuidad de $f$ existe $\delta>0$ tal que $d_X(x,x_0)<\delta$ implica $d_Z\bigl(f(x),f(x_0)\bigr)<\epsilon$ para todo $x\in X$. En particular, para todo $y\in Y$ con $d_Y(y,x_0)=d_X(y,x_0)<\delta$ se cumple
\[
  d_Z\bigl(f|_Y(y),f|_Y(x_0)\bigr)=d_Z\bigl(f(y),f(x_0)\bigr)<\epsilon.
\]
\end{proof}

\subsection{Ejercicios}

\begin{exercise}
Caracterice los puntos aislados de $\mathbb{Q}$ como subespacio de $\mathbb{R}$.
\end{exercise}

\section{Métricas equivalentes}

\subsection{Las métricas usuales de $\mathbb{R}$}

Sobre $\mathbb{R}$, todas las métricas de la familia $\ell^p$ se colapsan en una sola: para $1\leq p<\infty$,
\[
  d_p(x,y)=\bigl(|x-y|^p\bigr)^{1/p}=|x-y|=d_\infty(x,y).
\]
Por tanto, en la recta real las métricas euclidiana, taxicab y del máximo coinciden: no hay comparación alguna en una dimensión. La situación cambia en $\mathbb{R}^n$.

\begin{example}
\label{ej:comparacion-metricas-rn}
En $\mathbb{R}^n$, para todo $x,y\in\mathbb{R}^n$ se cumplen las estimaciones
\[
  d_\infty(x,y)\leq d_2(x,y)\leq d_1(x,y)\leq n\,d_\infty(x,y).
\]
La primera desigualdad se debe a que $|x_i-y_i|\leq\bigl(\sum_{j=1}^{n}|x_j-y_j|^2\bigr)^{1/2}$ para cada $i$; la segunda se sigue de Cauchy--Schwarz aplicada a los vectores $x-y$ y $(1,\dots,1)$; la tercera es inmediata.
\end{example}

Estas estimaciones muestran cómo las tres métricas se dominan mutuamente con constantes uniformes e independientes de los puntos, y son el modelo que motiva la definición general de equivalencia presentada a continuación.

\subsection{Métricas equivalentes}

\begin{definition}
\label{def:metricas-equivalentes}
Dos métricas $d$ y $d'$ sobre $X$ son \emph{equivalentes} si existen constantes $c,C>0$ tales que
\[
  c\,d(x,y)\leq d'(x,y)\leq C\,d(x,y)\quad\text{para todo }x,y\in X.
\]
\end{definition}

\begin{proposition}
\label{prop:equiv-misma-topologia}
Métricas equivalentes generan la misma familia de conjuntos abiertos (véase el \cref{ch:topologia}).
\end{proposition}

\begin{example}
En $\mathbb{R}^n$, todas las métricas $\ell^p$ son equivalentes, como se desprende de las estimaciones del \cref{ej:comparacion-metricas-rn}. Sin embargo, la métrica discreta no es equivalente a ninguna de ellas.
\end{example}

\subsection{Ejercicios}

\begin{exercise}
Sea $d$ una métrica sobre $X$. Pruebe que $d'(x,y)=\min\{1,d(x,y)\}$ es también una métrica y que induce la misma topología que $d$.
\end{exercise}

\begin{problem}
Sea $(X,d)$ un espacio métrico. Defina $\rho(x,y)=\frac{d(x,y)}{1+d(x,y)}$. Demuestre que $\rho$ es una métrica acotada que genera la misma topología que $d$. ¿Es $\rho$ equivalente a $d$ en el sentido de la \cref{def:metricas-equivalentes}?
\end{problem}

\section{Abiertos}

La distinción entre puntos ``interiores'', ``exteriores'' y ``frontera'' de un conjunto era intuitiva para los geométricos griegos, pero no adquirió un estatus formal preciso hasta el desarrollo de la topología en el siglo~XX. Henri Poincaré, en sus estudios sobre la geometría de dimensiones superiores, y más tarde Maurice Fréchet y Felix Hausdorff, comprendieron que las propiedades estructurales de los espacios ---abiertos, cerrados, convergencia, compacidad--- podían estudiarse abstractamente, sin necesidad de una distancia explícita. En este capítulo desarrollamos estas nociones en el contexto métrico, que es el puente natural hacia la topología general.

\subsection{Conjuntos abiertos}

\begin{definition}
\label{def:abierto}
Un subconjunto $U\subseteq X$ es \emph{abierto} si para cada $x\in U$ existe $r>0$ tal que $B(x,r)\subseteq U$.
\end{definition}

\begin{theorem}
\label{teo:topologia-metrica}
La familia $\tau$ de todos los conjuntos abiertos de un espacio métrico $(X,d)$ satisface:
\begin{enumerate}
  \item[(T1)] $\emptyset\in\tau$ y $X\in\tau$.
  \item[(T2)] La unión arbitraria de elementos de $\tau$ pertenece a $\tau$.
  \item[(T3)] La intersección finita de elementos de $\tau$ pertenece a $\tau$.
\end{enumerate}
\end{theorem}

\begin{proof}
(T1) es trivial. Para (T2), si $U=\bigcup_{i\in I}U_i$ con cada $U_i$ abierto, entonces para cada $x\in U$ existe $i$ tal que $x\in U_i$, y por tanto una bola $B(x,r)\subseteq U_i\subseteq U$. Para (T3), si $U_1,\dots,U_n$ son abiertos y $x\in\bigcap_{k=1}^n U_k$, elegimos $r_k>0$ con $B(x,r_k)\subseteq U_k$; entonces $B(x,r)\subseteq\bigcap U_k$ donde $r=\min\{r_1,\dots,r_n\}>0$.
\end{proof}

\subsection{Entornos}

\begin{definition}
\label{def:entorno}
Un \emph{entorno} (o \emph{vecindad}) de $x\in X$ es un conjunto $V\subseteq X$ tal que existe un abierto $U$ con $x\in U\subseteq V$. Equivalentemente, $V$ contiene una bola abierta centrada en $x$.
\end{definition}

La noción de entorno permite reformular la convergencia y la continuidad en términos topológicos, sin mencionar explícitamente el radio de la bola.

\subsection{Interior}

\begin{definition}
\label{def:interior}
El \emph{interior} de $A\subseteq X$ es el mayor abierto contenido en $A$:
\[
  \operatorname{int}(A)=\bigcup\{U\subseteq X:\,U\text{ abierto y }U\subseteq A\}.
\]
\end{definition}

\begin{proposition}
\label{prop:int-complemento}
$\operatorname{int}(A)=X\setminus\overline{X\setminus A}$.
\end{proposition}

\subsection{Puntos aislados}

\begin{definition}
\label{def:punto-aislado}
Un punto $x\in A$ es \emph{aislado} en $A$ si existe $r>0$ tal que $B(x,r)\cap A=\{x\}$.
\end{definition}

\subsection{Ejercicios}

\begin{exercise}
Pruebe que un conjunto es abierto si y solo si coincide con su interior.
\end{exercise}

\section{Cerrado}

\subsection{Conjuntos cerrados}

\begin{definition}
\label{def:cerrado}
Un subconjunto $F\subseteq X$ es \emph{cerrado} si su complementario $X\setminus F$ es abierto.
\end{definition}

\begin{proposition}
\label{prop:prop-cerrados}
En un espacio métrico $(X,d)$:
\begin{enumerate}
  \item[(i)] $\emptyset$ y $X$ son cerrados.
  \item[(ii)] La intersección arbitraria de cerrados es cerrada.
  \item[(iii)] La unión finita de cerrados es cerrada.
\end{enumerate}
\end{proposition}

\subsection{Convergencia de sucesiones}

La definición formal de convergencia utiliza la desigualdad triangular de la métrica y es la piedra angular sobre la que se construye el análisis en espacios métricos. Dejamos el estudio sistemático de las sucesiones para el \cref{ch:sucesiones}; aquí recogemos únicamente la definición y su relación con los conjuntos cerrados.

\begin{definition}
\label{def:convergencia-sucesiones}
Una sucesión $(x_n)$ en un espacio métrico $(X,d)$ \emph{converge} a $x\in X$ si para todo $\epsilon>0$ existe $N\in\mathbb{N}$ tal que
\[
  d(x_n,x)<\epsilon\quad\text{para todo }n\geq N.
\]
En tal caso escribimos $x_n\to x$, o $\lim_{n\to\infty}x_n=x$.
\end{definition}

En términos de bolas, $x_n\to x$ significa que para todo $\epsilon>0$ la bola $B(x,\epsilon)$ contiene a todos los términos de la sucesión a partir de cierto índice: la sucesión ``entra y permanece'' en cada entorno del límite. Nótese que $x_n\to x$ si y solo si la sucesión numérica $d(x_n,x)$ converge a $0$, de modo que la definición formal se reduce, en última instancia, a la convergencia en $\mathbb{R}$.

\begin{proposition}
\label{prop:cerrados-limites}
Sea $(X,d)$ un espacio métrico y $F\subseteq X$. Las siguientes afirmaciones son equivalentes:
\begin{enumerate}
  \item[(i)] $F$ es cerrado.
  \item[(ii)] Para toda sucesión $(x_n)$ en $F$ con $x_n\to x\in X$, se tiene $x\in F$.
\end{enumerate}
\end{proposition}
\begin{proof}
(i)$\Rightarrow$(ii): si $x\notin F$, entonces $x$ pertenece al abierto $X\setminus F$, luego existe $\epsilon>0$ con $B(x,\epsilon)\subseteq X\setminus F$. Pero $x_n\to x$ implica que $x_n\in B(x,\epsilon)$ para $n$ suficientemente grande, contradiciendo que $x_n\in F$.
(ii)$\Rightarrow$(i): probamos que $X\setminus F$ es abierto. Si no lo fuera, existiría $x\in X\setminus F$ tal que $B(x,1/n)\nsubseteq X\setminus F$ para todo $n$, es decir, existiría $x_n\in F\cap B(x,1/n)$; en tal caso $x_n\to x$, y por (ii) $x\in F$, contradicción.
\end{proof}

\subsection{Puntos de acumulación}

\begin{definition}
\label{def:punto-acumulacion}
Un punto $x\in X$ es un \emph{punto de acumulación} de $A\subseteq X$ si todo entorno de $x$ contiene al menos un punto de $A$ distinto de $x$.
\end{definition}

\begin{theorem}
\label{teo:bolzano-weierstrass-r}
Todo subconjunto infinito y acotado de $\mathbb{R}$ posee al menos un punto de acumulación en $\mathbb{R}$.
\end{theorem}

Este resultado, el teorema de Bolzano--Weierstrass para la recta real, será generalizado a espacios métricos compactos en el \cref{ch:compacidad}.

\subsection{Clausura}

\begin{definition}
\label{def:clausura}
La \emph{clausura} (o \emph{adherencia}) de $A\subseteq X$ es el menor cerrado que contiene a $A$, es decir,
\[
  \overline{A}=\bigcap\{F\subseteq X:\,F\text{ cerrado y }A\subseteq F\}.
\]
\end{definition}

\begin{theorem}[Caracterización secuencial de la clausura]
\label{teo:carac-clausura}
Un punto $x\in X$ pertenece a $\overline{A}$ si y solo si existe una sucesión $(a_n)$ en $A$ tal que $a_n\to x$.
\end{theorem}

\begin{proof}
Si $x\in\overline{A}$ y $x\notin A$, entonces para cada $n\in\mathbb{N}$ la bola $B(x,1/n)$ intersecta a $A$ (de lo contrario $X\setminus B(x,1/n)$ sería un cerrado que contiene a $A$ pero no a $x$). Elegimos $a_n\in A\cap B(x,1/n)$; entonces $d(a_n,x)<1/n\to 0$. Recíprocamente, si $a_n\to x$ con $a_n\in A$, todo cerrado que contenga a $A$ debe contener el límite de toda sucesión convergente en él, luego $x\in\overline{A}$.
\end{proof}

La \cref{teo:carac-clausura} pone de manifiesto la relación entre los puntos de contacto (los de $\overline{A}$) y los puntos límite de sucesiones: todo punto de $A$ y todo punto de acumulación de $A$ pertenecen a $\overline{A}$.

\subsection{Frontera}

\begin{definition}
\label{def:frontera}
La \emph{frontera} (o \emph{borde}) de $A\subseteq X$ es
\[
  \partial A=\overline{A}\setminus\operatorname{int}(A).
\]
\end{definition}

\begin{proposition}
\label{prop:frontera-carac}
$x\in\partial A$ si y solo si todo entorno de $x$ intersecta tanto $A$ como $X\setminus A$.
\end{proposition}

\subsection{Conjuntos densos}

\begin{definition}
\label{def:densidad}
Un subconjunto $D\subseteq X$ es \emph{denso} en $X$ si $\overline{D}=X$. Equivalentemente, todo abierto no vacío de $X$ contiene al menos un punto de $D$.
\end{definition}

\begin{theorem}
\label{teo:q-denso-r}
$\mathbb{Q}$ es denso en $\mathbb{R}$ con la métrica usual.
\end{theorem}
\begin{proof}
Ya se demostró en el \cref{teo:densidad-q} del \cref{ap:numeros-reales}.
\end{proof}

\subsection{Bases topológicas}

\begin{definition}
\label{def:base}
Una familia $\mathcal{B}$ de abiertos es una \emph{base} de la topología si todo abierto no vacío es unión de elementos de $\mathcal{B}$.
\end{definition}

\begin{theorem}
\label{teo:base-metrica}
Sea $(X,d)$ un espacio métrico separable y $D\subseteq X$ un subconjunto denso numerable. La familia de todas las bolas abiertas $B(d,r)$ con $d\in D$ y $r>0$ racional es una base numerable de la topología métrica.
\end{theorem}

\begin{proof}
El conjunto $\{B(d,r):d\in D,\ r\in\mathbb{Q}_{>0}\}$ es numerable como unión sobre $d\in D$ de las familias numerables de radios. Dado un abierto $U$ y $x\in U$, existe $\rho>0$ con $B(x,\rho)\subseteq U$. Elegimos $q\in\mathbb{Q}$ con $0<q<\rho/2$ y, por densidad, $d\in D\cap B(x,q)$. Entonces $x\in B(d,q)\subseteq B(x,2q)\subseteq B(x,\rho)\subseteq U$, luego $U$ se escribe como unión de bolas de la familia.
\end{proof}

Esta caracterización conecta las bases con la separabilidad, estudiada a continuación.

\subsection{Separabilidad}

\begin{definition}
\label{def:separabilidad}
Un espacio métrico $(X,d)$ es \emph{separable} si posee un subconjunto denso numerable.
\end{definition}

Por la \cref{teo:q-denso-r} y la numerabilidad de $\mathbb{Q}$, el espacio $\mathbb{R}$ es separable. La \cref{teo:base-metrica} muestra además que todo espacio métrico separable admite una base numerable; el recíproco también es cierto en espacios métricos, de modo que en este contexto la separabilidad es equivalente a la existencia de una base numerable. La separabilidad se estudia con detalle en el \cref{ch:espacios-separables}.

\subsection{Axiomas de separación en espacios métricos}

Los axiomas de separación clasifican los espacios topológicos según su capacidad para distinguir puntos o cerrados mediante abiertos. Aunque estos axiomas se estudian en profundidad en el \cref{ch:topologia-general}, los necesitamos ya en los capítulos de compacidad y en los grandes teoremas.

\begin{definition}
\label{def:separacion-metrico}
Un espacio topológico $(X,\tau)$ satisface:
\begin{itemize}
  \item $T_0$: para puntos distintos, al menos uno posee un entorno que no contiene al otro.
  \item $T_1$ (Fréchet): para puntos distintos, cada uno tiene un entorno que no contiene al otro.
  \item $T_2$ (Hausdorff): puntos distintos tienen entornos disjuntos.
  \item $T_3$ (regular): es $T_1$ y, para todo cerrado $F$ y $x\notin F$, existen entornos disjuntos de $F$ y $x$.
  \item $T_4$ (normal): es $T_1$ y, para cualesquiera cerrados disjuntos $A,B$, existen entornos disjuntos de $A$ y $B$.
\end{itemize}
\end{definition}

\begin{theorem}
\label{teo:metrico-normal}
Todo espacio métrico es normal (y por tanto $T_2$, $T_3$, $T_4$).
\end{theorem}
\begin{proof}
Sea $A,B$ cerrados disjuntos. Para cada $a\in A$ definimos $r_a=d(a,B)/2>0$, y para cada $b\in B$ definimos $s_b=d(b,A)/2>0$. Entonces $U=\bigcup_{a\in A}B(a,r_a)$ y $V=\bigcup_{b\in B}B(b,s_b)$ son abiertos disjuntos que contienen a $A$ y $B$ respectivamente.
\end{proof}

\begin{corollary}[Urysohn en espacios métricos]
\label{cor:urysohn-metrico}
Sean $A,B$ cerrados disjuntos en un espacio métrico $(X,d)$. La función $f\colon X\to[0,1]$ definida por
\[
  f(x)=\frac{d(x,A)}{d(x,A)+d(x,B)}
\]
es continua, $f|_A=0$ y $f|_B=1$.
\end{corollary}
\begin{proof}
La función $d(x,A)=\inf_{a\in A}d(x,a)$ es continua (Lipschitz con constante $1$), y el denominador nunca se anula porque $A$ y $B$ son disjuntos y cerrados.
\end{proof}

\subsection{Ejercicios}

\begin{exercise}
Demuestre que la frontera de una bola abierta en $\mathbb{R}^n$ con la métrica euclidiana es la esfera correspondiente.
\end{exercise}

\begin{problem}
Sea $X$ un espacio métrico y $A\subseteq X$. Demuestre que $\partial(\partial A)\subseteq\partial A$. ¿Es siempre $\partial(\partial A)=\partial A$?
\end{problem}

\section{Caracterización de $\mathbb{R}$ y el conjunto de Cantor}

\subsection{Estructura de los abiertos de $\mathbb{R}$}

El teorema que presentamos a continuación caracteriza por completo los subconjuntos abiertos de la recta real y es una de las primeras instancias de un \emph{teorema de estructura}: describe todos los objetos de una clase a partir de piezas elementales, en este caso, de intervalos abiertos.

\begin{theorem}[Caracterización de los abiertos de $\mathbb{R}$]
\label{teo:abiertos-r}
Todo subconjunto abierto no vacío $U\subseteq\mathbb{R}$ puede escribirse como una unión numerable de intervalos abiertos disjuntos por pares:
\[
  U=\bigcup_{n\in\mathbb{N}}(a_n,b_n).
\]
La descomposición es única salvo reordenación de los intervalos.
\end{theorem}
\begin{proof}
Definamos sobre $U$ la siguiente relación: $x\sim y$ si el intervalo cerrado de extremos $x$ e $y$ está contenido en $U$. Esta relación es de equivalencia. Cada clase de equivalencia es un intervalo, pues dados dos puntos de la clase, el intervalo que los une está contenido en $U$; es abierto porque $U$ lo es, y sus extremos no pertenecen a $U$ (si un extremo estuviera en $U$, se fusionaría con la clase vecina). Las clases son disjuntas por construcción y su familia es numerable porque, por la densidad de $\mathbb{Q}$ (véase \cref{teo:q-denso-r}), cada clase contiene al menos un racional, y el conjunto de racionales es numerable.
\end{proof}

\begin{remark}
Los intervalos $(a_n,b_n)$ son precisamente las componentes conexas de $U$ (véase el \cref{ch:conexidad}). Obsérvese que una descomposición análoga no es válida en $\mathbb{R}^n$ para $n\geq 2$: los abiertos del plano no son, en general, uniones disjuntas de rectángulos abiertos.
\end{remark}

\subsection{El conjunto de Cantor}

La caracterización anterior muestra que los abiertos de la recta se describen mediante intervalos. El conjunto de Cantor, por el contrario, es el ejemplo paradigmático de conjunto cerrado que desafía la intuición geométrica.

\begin{definition}[Conjunto de Cantor]
\label{def:conjunto-cantor}
Se define $K_0=[0,1]$. Dado $K_n$, formado por $2^n$ intervalos cerrados disjuntos de longitud $3^{-n}$, se construye $K_{n+1}$ eliminando de cada uno de esos intervalos el tercio central abierto. El \emph{conjunto de Cantor} es
\[
  K=\bigcap_{n=0}^{\infty}K_n.
\]
\end{definition}

Cada $K_n$ es una unión finita de conjuntos cerrados, luego es cerrado; $K$, como intersección numerable de cerrados, es cerrado en $[0,1]$ y, por tanto, en $\mathbb{R}$.

\begin{proposition}
\label{prop:propiedades-cantor}
El conjunto de Cantor $K$ satisface:
\begin{enumerate}
  \item[(i)] $K$ es compacto.
  \item[(ii)] $K$ tiene interior vacío: $\operatorname{int}(K)=\emptyset$.
  \item[(iii)] $K$ no tiene puntos aislados: todo punto de $K$ es punto de acumulación de $K$.
  \item[(iv)] $K$ no es numerable.
  \item[(v)] $K$ es totalmente disconexo (véase el \cref{ch:conexidad}).
\end{enumerate}
\end{proposition}
\begin{proof}
(i) es consecuencia del teorema de Heine--Borel, pues $K$ es cerrado y acotado. Para (ii), la suma de las longitudes de los intervalos que forman $K_n$ es $(2/3)^n\to 0$: si $K$ contuviera un intervalo abierto $(a,b)$, entonces la longitud de $K_n$, que contiene a $K$, sería al menos $b-a>0$ para todo $n$, contradicción. Para (iii), dado $x\in K$ y $\epsilon>0$, el punto $x$ pertenece a uno de los intervalos cerrados $I_n\subseteq K_n$, de longitud $3^{-n}$; eligiendo $n$ con $3^{-n}<\epsilon$, dos extremos de $I_n$ (que pertenecen a $K$) son puntos de $K\cap B(x,\epsilon)$ distintos de $x$, pues $x$ no puede ser al mismo tiempo ambos extremos. Para (iv), cada $x\in K$ admite un único desarrollo $x=\sum_{n=1}^{\infty}2\varepsilon_n3^{-n}$ con $\varepsilon_n\in\{0,1\}$, lo que identifica $K$ con el conjunto $\{0,1\}^{\mathbb{N}}$ de las sucesiones binarias, que es no numerable por el teorema de Cantor (véase el \cref{teo:cantor-uncountable} del \cref{ap:teoria-conjuntos}). La propiedad (v) se demuestra en la \cref{ch:conexidad}.
\end{proof}

Este conjunto singulariza la distancia entre lo \emph{pequeño} en medida y lo \emph{grande} en cardinalidad: $K$ es ``insignificante'' por su interior vacío y su longitud total nula, pero ``enorme'' por ser no numerable. El conjunto de Cantor aparecerá de nuevo en la \cref{ch:espacios-separables}, donde se muestra que es homeomorfo al producto numerable $\{0,1\}^{\mathbb{N}}$, y en el estudio de la completitud.
\chapter{Espacios topológicos}
\label{ch:topologia}
\index{Espacio topológico}\index{Topología}\index{Conjunto abierto}\index{Conjunto cerrado}\index{Adherencia}\index{Base topológica}\index{Subbase}

\section{Definición y ejemplos básicos}

Todo espacio métrico induce de manera natural una topología, cuyos
conjuntos abiertos son generados por las bolas de la \cref{def:metrica} del
capítulo anterior (véase \cref{ej:euclidea}). Sin embargo, la noción de
topología es más general y permite estudiar estructuras que no provienen de
ninguna métrica.

\begin{definition}\label{def:topologia}
Una \emph{topología} sobre un conjunto no vacío $X$ es una familia
$\tau\subseteq\mathcal{P}(X)$ que verifica:
\begin{enumerate}
  \item[(T1)] $\emptyset\in\tau$ y $X\in\tau$;
  \item[(T2)] la unión arbitraria de elementos de $\tau$ pertenece a $\tau$;
  \item[(T3)] la intersección finita de elementos de $\tau$ pertenece a $\tau$.
\end{enumerate}
Al par $(X,\tau)$ se le llama \emph{espacio topológico}, y los elementos de
$\tau$ son los \emph{conjuntos abiertos}.
\end{definition}

\begin{example}\label{ej:top-trivial}
La familia $\tau=\{\emptyset,X\}$ es una topología sobre $X$, llamada la
\emph{topología trivial} (o \emph{indiscreta}).
\end{example}

\begin{example}\label{ej:top-discreta}
La familia $\tau=\mathcal{P}(X)$ es una topología, llamada la \emph{topología
discreta}. Todo subconjunto es abierto (y también cerrado).
\end{example}

\begin{example}\label{ej:top-cofinita}
Sea $X$ infinito. La colección
\[
  \tau=\{\,U\subseteq X : U=\emptyset\text{ o }X\setminus U\text{ es finito}\,\}
\]
es la \emph{topología cofinita}. La condición (T2) se sigue de que la
intersección de una familia arbitraria de conjuntos finitos es finita; (T3)
de que la unión finita de finitos es finita.
\end{example}

\section{Conjuntos abiertos y cerrados}

\begin{definition}\label{def:abierto-metrico}
En un espacio métrico $(X,d)$, un subconjunto $U\subseteq X$ es \emph{abierto}
si para cada $x\in U$ existe $r>0$ tal que $B(x,r)\subseteq U$.
\end{definition}

La familia de los conjuntos abiertos en un espacio métrico constituye una
topología, denominada la \emph{topología métrica}. Así pues, todo espacio
métrico es un espacio topológico. El recíproco no es cierto: existen espacios
topológicos cuya topología no puede definirse mediante ninguna métrica
(\emph{espacios no metrizables}).

\begin{definition}\label{def:cerrado-top}
Un subconjunto $F\subseteq X$ de un espacio topológico es \emph{cerrado} si
su complementario $X\setminus F$ es abierto.
\end{definition}

\begin{proposition}\label{prop:prop-cerrados-top}
En un espacio topológico $(X,\tau)$ se cumple:
\begin{enumerate}
  \item[(i)] $\emptyset$ y $X$ son cerrados;
  \item[(ii)] la intersección arbitraria de cerrados es cerrada;
  \item[(iii)] la unión finita de cerrados es cerrada.
\end{enumerate}
\end{proposition}

\section{Interior, adherencia y frontera}

\begin{definition}\label{def:interior-top}
El \emph{interior} de un conjunto $A\subseteq X$ es el mayor abierto contenido
en $A$, esto es,
\[
  \operatorname{int}(A)=\bigcup\{\,U\in\tau : U\subseteq A\,\}.
\]
\end{definition}

\begin{definition}\label{def:adherencia}
La \emph{adherencia} (o \emph{clausura}) de $A\subseteq X$ es el menor cerrado
que contiene a $A$, esto es,
\[
  \overline{A}=\bigcap\{\,F\subseteq X : F\text{ cerrado y }A\subseteq F\,\}.
\]
\end{definition}

\begin{definition}\label{def:frontera-top}
La \emph{frontera} (o \emph{borde}) de $A\subseteq X$ se define por
\[
  \partial A = \overline{A}\setminus\operatorname{int}(A).
\]
\end{definition}

\begin{proposition}\label{prop:carac-adherencia}
Un punto $x\in X$ pertenece a $\overline{A}$ si y solo si todo abierto que
contiene a $x$ intersecta a $A$.
\end{proposition}

\section{Bases y subbases}

\begin{definition}\label{def:base-top}
Una familia $\mathcal{B}\subseteq\tau$ es una \emph{base} de la topología
$\tau$ si todo abierto no vacío es unión de elementos de $\mathcal{B}$.
\end{definition}

\begin{example}\label{ej:base-metrica}
En un espacio métrico, la familia de todas las bolas abiertas
$\{B(x,r):x\in X,\,r>0\}$ es una base de la topología métrica.
\end{example}

\begin{definition}\label{def:subbase-top}
Una familia $\mathcal{S}\subseteq\tau$ es una \emph{subbase} de $\tau$ si la
colección de intersecciones finitas de elementos de $\mathcal{S}$ forma una base.
\end{definition}

\section{Conexidad}

\begin{definition}\label{def:conexo-top}
Un espacio topológico $(X,\tau)$ es \emph{conexo} si los únicos subconjuntos
abiertos y cerrados simultáneamente son $\emptyset$ y $X$.
\end{definition}

\begin{proposition}\label{prop:r-conexo}
El espacio $\mathbb{R}$ con la topología usual es conexo.
\end{proposition}

\begin{proposition}\label{prop:imagen-conexa-top}
La imagen continua de un espacio conexo es conexa. Es decir, si $X$ es conexo
y $f\colon X\to Y$ es continua, entonces $f(X)$ es conexo en $Y$.
\end{proposition}

\begin{corollary}[Teorema del valor intermedio]\label{cor:valor-intermedio}
Sea $f\colon[a,b]\to\mathbb{R}$ continua. Si $c$ es un valor comprendido entre
$f(a)$ y $f(b)$, entonces existe $x\in[a,b]$ tal que $f(x)=c$.
\end{corollary}

\section{Ejercicios}

\begin{exercise}
Demuestre que la topología cofinita sobre $\mathbb{N}$ no es metrizable.
Sugerencia: estudie las sucesiones convergentes en esta topología.
\end{exercise}

\begin{exercise}
Sea $A\subseteq X$ en un espacio topológico. Pruebe que
$\operatorname{int}(A)=X\setminus\overline{X\setminus A}$.
\end{exercise}

\begin{exercise}
Muestre que si $X$ es un espacio métrico con la métrica discreta, entonces la
topología inducida es la discreta. Determine la base mínima de esta topología.
\end{exercise}

\begin{exercise}
Sea $f\colon X\to Y$ continua y suprayectiva. Si $X$ es conexo, demuestre que
$Y$ es conexo. Dé un ejemplo en el que $Y$ sea conexo pero $X$ no lo sea.
\end{exercise}

\begin{problem}
Caracterice todos los subconjuntos conexos de $\mathbb{Q}$ con la topología
heredada de $\mathbb{R}$.
\end{problem}

\chapter{Sucesiones}
\label{ch:sucesiones}
\index{Sucesión}\index{Convergencia}\index{Subsucesión}\index{Sucesión de Cauchy}\index{Límite superior}\index{Límite inferior}

\section{Convergencia}

La noción de sucesión convergente es quizá el ejemplo más claro de cómo el rigor del siglo~XIX transformó una intuición milenaria en una definición operativa. Los griegos ya manejaban procesos de aproximación sucesiva; Newton y Leibniz las usaron extensamente; pero fue Augustin-Louis Cauchy, en su \emph{Cours d'Analyse} de 1821, quien formuló la primera definición que se aproxima a la moderna. Sin embargo, la versión definitiva, con la cuantificación explícita sobre $\varepsilon$ y $N$, se debe a Karl Weierstrass y sus lecciones en Berlín a partir de 1861.

\begin{definition}
\label{def:convergencia}
Sea $(X,d)$ un espacio métrico. Una sucesión $(x_n)_{n\in\mathbb{N}}$ en $X$ \emph{converge} a $x\in X$, denotado $x_n\to x$, si
\begin{equation}\label{eq:convergencia}
  \forall\varepsilon>0\ \exists N\in\mathbb{N}\ \forall n\geq N:\quad d(x_n,x)<\varepsilon.
\end{equation}
En tal caso, $x$ es el \emph{límite} de $(x_n)$.
\end{definition}

\begin{theorem}[Unicidad del límite]
\label{teo:unicidad-limite}
En un espacio métrico, el límite de una sucesión convergente, si existe, es único.
\end{theorem}
\begin{proof}
Supongamos $x_n\to x$ y $x_n\to y$. Para todo $\varepsilon>0$, existen $N_1,N_2$ tales que $d(x_n,x)<\varepsilon/2$ para $n\geq N_1$ y $d(x_n,y)<\varepsilon/2$ para $n\geq N_2$. Tomando $n\geq\max\{N_1,N_2\}$, la desigualdad triangular da
\[
  d(x,y)\leq d(x,x_n)+d(x_n,y)<\varepsilon.
\]
Como $\varepsilon>0$ es arbitrario, $d(x,y)=0$, luego $x=y$.
\end{proof}

\begin{proposition}
\label{prop:convergencia-componentes}
En $\mathbb{R}^n$ con la métrica euclidiana, $x^{(k)}\to x$ si y solo si cada coordenada converge: $x^{(k)}_i\to x_i$ para $i=1,\dots,n$.
\end{proposition}

\section{Subsucesiones}

\begin{definition}
\label{def:subsucesion}
Una \emph{subsucesión} de $(x_n)$ es una sucesión de la forma $(x_{n_k})_{k\in\mathbb{N}}$ donde $(n_k)$ es una sucesión estrictamente creciente de índices.
\end{definition}

\begin{proposition}
\label{prop:limite-subsucesion}
Si $x_n\to x$, entonces toda subsucesión de $(x_n)$ converge a $x$.
\end{proposition}

\begin{theorem}[Bolzano--Weierstrass en $\mathbb{R}^n$]
\label{teo:bw-rn}
Toda sucesión acotada en $\mathbb{R}^n$ posee una subsucesión convergente.
\end{theorem}
\begin{proof}[Bosquejo]
Por inducción sobre la dimensión. En $\mathbb{R}$, una sucesión acotada posee una subsucesión monótona (teorema de Weierstrass), la cual converge por completitud. En $\mathbb{R}^n$, se extrae primero una subsucesión convergente en la primera coordenada, luego de esta una subsucesión convergente en la segunda, y así sucesivamente.
\end{proof}

\section{Sucesiones de Cauchy}

\begin{definition}
\label{def:cauchy}
Una sucesión $(x_n)$ en $(X,d)$ es de \emph{Cauchy} si
\begin{equation}\label{eq:cauchy}
  \forall\varepsilon>0\ \exists N\in\mathbb{N}\ \forall m,n\geq N:\quad d(x_m,x_n)<\varepsilon.
\end{equation}
\end{definition}

\begin{proposition}
\label{prop:convergente-cauchy}
Toda sucesión convergente en un espacio métrico es de Cauchy.
\end{proposition}
\begin{proof}
Si $x_n\to x$, dado $\varepsilon>0$ existe $N$ tal que $d(x_n,x)<\varepsilon/2$ para $n\geq N$. Entonces para $m,n\geq N$,
\[
  d(x_m,x_n)\leq d(x_m,x)+d(x,x_n)<\varepsilon.
\]
\end{proof}

El recíproco del \cref{prop:convergente-cauchy} no es cierto en general. La sucesión de aproximaciones racionales de $\sqrt{2}$ es de Cauchy en $\mathbb{Q}$ pero no converge en $\mathbb{Q}$. Los espacios donde sí vale el recíproso son precisamente los \emph{completos}, estudiados en el \cref{ch:completitud}.

\section{Límite superior e inferior}

\begin{definition}
\label{def:limsup-liminf}
Sea $(a_n)$ una sucesión acotada en $\mathbb{R}$. Se definen:
\[
  \limsup_{n\to\infty} a_n=\lim_{n\to\infty}\sup_{k\geq n}a_k,\qquad
  \liminf_{n\to\infty} a_n=\lim_{n\to\infty}\inf_{k\geq n}a_k.
\]
\end{definition}

\begin{proposition}
\label{prop:limsup-liminf-carac}
\begin{enumerate}
  \item[(i)] $\liminf a_n\leq\limsup a_n$.
  \item[(ii)] $a_n\to L$ si y solo si $\liminf a_n=\limsup a_n=L$.
\end{enumerate}
\end{proposition}

\section{Criterios de convergencia}

\begin{theorem}[Criterio de comparación]
\label{teo:comparacion-sucesiones}
Sean $(a_n)$, $(b_n)$ sucesiones en $\mathbb{R}$ con $0\leq a_n\leq b_n$ para todo $n$. Si $b_n\to 0$, entonces $a_n\to 0$.
\end{theorem}

\begin{theorem}[Criterio de la razón]
\label{teo:cociente}
Sea $(a_n)$ una sucesión en $\mathbb{R}$ tal que $\lim_{n\to\infty}|a_{n+1}/a_n|=L$. Si $L<1$, entonces $a_n\to 0$.
\end{theorem}

\begin{theorem}[Criterio de Cauchy para sucesiones reales]
\label{teo:cauchy-reales}
Una sucesión de números reales converge si y solo si es de Cauchy.
\end{theorem}
\begin{proof}
La implicación directa es el \cref{prop:convergente-cauchy}. Para el recíproco, sea $(a_n)$ de Cauchy. La sucesión está acotada, luego por Bolzano--Weierstrass posee una subsucesión convergente $a_{n_k}\to L$. Dado $\varepsilon>0$, sea $N$ tal que $|a_m-a_n|<\varepsilon/2$ para $m,n\geq N$, y $K$ tal que $n_K\geq N$ y $|a_{n_K}-L|<\varepsilon/2$. Entonces para $n\geq N$,
\[
  |a_n-L|\leq |a_n-a_{n_K}|+|a_{n_K}-L|<\varepsilon.
\]
\end{proof}

\begin{exercise}
Demuestre que una sucesión de Cauchy está acotada.
\end{exercise}

\begin{exercise}
Sea $(x_n)$ una sucesión en un espacio métrico tal que $\sum_{n=1}^{\infty}d(x_n,x_{n+1})<\infty$. Pruebe que $(x_n)$ es de Cauchy.
\end{exercise}

\begin{exercise}
Calcule $\limsup$ y $\liminf$ de la sucesión $a_n=(1+(-1)^n)n/(n+1)$.
\end{exercise}

\begin{problem}
Sea $(X,d)$ un espacio métrico. Una sucesión $(x_n)$ se dice \emph{contractiva} si existe $c\in[0,1)$ tal que $d(x_{n+1},x_n)\leq c\,d(x_n,x_{n-1})$ para todo $n\geq 2$. Demuestre que toda sucesión contractiva es de Cauchy y que, si $X$ es completo, converge.
\end{problem}

\chapter{Funciones continuas}
\label{ch:funciones-continuas}
\index{Función continua}\index{Continuidad}\index{Continuidad secuencial}\index{Homeomorfismo}\index{Continuidad uniforme}

\section{Continuidad}

La continuidad es uno de los conceptos más fecundos de toda la matemática, pero también uno de los que más tardó en alcanzar su forma definitiva. Euler, en el siglo~XVIII, la identificaba informalmente con la posibilidad de trazar la gráfica de una función ``sin levantar el lápiz del papel''. Esta imagen geométrica, aunque intuitiva, no permite demostrar teoremas generales. La definición moderna, debida a Bolzano (1817) y Weierstrass (1861), desplaza el énfasis de la curva a la proximidad entre puntos: una función es continua si puntos cercanos en el dominio se aplican en puntos cercanos en el codominio.

\begin{definition}
\label{def:continuidad}
Sean $(X,d_X)$ e $(Y,d_Y)$ espacios métricos. Una función $f\colon X\to Y$ es \emph{continua} en $x_0\in X$ si
\begin{equation}\label{eq:continuidad}
  \forall\varepsilon>0\ \exists\delta>0:\quad d_X(x,x_0)<\delta\Rightarrow d_Y(f(x),f(x_0))<\varepsilon.
\end{equation}
Se dice que $f$ es continua en $X$ cuando lo es en todo punto.
\end{definition}

\begin{example}
\label{ej:polinomios-continuos}
Todo polinomio $p\colon\mathbb{R}\to\mathbb{R}$ es continuo. La función $f(x)=1/x$ es continua en $\mathbb{R}\setminus\{0\}$ pero no admite extensión continua a $\mathbb{R}$.
\end{example}

\section{Continuidad secuencial}

\begin{theorem}[Caracterización secuencial]
\label{teo:continuidad-secuencial}
$f\colon X\to Y$ es continua en $x_0$ si y solo si para toda sucesión $x_n\to x_0$ en $X$ se tiene $f(x_n)\to f(x_0)$ en $Y$.
\end{theorem}
\begin{proof}
$(\Rightarrow)$ Dado $\varepsilon>0$, sea $\delta>0$ como en \eqref{eq:continuidad}. Si $x_n\to x_0$, existe $N$ tal que $d_X(x_n,x_0)<\delta$ para $n\geq N$, luego $d_Y(f(x_n),f(x_0))<\varepsilon$.

$(\Leftarrow)$ Si $f$ no fuera continua en $x_0$, existiría $\varepsilon>0$ tal que para cada $n\in\mathbb{N}$ existe $x_n$ con $d_X(x_n,x_0)<1/n$ pero $d_Y(f(x_n),f(x_0))\geq\varepsilon$. Entonces $x_n\to x_0$ pero $f(x_n)\not\to f(x_0)$.
\end{proof}

\section{Caracterización topológica}

\begin{theorem}
\label{teo:continuidad-topologica}
$f\colon X\to Y$ es continua en $X$ si y solo si la imagen inversa de todo abierto de $Y$ es un abierto de $X$.
\end{theorem}
\begin{proof}
$(\Rightarrow)$ Sea $V\subseteq Y$ abierto y $x_0\in f^{-1}(V)$. Como $V$ es abierto, existe $\varepsilon>0$ con $B(f(x_0),\varepsilon)\subseteq V$. Por continuidad, existe $\delta>0$ tal que $B(x_0,\delta)\subseteq f^{-1}(V)$.

$(\Leftarrow)$ Sea $x_0\in X$ y $\varepsilon>0$. El conjunto $V=B(f(x_0),\varepsilon)$ es abierto, luego $f^{-1}(V)$ es abierto y contiene a $x_0$. Existe $\delta>0$ con $B(x_0,\delta)\subseteq f^{-1}(V)$, lo cual implica \eqref{eq:continuidad}.
\end{proof}

Esta caracterización es la clave de la generalización topológica: en un espacio topológico abstracto, donde no hay métrica, la continuidad se \emph{define} mediante la preservación de abiertos por imagen inversa.

\section{Homeomorfismos}

\begin{definition}
\label{def:homeomorfismo}
Una función $f\colon X\to Y$ es un \emph{homeomorfismo} si es biyectiva, continua, y su inversa $f^{-1}$ es también continua. Cuando existe un homeomorfismo entre $X$ e $Y$, se dice que son \emph{homeomorfos}.
\end{definition}

Los homeomorfismos preservan todas las propiedades topológicas: abiertos, cerrados, clausura, interior, compacidad, conexidad, etc. La topología se ocupa precisamente de clasificar espacios ``salvo homeomorfismo''.

\begin{example}
\label{ej:homeomorfismo-r}
La función $f(x)=\arctan x$ es un homeomorfismo entre $\mathbb{R}$ y $(-\pi/2,\pi/2)$. Sin embargo, $\mathbb{R}$ no es homeomorfo a $[0,1]$ (la compacidad, que estudiaremos en el \cref{ch:compacidad}, es un invariante topológico).
\end{example}

\section{Continuidad uniforme}

\begin{definition}
\label{def:uniforme}
Una función $f\colon(X,d_X)\to(Y,d_Y)$ es \emph{uniformemente continua} si
\begin{equation}\label{eq:uniforme}
  \forall\varepsilon>0\ \exists\delta>0\ \forall x_1,x_2\in X:\quad d_X(x_1,x_2)<\delta\Rightarrow d_Y(f(x_1),f(x_2))<\varepsilon.
\end{equation}
\end{definition}

La diferencia crucial con la continuidad ordinaria es que $\delta$ no depende del punto $x_0$, sino únicamente de $\varepsilon$.

\begin{theorem}[de Heine]
\label{teo:heine}
Toda función continua definida sobre un espacio métrico compacto (\cref{ch:compacidad}) y con valores en otro espacio métrico es uniformemente continua.
\end{theorem}

\begin{proof}[Bosquejo]
Para cada $x\in X$, sea $\delta_x$ como en la definición de continuidad. Las bolas $B(x,\delta_x/2)$ recubren $X$. Por compacidad, existe un subrecubrimiento finito. El mínimo de los radios correspondientes proporciona un $\delta$ uniforme.
\end{proof}

\begin{exercise}
Demuestre que $f(x)=x^2$ es continua en $\mathbb{R}$ pero no uniformemente continua.
\end{exercise}

\begin{exercise}
Sea $f\colon X\to Y$ continua e inyectiva, con $X$ compacto. Demuestre que $f$ es un homeomorfismo sobre su imagen.
\end{exercise}

\begin{exercise}
Pruebe que la composición de funciones uniformemente continuas es uniformemente continua.
\end{exercise}

\begin{problem}
Sea $f\colon\mathbb{R}\to\mathbb{R}$ continua y periódica de período $T>0$. Demuestre que $f$ es uniformemente continua en $\mathbb{R}$.
\end{problem}

\chapter{Compacidad}
\label{ch:compacidad}
\index{Compacidad}\index{Cubierta abierta}\index{Compacidad secuencial}\index{Totalmente acotado}\index{Teorema de Heine--Borel}\index{Teorema de Bolzano--Weierstrass}\index{Punto límite}

\section{Cubiertas abiertas}

La compacidad es la propiedad topológica que generaliza, en espacios métricos arbitrarios, la combinación de ``cerrado y acotado'' que caracteriza a los subconjuntos compactos de $\mathbb{R}^n$. Esta generalización no fue inmediata: durante décadas, los matemáticos trabajaron con nociones de ``acotación'' que resultaban insuficientes en dimensiones infinitas. La definición moderna, basada en cubiertas abiertas, fue introducida por Maurice Fréchet y refinada por Pavel Alexandrov y Pavel Urysohn en la década de 1920.

\begin{definition}
\label{def:cubierta}
Una \emph{cubierta abierta} de un espacio métrico $X$ es una familia $\{U_i\}_{i\in I}$ de abiertos tal que $X=\bigcup_{i\in I}U_i$.
\end{definition}

\begin{definition}
\label{def:compacto}
Un espacio métrico $X$ es \emph{compacto} si de toda cubierta abierta se puede extraer una subcubierta finita, es decir, existen $i_1,\dots,i_n\in I$ tales que $X=U_{i_1}\cup\dots\cup U_{i_n}$.
\end{definition}

\begin{example}
\label{ej:intervalo-cerrado}
El intervalo $[0,1]\subseteq\mathbb{R}$ es compacto. Este es el contenido del \emph{lema de Heine--Borel} para la recta, que generalizaremos más adelante.
\end{example}

\begin{example}
\label{ej:bola-unitaria-l2}
La bola unidad cerrada de $\ell^2$ no es compacta. La base canónica $(e_n)$, donde $e_n$ tiene un $1$ en la posición $n$ y $0$ en las demás, es una sucesión acotada sin subsucesiones convergentes.
\end{example}

\section{Compacidad}

\begin{proposition}
\label{prop:compacto-cerrado}
Todo subespacio compacto de un espacio métrico es cerrado.
\end{proposition}
\begin{proof}
Sea $K\subseteq X$ compacto. Para probar que $K$ es cerrado, mostramos que $X\setminus K$ es abierto. Sea $x\notin K$. Para cada $y\in K$, sean $U_y=B(y,d(x,y)/2)$ y $V_y=B(x,d(x,y)/2)$. Los $U_y$ forman una cubierta abierta de $K$, de la cual extraemos una subcubierta finita $U_{y_1},\dots,U_{y_n}$. Entonces $V=\bigcap_{i=1}^n V_{y_i}$ es un abierto que contiene a $x$ y no intersecta a $K$.
\end{proof}

\begin{proposition}
\label{prop:imagen-compacta}
La imagen continua de un compacto es compacta.
\end{proposition}
\begin{proof}
Sea $f\colon X\to Y$ continua, $K\subseteq X$ compacto, y $\{V_j\}$ una cubierta abierta de $f(K)$. Entonces $\{f^{-1}(V_j)\}$ cubre $K$; extrayendo una subcubierta finita y aplicando $f$, obtenemos una subcubierta finita de $f(K)$.
\end{proof}

\section{Compacidad secuencial}

\begin{definition}
\label{def:compacto-sucesiones}
Un espacio métrico $(X,d)$ es \emph{secuencialmente compacto} si toda sucesión en $X$ posee una subsucesión convergente.
\end{definition}

\begin{theorem}
\label{teo:equivalencia-compacidad}
En un espacio métrico, las siguientes afirmaciones son equivalentes:
\begin{enumerate}
  \item[(i)] $X$ es compacto (por cubiertas).
  \item[(ii)] $X$ es secuencialmente compacto.
  \item[(iii)] $X$ es completo y totalmente acotado.
\end{enumerate}
\end{theorem}
\begin{proof}[Bosquejo]
(i)$\Rightarrow$(ii): Si existiera una sucesión sin subsucesión convergente, cada punto tendría un entorno conteniendo solo finitos términos, proporcionando una cubierta sin subcubierta finita.

(ii)$\Rightarrow$(iii): La completitud se sigue de que toda sucesión de Cauchy con subsucesión convergente converge. La total acotación se prueba por reducción al absurdo: si existiera $\varepsilon>0$ sin recubrimiento finito por $\varepsilon$-bolas, se construiría inductivamente una sucesión con todos sus términos a distancia $\geq\varepsilon$.

(iii)$\Rightarrow$(i): Dada una cubierta abierta, se prueba que existe un $\varepsilon>0$ de Lebesgue tal que toda bola de radio $\varepsilon$ está contenida en algún abierto de la cubierta. La total acotación proporciona un recubrimiento finito por tales bolas.
\end{proof}

\section{Compacidad por punto límite}

\begin{definition}
\label{def:punto-limite}
Un espacio métrico $X$ tiene la \emph{propiedad del punto límite} si todo subconjunto infinito de $X$ posee al menos un punto de acumulación en $X$.
\end{definition}

\begin{proposition}
\label{prop:punto-limite-equivalente}
En espacios métricos, la propiedad del punto límite es equivalente a la compacidad.
\end{proposition}

\section{Totalmente acotado}

\begin{definition}
\label{def:totalmente-acotado}
Un espacio métrico $X$ es \emph{totalmente acotado} si para todo $\varepsilon>0$ existe un recubrimiento finito de $X$ por bolas de radio $\varepsilon$.
\end{definition}

\begin{proposition}
\label{prop:totalmente-acotado-acotado}
Todo espacio totalmente acotado está acotado. En $\mathbb{R}^n$, un conjunto está totalmente acotado si y solo si está acotado.
\end{proposition}

\section{Teorema de Heine--Borel}

\begin{theorem}[Heine--Borel]
\label{teo:heine-borel}
Un subconjunto de $\mathbb{R}^n$ (con la métrica euclidiana) es compacto si y solo si es cerrado y acotado.
\end{theorem}
\begin{proof}
Si $K\subseteq\mathbb{R}^n$ es compacto, es cerrado (\cref{prop:compacto-cerrado}) y acotado (todo compacto en un espacio métrico está acotado). Recíprocamente, si $K$ es cerrado y acotado, está contenido en un cubo $[-M,M]^n$, que es compacto por el teorema de Tychonoff para productos finitos. Como subconjunto cerrado de un compacto, $K$ es compacto.
\end{proof}

\section{Teorema de Bolzano--Weierstrass}

\begin{theorem}[Bolzano--Weierstrass]
\label{teo:bw-compacidad}
Todo subconjunto infinito acotado de $\mathbb{R}^n$ posee un punto de acumulación.
\end{theorem}
\begin{proof}
Sea $A\subseteq\mathbb{R}^n$ infinito y acotado. Su clausura $\overline{A}$ es cerrada y acotada, luego compacta por Heine--Borel. Como subconjunto infinito de un compacto, $A$ posee un punto de acumulación.
\end{proof}

\section{Funciones continuas sobre compactos}

\begin{theorem}[del máximo y mínimo]
\label{teo:max-min}
Sea $X$ compacto y $f\colon X\to\mathbb{R}$ continua. Entonces $f$ alcanza su máximo y su mínimo en $X$.
\end{theorem}
\begin{proof}
Por \cref{prop:imagen-compacta}, $f(X)$ es compacto en $\mathbb{R}$, luego cerrado y acotado. Todo conjunto cerrado y acotado de $\mathbb{R}$ contiene su supremo e ínfimo.
\end{proof}

\begin{exercise}
Demuestre que la intersección arbitraria de compactos es compacta, y que la unión finita de compactos es compacta.
\end{exercise}

\begin{exercise}
Sea $X$ compacto e $Y$ Hausdorff. Demuestre que toda biyección continua $f\colon X\to Y$ es un homeomorfismo.
\end{exercise}

\begin{exercise}
Muestre que $C([a,b])$ con la métrica del supremo no es totalmente acotado.
\end{exercise}

\begin{problem}
Sea $(X,d)$ un espacio métrico. Pruebe que $X$ es compacto si y solo si toda función continua $f\colon X\to\mathbb{R}$ está acotada superiormente.
\end{problem}

\chapter{Conexidad}
\label{ch:conexidad}
\index{Conexidad}\index{Conjunto conexo}\index{Conexión por caminos}\index{Componente conexa}\index{Arco}

\section{Conjuntos conexos}

La conexidad es la propiedad topológica que formaliza la idea intuitiva de que un espacio "está todo en una pieza''. Aunque parezca obvia en contextos geométricos, su formulación exacta requiere cuidado: el intervalo $[0,1]\cup[2,3]$ no está ``en una pieza'', pero ¿cómo capturar esto sin apelar a la intuición visual? La respuesta, debida esencialmente a los primeros topólogos del siglo~XX, consiste en observar que $[0,1]\cup[2,3]$ puede descomponerse en dos abiertos disjuntos no vacíos, mientras que $[0,1]$ no admite tal descomposición.

\begin{definition}
\label{def:conexo}
Un espacio métrico $X$ es \emph{conexo} si los únicos subconjuntos de $X$ que son simultáneamente abiertos y cerrados son $\emptyset$ y $X$. Un subconjunto $A\subseteq X$ es conexo si lo es como subespacio métrico.
\end{definition}

\begin{theorem}
\label{teo:r-conexo}
El espacio $\mathbb{R}$ con la métrica usual es conexo.
\end{theorem}
\begin{proof}
Supongamos que $\mathbb{R}=U\cup V$ con $U,V$ abiertos disjuntos no vacíos. Sea $a\in U$ y $b\in V$; podemos suponer $a<b$. Sea $s=\sup(U\cap[a,b])$. Como $U$ es abierto, $s\notin U$ (de lo contrario existiría una bola alrededor de $s$ contenida en $U$, contradiciendo la definición de supremo). Como $V$ es abierto, $s\notin V$ (pues de lo contrario $s>a$ y una bola alrededor de $s$ evitaría que puntos menores que $s$ y cercanos estuvieran en $U$). Contradicción.
\end{proof}

\begin{theorem}
\label{teo:intervalos-conexos}
Los únicos subconjuntos conexos de $\mathbb{R}$ son los intervalos (abiertos, cerrados, semiabiertos, acotados o no, incluyendo los puntos).
\end{theorem}

\section{Conexión por caminos}

\begin{definition}
\label{def:conexo-por-caminos}
Un espacio métrico $X$ es \emph{conexo por caminos} si para cualesquiera $x,y\in X$ existe una función continua $\gamma\colon[0,1]\to X$ tal que $\gamma(0)=x$ y $\gamma(1)=y$. A tal función se le llama \emph{camino} de $x$ a $y$.
\end{definition}

\begin{proposition}
\label{prop:caminos-implica-conexo}
Todo espacio conexo por caminos es conexo.
\end{proposition}
\begin{proof}
Supongamos $X=U\cup V$ con $U,V$ abiertos disjuntos no vacíos. Sea $x\in U$, $y\in V$ y $\gamma$ un camino de $x$ a $y$. Entonces $\gamma^{-1}(U)$ y $\gamma^{-1}(V)$ son abiertos disjuntos no vacíos de $[0,1]$ cuya unión es $[0,1]$, contradiciendo la conexidad de $[0,1]$.
\end{proof}

El recíproco no es cierto en general. El \emph{peine del topólogo} y la \emph{curva del seno del topólogo} son ejemplos clásicos de espacios conexos pero no conexos por caminos.

\section{Componentes conexas}

\begin{definition}
\label{def:componente}
La \emph{componente conexa} de un punto $x\in X$ es la unión de todos los subconjuntos conexos de $X$ que contienen a $x$.
\end{definition}

\begin{proposition}
\label{prop:componentes-propiedades}
Las componentes conexas de un espacio métrico forman una partición en subconjuntos cerrados, maximalmente conexos y disjuntos dos a dos.
\end{proposition}

\section{Aplicaciones de la conexidad}

\begin{theorem}[del valor intermedio]
\label{teo:valor-intermedio}
Sea $f\colon[a,b]\to\mathbb{R}$ continua. Si $c$ es un valor entre $f(a)$ y $f(b)$, entonces existe $x\in[a,b]$ tal que $f(x)=c$.
\end{theorem}
\begin{proof}
Si $f$ es continua y $[a,b]$ es conexo, entonces $f([a,b])$ es conexo en $\mathbb{R}$, luego es un intervalo. Como contiene a $f(a)$ y $f(b)$, contiene a todo valor intermedio.
\end{proof}

\begin{theorem}
\label{teo:homeomorfismo-intervalos}
Dos intervalos compactos $[a,b]$ y $[c,d]$ son homeomorfos. Sin embargo, $[a,b]$ no es homeomorfo a $[a,b)\cup(b,c]$ (la conexidad es un invariante topológico).
\end{theorem}

\begin{exercise}
Demuestre que la circunferencia $S^1=\{(x,y)\in\mathbb{R}^2:\,x^2+y^2=1\}$ es conexa. Sugerencia: considere la parametrización $t\mapsto(\cos t,\sin t)$.
\end{exercise}

\begin{exercise}
Pruebe que el conjunto de Cantor no es conexo. De hecho, determine sus componentes conexas.
\end{exercise}

\begin{exercise}
Sea $f\colon\mathbb{R}\to\mathbb{R}$ continua tal que $f(x)\in\mathbb{Q}$ para todo $x$. Demuestre que $f$ es constante.
\end{exercise}

\begin{problem}
Sea $X$ conexo por caminos. Demuestre que para todo $n\in\mathbb{N}$, el espacio producto $X^n$ es conexo por caminos.
\end{problem}

\chapter{Completitud}
\label{ch:completitud}
\index{Espacio completo}\index{Completación}\index{Criterio de Cantor}\index{Principio de contracción de Banach}\index{Punto fijo}

\section{Espacios completos}

La completitud es la propiedad que distingue a $\mathbb{R}$ de $\mathbb{Q}$, y que hace posible gran parte del análisis matemático moderno. Sin ella, las sucesiones de Cauchy serían huérfanas de límite, los operadores contractivos no tendrían punto fijo garantizado, y las ecuaciones diferenciales carecerían de solución en general. El concepto fue clarificado por Cauchy en 1821, pero su formulación axiomática y la construcción de completaciones arbitrarias son obra de Cantor, Dedekind y, en el contexto métrico, de Fréchet y Hausdorff.

\begin{definition}
\label{def:espacio-completo}
Un espacio métrico $(X,d)$ es \emph{completo} si toda sucesión de Cauchy en $X$ converge a un punto de $X$.
\end{definition}

\begin{example}
\label{ej:rn-completo}
$\mathbb{R}^n$ con la métrica euclidiana es completo. Este hecho, heredado de la completitud de $\mathbb{R}$, es la base de casi toda el análisis en dimensiones finitas.
\end{example}

\begin{example}
\label{ej:cb-completo}
El espacio $C([a,b])$ con la métrica del supremo es completo. La convergencia en esta métrica es la convergencia uniforme, y el límite uniforme de funciones continuas es continua.
\end{example}

\begin{example}
\label{ej:q-incompleto}
$\mathbb{Q}$ con la métrica usual no es completo. La sucesión de aproximaciones decimales de $\sqrt{2}$ es de Cauchy pero no converge en $\mathbb{Q}$.
\end{example}

\section{Completación de espacios métricos}

\begin{theorem}
\label{teo:completacion}
Todo espacio métrico $(X,d)$ admite una \emph{completación}: existe un espacio métrico completo $(\widetilde{X},\widetilde{d})$ y una isometría $i\colon X\to\widetilde{X}$ tal que $i(X)$ es denso en $\widetilde{X}$. La completación es única salvo isometría.
\end{theorem}
\begin{proof}[Bosquejo]
Se considera el conjunto de todas las sucesiones de Cauchy en $X$, dotado de la métrica $d((x_n),(y_n))=\lim_{n\to\infty}d(x_n,y_n)$. Se identifican sucesiones cuya distancia tiende a cero, obteniendo $\widetilde{X}$. La isometría $i$ envía cada $x\in X$ a la sucesión constante $(x,x,\dots)$.
\end{proof}

\begin{remark}
La construcción de Cantor de $\mathbb{R}$ a partir de $\mathbb{Q}$ es exactamente este procedimiento de completación aplicado al espacio métrico de los racionales.
\end{remark}

\section{Teorema de intersección de Cantor}

\begin{theorem}[de intersección de Cantor]
\label{teo:cantor-interseccion}
Un espacio métrico $(X,d)$ es completo si y solo si toda sucesión $(F_n)$ de conjuntos cerrados no vacíos, encajados ($F_{n+1}\subseteq F_n$) y tales que $\operatorname{diam}(F_n)\to 0$, tiene intersección no vacía $\bigcap_n F_n=\{x^*\}$.
\end{theorem}
\begin{proof}
$(\Rightarrow)$ Elijamos $x_n\in F_n$. La condición $\operatorname{diam}(F_n)\to 0$ implica que $(x_n)$ es de Cauchy; por completitud, $x_n\to x^*$. Como cada $F_n$ es cerrado y contiene a todos los $x_k$ con $k\geq n$, se sigue $x^*\in F_n$ para todo $n$. La unicidad se deduce de que $\operatorname{diam}(\bigcap F_n)\leq\operatorname{diam}(F_n)\to 0$.

$(\Leftarrow)$ Sea $(x_n)$ de Cauchy. Sea $F_n=\overline{\{x_k:\,k\geq n\}}$. Los $F_n$ son cerrados, encajados, y $\operatorname{diam}(F_n)\to 0$. Por hipótesis, $\bigcap F_n=\{x^*\}$, y se verifica que $x_n\to x^*$.
\end{proof}

\section{Principio de contracción de Banach}

\begin{definition}
\label{def:contractiva}
Una función $f\colon X\to X$ es \emph{contractiva} (o de \emph{Lipschitz} con constante $L<1$) si existe $L\in[0,1)$ tal que
\begin{equation}\label{eq:contractiva}
  d(f(x),f(y))\leq L\,d(x,y)\quad\text{para todo }x,y\in X.
\end{equation}
\end{definition}

\begin{theorem}[del punto fijo de Banach]
\label{teo:banach}
Sea $(X,d)$ un espacio métrico completo no vacío. Toda aplicación contractiva $f\colon X\to X$ posee un único punto fijo $x^*\in X$, es decir, $f(x^*)=x^*$. Además, para cualquier $x_0\in X$, la sucesión $x_{n+1}=f(x_n)$ converge a $x^*$, con la estimación de error
\begin{equation}\label{eq:error-banach}
  d(x_n,x^*)\leq\frac{L^n}{1-L}\,d(x_0,x_1).
\end{equation}
\end{theorem}
\begin{proof}
Definimos $x_{n+1}=f(x_n)$. De \eqref{eq:contractiva} se deduce inductivamente $d(x_n,x_{n+1})\leq L^n d(x_0,x_1)$. Para $m>n$,
\[
  d(x_n,x_m)\leq\sum_{k=n}^{m-1}d(x_k,x_{k+1})\leq d(x_0,x_1)\sum_{k=n}^{\infty}L^k=\frac{L^n}{1-L}d(x_0,x_1).
\]
Luego $(x_n)$ es de Cauchy. Por completitud, $x_n\to x^*$. La continuidad de $f$ implica $f(x^*)=\lim f(x_n)=\lim x_{n+1}=x^*$. La unicidad: si $f(x)=x$ y $f(y)=y$, entonces $d(x,y)=d(f(x),f(y))\leq L\,d(x,y)$, luego $d(x,y)=0$.
\end{proof}

\section{Aplicaciones}

\begin{example}[Ecuaciones diferenciales]
\label{ej:edo-banach}
El teorema de existencia y unicidad de Picard--Lindelöf para ecuaciones diferenciales ordinarias se prueba transformando el problema de valor inicial $y'=F(t,y)$, $y(t_0)=y_0$ en una ecuación integral y definiendo un operador de Picard sobre un espacio $C([t_0-\delta,t_0+\delta])$ adecuado. Bajo una condición de Lipschitz en $F$, este operador es contractivo para $\delta$ suficientemente pequeño.
\end{example}

\begin{example}[Método de Newton]
\label{ej:newton}
El método iterativo $x_{n+1}=x_n-f(x_n)/f'(x_n)$ para hallar raíces de $f$ puede interpretarse como una aplicación de Banach en un entorno adecuado de la raíz, siempre que $f'$ no se anule y ciertas cotas de convexidad se cumplan.
\end{example}

\begin{exercise}
Demuestre que toda función lipschitziana es uniformemente continua, pero que el recíproco no es cierto. Sugerencia: considere $f(x)=\sqrt{x}$ en $[0,1]$.
\end{exercise}

\begin{exercise}
Sea $X$ completo y $f\colon X\to X$ tal que $f^k$ (la $k$-ésima iterada) es contractiva para algún $k\in\mathbb{N}$. Demuestre que $f$ posee un único punto fijo.
\end{exercise}

\begin{exercise}
Sea $K\subseteq\mathbb{R}^n$ compacto y $f\colon K\to K$ contractiva. Demuestre que la convergencia de $x_{n+1}=f(x_n)$ es exponencialmente rápida, y estime el número de iteraciones necesarias para alcanzar una precisión $\varepsilon$.
\end{exercise}

\begin{problem}
Sea $(X,d)$ completo y $f\colon X\to X$ tal que $d(f(x),f(y))<d(x,y)$ para todo $x\neq y$. ¿Tiene $f$ necesariamente un punto fijo? Si no, dé un contraejemplo.
\end{problem}

\chapter{Series}
\label{ch:series}
\index{Serie}\index{Convergencia absoluta}\index{Convergencia condicional}\index{Criterios de convergencia}\index{Serie de potencias}\index{Convergencia uniforme}

\section{Series numéricas}

Las series infinitas aparecieron de manera natural en el cálculo de Newton y Leibniz como sumas de infinitos términos, pero su manejo riguroso fue una de las motivaciones centrales del análisis del siglo~XIX. Series como la armónica $\sum 1/n$, que ``crece sin bounds'' aunque sus términos tiendan a cero, o la serie de Grandi $\sum (-1)^n$, cuya suma parecía ser $0$, $1$ o $1/2$ según el método de agrupación, mostraban que la intuición finita no podía extenderse al infinito sin controles explícitos. Fue Cauchy quien, en su \emph{Cours d'Analyse}, estableció la definición moderna: una serie converge si la sucesión de sus sumas parciales converge.

\begin{definition}
\label{def:serie}
Dada una sucesión $(a_n)$ en $\mathbb{R}$ (o en $\mathbb{C}$), la \emph{serie} $\sum_{n=1}^{\infty}a_n$ es la sucesión de \emph{sumas parciales}
\[
  S_N=\sum_{n=1}^{N}a_n.
\]
Se dice que la serie \emph{converge} si $(S_N)$ converge; en tal caso, su límite se denota $\sum_{n=1}^{\infty}a_n$.
\end{definition}

\begin{theorem}[Condición necesaria de convergencia]
\label{teo:cond-necesaria}
Si $\sum a_n$ converge, entonces $a_n\to 0$.
\end{theorem}
\begin{proof}
Si $S_N\to S$, entonces $a_N=S_N-S_{N-1}\to S-S=0$.
\end{proof}

\begin{example}
\label{ej:armonica}
La serie armónica $\sum_{n=1}^{\infty}1/n$ diverge, aunque $1/n\to 0$. Esto muestra que la condición necesaria no es suficiente.
\end{example}

\section{Series absolutamente convergentes}

\begin{definition}
\label{def:absoluta}
Una serie $\sum a_n$ converge \emph{absolutamente} si $\sum|a_n|$ converge.
\end{definition}

\begin{theorem}
\label{teo:absoluta-implica-convergencia}
Toda serie absolutamente convergente es convergente.
\end{theorem}
\begin{proof}
Si $\sum|a_n|$ converge, entonces las sumas parciales de $\sum|a_n|$ forman una sucesión de Cauchy. Dado $\varepsilon>0$, existe $N$ tal que $\sum_{k=m+1}^{n}|a_k|<\varepsilon$ para $n>m\geq N$. Entonces $|\sum_{k=m+1}^{n}a_k|\leq\sum_{k=m+1}^{n}|a_k|<\varepsilon$, luego las sumas parciales de $\sum a_n$ son de Cauchy y convergen.
\end{proof}

\section{Series condicionalmente convergentes}

\begin{definition}
\label{def:condicional}
Una serie convergente $\sum a_n$ es \emph{condicionalmente convergente} si $\sum|a_n|$ diverge.
\end{definition}

\begin{theorem}[Riemann]
\label{teo:riemann}
Si $\sum a_n$ es condicionalmente convergente, entonces para todo $L\in\mathbb{R}\cup\{\pm\infty\}$ existe una reordenación $(a_{\sigma(n)})$ tal que $\sum a_{\sigma(n)}=L$.
\end{theorem}

Este teorema, publicado por Bernhard Riemann en 1854, es uno de los resultados más sorprendentes del análisis clásico: demuestra que la conmutatividad de la suma, evidente en cantidades finitas, se pierde radicalmente en el infinito si la convergencia no es absoluta.

\section{Criterios de convergencia}

\begin{theorem}[Criterio de comparación]
\label{teo:comparacion-series}
Sean $\sum a_n$ y $\sum b_n$ series de términos no negativos con $a_n\leq b_n$ para todo $n$. Si $\sum b_n$ converge, entonces $\sum a_n$ converge.
\end{theorem}

\begin{theorem}[Criterio de la razón]
\label{teo:cociente-series}
Sea $\sum a_n$ con $a_n\neq 0$. Si $\lim_{n\to\infty}|a_{n+1}/a_n|=L$, entonces:
\begin{itemize}
  \item Si $L<1$, la serie converge absolutamente.
  \item Si $L>1$, la serie diverge.
  \item Si $L=1$, el criterio no decide.
\end{itemize}
\end{theorem}

\begin{theorem}[Criterio de la raíz]
\label{teo:raiz}
Sea $\sum a_n$. Si $\limsup_{n\to\infty}\sqrt[n]{|a_n|}=L$, entonces la serie converge absolutamente si $L<1$ y diverge si $L>1$.
\end{theorem}

\begin{theorem}[Criterio de Dirichlet]
\label{teo:dirichlet}
Sean $(a_n)$ y $(b_n)$ sucesiones reales tales que las sumas parciales de $\sum a_n$ están acotadas, $b_n\to 0$ monótonamente. Entonces $\sum a_nb_n$ converge.
\end{theorem}

\section{Series de potencias}

\begin{definition}
\label{def:serie-potencias}
Una \emph{serie de potencias} centrada en $x_0$ es una serie de la forma
\[
  \sum_{n=0}^{\infty}c_n(x-x_0)^n,
\]
donde $(c_n)$ es una sucesión de coeficientes.
\end{definition}

\begin{theorem}
\label{teo:radio-convergencia}
Para toda serie de potencias existe un número $R\in[0,+\infty]$, llamado \emph{radio de convergencia}, tal que la serie converge absolutamente si $|x-x_0|<R$ y diverge si $|x-x_0|>R$.
\end{theorem}
\begin{proof}
Sea $R=1/\limsup_{n\to\infty}\sqrt[n]{|c_n|}$ (con las convenciones $1/0=+\infty$, $1/\infty=0$). El criterio de la raíz da el resultado.
\end{proof}

\section{Convergencia uniforme}

\begin{definition}
\label{def:convergencia-uniforme}
Una serie de funciones $\sum f_n(x)$ converge \emph{uniformemente} en un conjunto $A$ si la sucesión de sumas parciales converge uniformemente en $A$.
\end{definition}

\begin{theorem}[Weierstrass $M$]
\label{teo:weierstrass-m}
Si $|f_n(x)|\leq M_n$ para todo $x\in A$ y $\sum M_n$ converge, entonces $\sum f_n$ converge absoluta y uniformemente en $A$.
\end{theorem}

\section{Intercambio de límites}

\begin{theorem}
\label{teo:intercambio}
Si $\sum f_n$ converge uniformemente en $[a,b]$ y cada $f_n$ es continua, entonces la función suma $S(x)=\sum f_n(x)$ es continua. Además,
\[
  \int_a^b\sum_{n=1}^{\infty}f_n(x)\,dx=\sum_{n=1}^{\infty}\int_a^bf_n(x)\,dx.
\]
\end{theorem}

\begin{exercise}
Demuestre que $\sum_{n=1}^{\infty}1/n^2$ converge. Calcule su valor (sugerencia: relación con la función zeta de Riemann).
\end{exercise}

\begin{exercise}
Determine el radio de convergencia de $\sum_{n=0}^{\infty}n!x^n$ y de $\sum_{n=0}^{\infty}x^n/n!$.
\end{exercise}

\begin{exercise}
Dé una reordenación de $\sum_{n=1}^{\infty}(-1)^{n+1}/n$ que converja a $0$.
\end{exercise}

\begin{problem}
Sea $\sum a_n$ una serie de números reales. Pruebe que si $\sum|a_n|$ diverge, entonces existe una sucesión $(b_n)$ acotada tal que $\sum a_nb_n$ diverge.
\end{problem}

\chapter{Grandes teoremas del análisis}
\label{ch:grandes-teoremas}
\index{Principio de contracción de Banach}\index{Categoría de Baire}\index{Teorema de Ascoli--Arzelà}\index{Teorema de Stone--Weierstrass}\index{Teorema de Dini}\index{Lema de Urysohn}\index{Teorema de Tietze}

\section{Principio de contracción de Banach}

Ya demostrado en el \cref{ch:completitud}, el teorema del punto fijo de Banach merece ser recordado aquí como uno de los grandes teoremas estructurales del análisis. No es un resultado de cálculo: no calcula nada, sino que \emph{garantiza} la existencia y unicidad de una solución bajo condiciones cualitativas. Es el paradigma de cómo el rigor transforma la matemática aplicada.

\begin{theorem}[Banach, recapitulación]
\label{teo:banach-recap}
Sea $(X,d)$ completo no vacío. Toda contracción $f\colon X\to X$ tiene un único punto fijo, y las iteradas $x_{n+1}=f(x_n)$ convergen a él desde cualquier punto inicial.
\end{theorem}

\section{Teorema de la Categoría de Baire}

El teorema de Baire, publicado por René-Louis Baire en 1899, es uno de los resultados más profundos y útiles del análisis. Establece que un espacio completo ``no puede ser cubierto por polvo'': no es union numerable de conjuntos ``magros''. Esta aparente sutileza tiene consecuencias espectaculares: la existencia de funciones continuas no derivables en ningún punto, la no completitud de ciertos espacios de funciones suaves, y el teorema de la función abierta en análisis funcional.

\begin{definition}
\label{def:nowhere-dense}
Un subconjunto $A\subseteq X$ es \emph{nowhere dense} (de primera categoría en sí mismo) si $\operatorname{int}(\overline{A})=\emptyset$.
\end{definition}

\begin{definition}
\label{def:baire-space}
Un espacio métrico $X$ es un \emph{espacio de Baire} si la intersección de toda familia numerable de abiertos densos es densa en $X$.
\end{definition}

\begin{theorem}[Baire]
\label{teo:baire}
Todo espacio métrico completo es un espacio de Baire.
\end{theorem}
\begin{proof}[Bosquejo]
Sea $\{U_n\}$ una sucesión de abiertos densos y $V$ un abierto no vacío. Como $U_1$ es denso, $V\cap U_1$ contiene una bola cerrada $B_1$ de radio $r_1<1$. Como $U_2$ es denso, $B_1^{\circ}\cap U_2$ contiene una bola cerrada $B_2$ de radio $r_2<1/2$, etc. Las bolas $B_n$ son encajadas con diámetros tendiendo a cero. Por completitud, $\bigcap B_n=\{x^*\}$, y $x^*\in V\cap\bigcap U_n$.
\end{proof}

\begin{theorem}[Conjuntos residuales]
\label{teo:residual}
En un espacio de Baire, un conjunto que contiene una intersección numerable de abiertos densos es \emph{residual} (o \emph{co-magro}). Su complemento es de primera categoría. Los conjuntos residuales son densos.
\end{theorem}

\section{Teorema de Ascoli--Arzelà}

La compacidad en espacios de funciones es más sutil que en espacios finito-dimensionales. El teorema de Ascoli--Arzelà, que demostraremos en el \cref{ch:familias-funciones}, caracteriza exactamente cuándo una familia de funciones continuas tiene la propiedad de que toda sucesión posee una subsucesión convergente. La clave está en la combinación de \emph{acotación uniforme} y \emph{equicontinuidad}.

\begin{theorem}[Ascoli--Arzelà, enunciado]
\label{teo:ascoli-arzela-enunciado}
Sea $X$ un espacio métrico compacto. Una familia $\mathcal{F}\subseteq C(X)$ es relativamente compacta en la métrica del supremo si y solo si está uniformemente acotada y equicontinua.
\end{theorem}

\section{Teorema de Stone--Weierstrass}

El teorema de aproximación de Weierstrass (1885) establece que todo función continua en un intervalo compacto puede aproximarse uniformemente por polinomios. Marshall Stone generalizó este resultado en 1937 a álgebras de funciones sobre espacios compactos Hausdorff, mostrando que la propiedad clave no es la forma polinomial, sino la capacidad de separar puntos.

\begin{theorem}[Weierstrass clásico]
\label{teo:weierstrass}
Para toda $f\in C([a,b])$ y todo $\varepsilon>0$, existe un polinomio $p$ tal que $\sup_{x\in[a,b]}|f(x)-p(x)|<\varepsilon$.
\end{theorem}

\begin{theorem}[Stone--Weierstrass, enunciado]
\label{teo:stone-weierstrass-enunciado}
Sea $X$ compacto Hausdorff y $\mathcal{A}\subseteq C(X)$ una subálgebra que separa puntos y contiene las constantes. Entonces $\mathcal{A}$ es densa en $C(X)$ con la métrica del supremo.
\end{theorem}

\section{Teorema de Dini}

\begin{theorem}[Dini]
\label{teo:dini}
Sea $X$ compacto y $(f_n)$ una sucesión monótona (creciente o decreciente) de funciones continuas que converge puntualmente a una función continua $f$. Entonces la convergencia es uniforme.
\end{theorem}
\begin{proof}[Bosquejo]
Supongamos $f_n\downarrow f$. Sea $g_n=f_n-f\geq 0$, continuas, decrecientes a cero puntualmente. Dado $\varepsilon>0$, los conjuntos $K_n=\{x\in X:\,g_n(x)\geq\varepsilon\}$ son compactos, encajados, con intersección vacía. Por la propiedad de intersección finita, algún $K_N$ es vacío, luego $g_n(x)<\varepsilon$ para todo $n\geq N$ y todo $x$.
\end{proof}

\section{Lema de Urysohn (introducción)}

\begin{theorem}[Urysohn, enunciado]
\label{teo:urysohn}
Sea $X$ un espacio topológico normal (es decir, $T_1$ y donde cualesquiera dos cerrados disjuntos pueden separarse por abiertos disjuntos). Para cualesquiera cerrados disjuntos $A,B\subseteq X$, existe una función continua $f\colon X\to[0,1]$ tal que $f|_A=0$ y $f|_B=1$.
\end{theorem}

Este lema, fundamental en topología general, muestra que la separación de cerrados por abiertos implica la separación por funciones continuas. Su demostración en espacios métricos es mucho más sencilla que en el caso general.

\section{Teorema de extensión de Tietze (introducción)}

\begin{theorem}[Tietze, enunciado]
\label{teo:tietze}
Sea $X$ un espacio topológico normal y $A\subseteq X$ cerrado. Toda función continua $f\colon A\to\mathbb{R}$ admite una extensión continua $\widetilde{f}\colon X\to\mathbb{R}$.
\end{theorem}

\begin{remark}
Los teoremas de Urysohn y Tietze serán revisitados en el \cref{ch:topologia-general}, donde se demostrarán en el contexto de espacios topológicos normales.
\end{remark}

\begin{exercise}
Use el teorema de Baire para probar que el conjunto de funciones continuas no derivables en ningún punto de $[0,1]$ es residual en $C([0,1])$ con la métrica del supremo.
\end{exercise}

\begin{exercise}
Demuestre que si $X$ es un espacio métrico completo sin puntos aislados, entonces $X$ no es numerable.
\end{exercise}

\begin{exercise}
Sea $f_n(x)=x^n$ en $[0,1]$. ¿Aplica el teorema de Dini? ¿Por qué la convergencia no es uniforme en $[0,1]$?
\end{exercise}

\begin{problem}
Demuestre el teorema de Urysohn en espacios métricos construyendo explícitamente $f(x)=d(x,A)/(d(x,A)+d(x,B))$.
\end{problem}

\chapter{Espacios separables y espacios completos}
\label{ch:espacios-separables}
\index{Espacio separable}\index{Conjunto denso}\index{Espacio de Baire}\index{Segunda categoría}

\section{Espacios separables}

La separabilidad es una propiedad de ``tamaño controlable'': un espacio es separable si contiene un subconjunto numerable que aproxima a todos sus puntos. Esta condición, aparentemente técnica, tiene consecuencias profundas: en espacios separables, toda sucesión acotada posee una subsucesión débilmente convergente (teorema de Eberlein--Šmulian); las medidas de probabilidad pueden aproximarse por combinaciones convexas de deltas de Dirac; y muchos espacios de funciones usuales son separables, lo cual facilita su implementación computacional.

\begin{definition}
\label{def:separable}
Un espacio métrico $(X,d)$ es \emph{separable} si posee un subconjunto denso numerable.
\end{definition}

\begin{example}
\label{ej:rn-separable}
$\mathbb{R}^n$ es separable, pues $\mathbb{Q}^n$ es denso y numerable.
\end{example}

\begin{example}
\label{ej:cb-separable}
$C([a,b])$ con la métrica del supremo es separable. Los polinomios con coeficientes racionales forman un conjunto denso numerable (por el teorema de Weierstrass, \cref{teo:weierstrass-aprox}).
\end{example}

\begin{example}
\label{ej:linf-no-separable}
El espacio $\ell^{\infty}$ de las sucesiones acotadas con la métrica del supremo no es separable. La familia de sucesiones de ceros y unos es no numerable y cualquier dos elementos distintos están a distancia $1$.
\end{example}

\section{Conjuntos densos numerables}

\begin{theorem}
\label{teo:base-numerable}
Un espacio métrico es separable si y solo si su topología posee una base numerable.
\end{theorem}
\begin{proof}
Si $\{x_n\}$ es denso, las bolas $B(x_n,q)$ con $q\in\mathbb{Q}^+$ forman una base numerable. Recíprocamente, si $\{U_n\}$ es una base numerable, eligiendo un punto $x_n\in U_n$ en cada elemento no vacío obtenemos un conjunto denso numerable.
\end{proof}

\begin{corollary}
\label{cor:lindelof}
Todo espacio métrico separable es \emph{Lindelöf}: de todo recubrimiento abierto se puede extraer un subrecubrimiento numerable.
\end{corollary}

\section{Espacios completamente métricos}

\begin{definition}
\label{def:completamente-metrico}
Un espacio métrico $(X,d)$ es \emph{completamente métrico} si admite una métrica completa que genere la misma topología. (Nótese que no exigimos que la métrica original sea completa.)
\end{definition}

\begin{theorem}[Alexandrov]
\label{teo:alexandrov}
Un subespacio de un espacio métrico completo es completamente métrico si y solo si es un $G_\delta$ (intersección numerable de abiertos).
\end{theorem}

\section{Espacios de Baire}

Ya definidos en el \cref{ch:grandes-teoremas}, los espacios de Baire merecen un estudio más detallado.

\begin{theorem}[Baire, recapitulación]
\label{teo:baire-recap}
Todo espacio métrico completo es de Baire. La intersección de abiertos densos es densa; ningún espacio de Baire es union numerable de conjuntos nowhere dense.
\end{theorem}

\begin{proposition}
\label{prop:abierto-denso-baire}
En un espacio de Baire, todo abierto no vacío es de segunda categoría en sí mismo (es decir, no es union numerable de nowhere dense).
\end{proposition}

\section{Propiedades fundamentales}

\begin{theorem}
\label{teo:separable-herencia}
Todo subespacio de un espacio métrico separable es separable.
\end{theorem}
\begin{proof}
Si $D$ es denso numerable en $X$ y $Y\subseteq X$, para cada $x\in D$ y $n\in\mathbb{N}$ con $B(x,1/n)\cap Y\neq\emptyset$, elegimos $y_{x,n}\in B(x,1/n)\cap Y$. El conjunto de tales $y_{x,n}$ es denso en $Y$.
\end{proof}

\begin{theorem}
\label{teo:producto-separable}
El producto numerable de espacios métricos separables es separable.
\end{theorem}

\section{Ejemplos importantes}

\begin{example}[Espacio de Hilbert L2]
\label{ej:l2-separable}
El espacio de Lebesgue $L^2([0,1])$ es separable. Los polinomios trigonométricos con coeficientes racionales son densos. En efecto, por el teorema de Stone--Weierstrass (\cref{teo:stone-weierstrass}) aplicado a funciones continuas sobre la circunferencia $[0,2\pi]$ con $0$ y $2\pi$ identificados, el álgebra generada por $\{1,\cos nx,\sin nx:\,n\in\mathbb{N}\}$ es densa en $C([0,2\pi])$ con la métrica del supremo; por densidad, también lo es en $L^2([0,2\pi])$ con la norma cuadrática.
\end{example}

\begin{example}[Espacio de Cantor]
\label{ej:cantor-separable}
El conjunto de Cantor, como subespacio de $\mathbb{R}$, es compacto, perfecto, totalmente disconexo y separable (de hecho, es homeomorfo al producto numerable $\{0,1\}^{\mathbb{N}}$).
\end{example}

\begin{exercise}
Demuestre que todo espacio métrico compacto es separable.
\end{exercise}

\begin{exercise}
Sea $X$ separable y $f\colon X\to Y$ sobreyectiva continua. Demuestre que $Y$ es separable.
\end{exercise}

\begin{exercise}
Pruebe que $\ell^p$ es separable para $1\leq p<\infty$, pero $\ell^{\infty}$ no lo es.
\end{exercise}

\begin{problem}
Un espacio métrico se dice \emph{polaco} si es separable y completamente métrico. Demuestre que $\mathbb{R}\setminus\mathbb{Q}$ es polaco, pero no es completo con la métrica usual.
\end{problem}

\chapter{Familias de funciones}
\label{ch:familias-funciones}
\index{Equicontinuidad}\index{Acotación uniforme}\index{Teorema de Ascoli--Arzelà}\index{Espacio normado}\index{Dualidad}

\section{Equicontinuidad}

Cuando estudiamos una familia de funciones, la continuidad individual de cada miembro no garantiza ningún comportamiento uniforme de la familia como conjunto. La noción de equicontinuidad, introducida por Ascoli en 1883 y perfeccionada por Arzelà en 1889, precisamente impone que la elección de $\delta$ en la definición de continuidad pueda hacerse de manera uniforme para toda la familia simultáneamente.

\begin{definition}
\label{def:equicontinuidad}
Una familia $\mathcal{F}$ de funciones de $(X,d_X)$ en $(Y,d_Y)$ es \emph{equicontinua} en $x_0\in X$ si para todo $\varepsilon>0$ existe $\delta>0$ tal que
\[
  d_X(x,x_0)<\delta\quad\Longrightarrow\quad d_Y(f(x),f(x_0))<\varepsilon\quad\text{para todo }f\in\mathcal{F}.
\]
Se dice que $\mathcal{F}$ es equicontinua en $X$ si lo es en cada punto.
\end{definition}

\begin{example}
\label{ej:equicontinua}
La familia de funciones $\{f_n(x)=\sin(nx)/n:\,n\in\mathbb{N}\}$ en $C([0,2\pi])$ es equicontinua, pues $|f_n(x)-f_n(y)|\leq |x-y|$ para todo $n$.
\end{example}

\begin{example}
\label{ej:no-equicontinua}
La familia $\{f_n(x)=x^n:\,n\in\mathbb{N}\}$ en $C([0,1])$ no es equicontinua en $x=1$. Para $\varepsilon=1/2$, cualquier $\delta>0$ permite encontrar $n$ suficientemente grande tal que $|1^n-(1-\delta/2)^n|>1/2$.
\end{example}

\section{Acotación uniforme}

\begin{definition}
\label{def:acotacion-uniforme}
Una familia $\mathcal{F}\subseteq C(X,Y)$ está \emph{uniformemente acotada} si existe $M>0$ tal que $d_Y(f(x),g(x))\leq M$ para todo $f,g\in\mathcal{F}$ y todo $x\in X$. (En particular, cada $f\in\mathcal{F}$ tiene rango acotado por una constante común.)
\end{definition}

\section{Compacidad en espacios de funciones}

El teorema de Ascoli--Arzelà responde a la pregunta: ¿cuándo una familia de funciones continuas se comporta como un conjunto compacto?

\begin{theorem}[Ascoli--Arzelà]
\label{teo:ascoli-arzela}
Sea $X$ un espacio métrico compacto. Una familia $\mathcal{F}\subseteq C(X,\mathbb{R}^n)$ es relativamente compacta en la métrica del supremo si y solo si está puntualmente acotada y equicontinua.
\end{theorem}
\begin{proof}[Demostración completa]
$(\Rightarrow)$ Si $\overline{\mathcal{F}}$ es compacta, entonces está acotada (compacto implica acotado) y equicontinua. La equicontinuidad se prueba por reducción al absurdo: si no fuera equicontinua en algún punto, existirían $\varepsilon>0$, puntos $x_n\to x_0$ y funciones $f_n\in\mathcal{F}$ con $|f_n(x_n)-f_n(x_0)|\geq\varepsilon$. Por compactidad, $f_n$ tiene una subsucesión convergente $f_{n_k}\to f$ uniformemente, lo cual contradice la continuidad uniforme de $f$.

$(\Leftarrow)$ Sea $(f_n)$ una sucesión en $\mathcal{F}$. Como $X$ es compacto, existe un subconjunto denso numerable $\{x_k\}$ (por separabilidad). Por acotación puntual, la sucesión $(f_n(x_1))$ está acotada, luego tiene una subsucesión convergente. Por diagonalización, extraemos una subsucesión $(f_{n_j})$ que converge puntualmente en todo $\{x_k\}$. La equicontinuidad permite extender la convergencia puntual a convergencia uniforme en $X$ mediante un argumento de $\varepsilon/3$.
\end{proof}

\section{Aplicaciones del teorema de Ascoli--Arzelà}

\begin{example}[Existencia en ecuaciones diferenciales]
\label{ej:edo-ascoli}
En la prueba del teorema de Peano de existencia de soluciones para ecuaciones diferenciales sin hipótesis de Lipschitz, se construye una sucesión de aproximaciones de Euler. La equicontinuidad se deduce de la acotación de la derivada, y el teorema de Ascoli--Arzelà proporciona una subsucesión convergente cuyo límite es la solución buscada.
\end{example}

\begin{example}[Teorema de Montel]
\label{ej:montel}
En análisis complejo, el teorema de Montel establece que una familia de funciones holomorfas localmente acotadas es normal (relativamente compacta en la topología de la convergencia uniforme sobre compactos). Es esencialmente el teorema de Ascoli--Arzelà combinado con la estimación de Cauchy.
\end{example}

\section{Recordatorio de espacios normados y dualidad}

El teorema de Ascoli--Arzelà estudia compacidad en espacios de funciones continuas con la métrica del supremo, que es un caso particular de espacio normado. Para la última sección de este capítulo necesitamos los conceptos básicos de norma y dualidad, que se desarrollan en el \cref{ch:analisis-funcional}; los recordamos aquí de forma autocontenida.

\begin{definition}
\label{def:norma-recordatorio}
Una \emph{norma} sobre un espacio vectorial real $E$ es una función $\|\cdot\|\colon E\to\mathbb{R}$ tal que:
\begin{enumerate}
  \item[(i)] $\|x\|\geq 0$, y $\|x\|=0\Leftrightarrow x=0$;
  \item[(ii)] $\|\lambda x\|=|\lambda|\,\|x\|$ para todo $\lambda\in\mathbb{R}$;
  \item[(iii)] $\|x+y\|\leq\|x\|+\|y\|$.
\end{enumerate}
Al par $(E,\|\cdot\|)$ se le llama \emph{espacio normado}.
\end{definition}

\begin{definition}
\label{def:dual-recordatorio}
El \emph{dual topológico} $E'$ de un espacio normado $E$ es el espacio de todos los funcionales lineales continuos $\phi\colon E\to\mathbb{R}$. La norma en $E'$ es $\|\phi\|=\sup_{\|x\|\leq 1}|\phi(x)|$.
\end{definition}

\section{Convergencia de familias de funciones}

\begin{definition}
\label{def:convergencia-debil}
Sea $X$ un espacio normado. Una sucesión $(f_n)$ en el espacio dual $X'$ converge \emph{débilmente*} a $f$ si $f_n(x)\to f(x)$ para todo $x\in X$.
\end{definition}

\begin{theorem}[Banach--Alaoglu, enunciado]
\label{teo:banach-alaoglu}
La bola unidad cerrada del dual de un espacio normado es compacta en la topología débil*.
\end{theorem}

\begin{exercise}
Demuestre que si $\mathcal{F}\subseteq C([a,b])$ es una familia de funciones lipschitzianas con la misma constante $L$, entonces $\mathcal{F}$ es equicontinua.
\end{exercise}

\begin{exercise}
Sea $\mathcal{F}=\{f\in C^1([0,1]):\,|f'(x)|\leq 1\text{ para todo }x\}$. Demuestre que $\mathcal{F}$ es relativamente compacta en $C([0,1])$.
\end{exercise}

\begin{exercise}
Use el teorema de Ascoli--Arzelà para probar que toda sucesión de funciones holomorfas uniformemente acotadas en un disco posee una subsucesión que converge uniformemente sobre compactos.
\end{exercise}

\begin{problem}
Sea $X$ compacto y $\mathcal{F}\subseteq C(X)$ equicontinua. Demuestre que la clausura de $\mathcal{F}$ en la topología de la convergencia puntual coincide con su clausura en la topología del supremo.
\end{problem}

\chapter{Funciones de variación acotada}
\label{ch:variacion-acotada}
\index{Variación total}\index{Función de variación acotada}\index{Función absolutamente continua}\index{Conjunto de medida cero}\index{Casi en todas partes}

\section{Variación total}

La noción de variación acotada, introducida por Camille Jordan en 1881, nació del intento de caracterizar aquellas funciones para las cuales la gráfica tiene ``longitud finita''. Jordan observó que una función continua podía oscilar infinitamente en un intervalo compacto, como la función seno del topólogo, y que el cálculo elemental de longitud de arco no era aplicable sin un control global de las oscilaciones.

\begin{definition}
\label{def:variacion}
Sea $f\colon[a,b]\to\mathbb{R}$. Dada una partición $P=\{a=x_0<x_1<\dots<x_n=b\}$, se define la \emph{variación} de $f$ sobre $P$ como
\[
  V(f,P)=\sum_{i=1}^{n}|f(x_i)-f(x_{i-1})|.
\]
La \emph{variación total} de $f$ en $[a,b]$ es
\[
  V_a^b(f)=\sup\{V(f,P):\,P\text{ partición de }[a,b]\}.
\]
Se dice que $f$ es de \emph{variación acotada} si $V_a^b(f)<\infty$.
\end{definition}

\begin{example}
\label{ej:monotona-va}
Toda función monótona en $[a,b]$ es de variación acotada, y $V_a^b(f)=|f(b)-f(a)|$.
\end{example}

\begin{example}
\label{ej:lipschitz-va}
Toda función lipschitziana es de variación acotada.
\end{example}

\begin{example}
\label{ej:no-va}
La función $f(x)=x\sin(1/x)$ para $x\neq 0$, $f(0)=0$, no es de variación acotada en $[0,1]$.
\end{example}

\section{Funciones monótonas}

\begin{theorem}
\label{teo:monotona-discontinuidades}
Toda función monótona $f\colon[a,b]\to\mathbb{R}$ tiene a lo sumo una cantidad numerable de discontinuidades, todas de salto.
\end{theorem}
\begin{proof}
En cada punto de discontinuidad $x$, el salto $f(x^+)-f(x^-)$ es positivo. Como la suma de los saltos no puede exceder $|f(b)-f(a)|$, solo puede haber una cantidad finita de saltos mayores que $1/n$ para cada $n$. La union numerable de conjuntos finitos es numerable.
\end{proof}

\section{Funciones de variación acotada}

\begin{theorem}[Jordan]
\label{teo:jordan}
Una función $f\colon[a,b]\to\mathbb{R}$ es de variación acotada si y solo si puede escribirse como diferencia de dos funciones crecientes: $f=g-h$, con $g,h$ crecientes.
\end{theorem}
\begin{proof}
$(\Rightarrow)$ Definimos $g(x)=V_a^x(f)$ y $h(x)=V_a^x(f)-f(x)$. Se verifica que $g$ es creciente y $h$ también.

$(\Leftarrow)$ Si $f=g-h$, entonces $V_a^b(f)\leq V_a^b(g)+V_a^b(h)=|g(b)-g(a)|+|h(b)-h(a)|<\infty$.
\end{proof}

\section{Funciones absolutamente continuas}

\begin{definition}
\label{def:absolutamente-continua}
Una función $f\colon[a,b]\to\mathbb{R}$ es \emph{absolutamente continua} si para todo $\varepsilon>0$ existe $\delta>0$ tal que, para toda familia finita de intervalos disjuntos $(a_i,b_i)$ con $\sum(b_i-a_i)<\delta$, se tiene $\sum|f(b_i)-f(a_i)|<\varepsilon$.
\end{definition}

\begin{proposition}
\label{prop:ac-implica-va}
Toda función absolutamente continua es de variación acotada.
\end{proposition}

\section{Derivada puntual y conjuntos de medida cero}

Antes de enunciar los teoremas de derivación de Lebesgue, necesitamos dos nociones que se desarrollan con plenitud en el \cref{ch:teoria-medida}, pero cuyo significado básico puede precisarse aquí mismo.

\begin{definition}
\label{def:derivada-puntual}
Sea $f\colon[a,b]\to\mathbb{R}$. Se dice que $f$ es \emph{derivable} en $x_0\in(a,b)$ si existe
\[
  f'(x_0)=\lim_{h\to 0}\frac{f(x_0+h)-f(x_0)}{h}.
\]
\end{definition}

\begin{definition}
\label{def:medida-cero-elemental}
Un conjunto $N\subseteq\mathbb{R}$ tiene \emph{medida cero} (en el sentido de Lebesgue) si para todo $\varepsilon>0$ existe una familia numerable de intervalos abiertos $(I_n)$ tal que $N\subseteq\bigcup_n I_n$ y $\sum_n \ell(I_n)<\varepsilon$, donde $\ell(I_n)$ es la longitud del intervalo.
\end{definition}

\begin{proposition}
\label{prop:medida-cero-numerable}
Todo conjunto numerable tiene medida cero. La unión numerable de conjuntos de medida cero tiene medida cero.
\end{proposition}

Diremos que una propiedad se cumple \emph{casi en todas partes} (c.e.p.) si se cumple en todo punto salvo quizás en un conjunto de medida cero.

\section{Derivadas casi en todas partes}

\begin{theorem}[Lebesgue]
\label{teo:lebesgue-derivada}
Toda función de variación acotada en $[a,b]$ es derivable casi en todas partes y su derivada es integrable (en el sentido de Lebesgue; véase el \cref{ch:teoria-medida}).
\end{theorem}

Este teorema, profundamente vinculado a la teoría de la medida, muestra que la regularidad de una función no es una propiedad puntual sino ``casi'' global.

\section{Relación entre continuidad absoluta y variación acotada}

\begin{theorem}
\label{teo:carac-ac}
Una función $f$ es absolutamente continua en $[a,b]$ si y solo si existe una función integrable $g$ tal que
\[
  f(x)=f(a)+\int_a^x g(t)\,dt.
\]
En tal caso, $g=f'$ casi en todas partes.
\end{theorem}

\section{Teorema fundamental del cálculo para funciones absolutamente continuas}

\begin{theorem}
\label{teo:fundamental-ac}
Si $f$ es absolutamente continua en $[a,b]$, entonces
\[
  f(x)-f(a)=\int_a^x f'(t)\,dt\quad\text{para todo }x\in[a,b].
\]
Recíprocamente, si $g$ es integrable, entonces $f(x)=\int_a^x g(t)\,dt$ es absolutamente continua y $f'=g$ casi en todas partes.
\end{theorem}

\begin{exercise}
Demuestre que toda función absolutamente continua es uniformemente continua, pero que el recíproco no es cierto (sugerencia: función de Cantor).
\end{exercise}

\begin{exercise}
Calcule la variación total de $f(x)=x^2$ en $[0,1]$.
\end{exercise}

\begin{exercise}
Demuestre que la composición de una función absolutamente continua con una lipschitziana es absolutamente continua.
\end{exercise}

\begin{problem}
Sea $f$ de variación acotada. Demuestre que $f$ es continua si y solo si $x\mapsto V_a^x(f)$ es continua.
\end{problem}

\chapter{Aproximación de funciones}
\label{ch:aproximacion}
\index{Aproximación uniforme}\index{Polinomios de Bernstein}\index{Teorema de Stone--Weierstrass}\index{Función suave}

\section{Aproximación uniforme}

Una de las preguntas centrales del análisis matemático es hasta qué punto una función complicada puede sustituirse por una más sencilla sin cometer un error mayor que una tolerancia prescrita. Karl Weierstrass, en 1885, respondió esta pregunta de manera definitiva para funciones continuas: todo función continua en un intervalo compacto puede aproximarse uniformemente por polinomios. Este resultado, que parecía contradecir la existencia de funciones continuas no derivables, mostró que la ``irregularidad'' no es un obstáculo para la aproximación polinomial, sino solo para la representación en serie de Taylor.

\begin{theorem}[Weierstrass]
\label{teo:weierstrass-aprox}
Para toda $f\in C([a,b])$ y todo $\varepsilon>0$, existe un polinomio $p$ tal que $\sup_{x\in[a,b]}|f(x)-p(x)|<\varepsilon$.
\end{theorem}

\section{Polinomios de Bernstein}

Sergei Bernstein, en 1912, proporcionó una prueba constructiva del teorema de Weierstrass mediante polinomios que hoy llevan su nombre. A diferencia de la interpolación de Lagrange, los polinomios de Bernstein aproximan globalmente y preservan propiedades cualitativas como la monotonía y la convexidad.

\begin{definition}
\label{def:bernstein}
Sea $f\in C([0,1])$. El \emph{polinomio de Bernstein} de grado $n$ asociado a $f$ es
\[
  B_n(f,x)=\sum_{k=0}^{n}f\Bigl(\frac{k}{n}\Bigr)\binom{n}{k}x^k(1-x)^{n-k}.
\]
\end{definition}

\begin{theorem}[Bernstein]
\label{teo:bernstein}
La sucesión $B_n(f,\cdot)$ converge uniformemente a $f$ en $[0,1]$.
\end{theorem}
\begin{proof}[Bosquejo analítico]
Sean $r_k(x)=\binom{n}{k}x^k(1-x)^{n-k}$ los coeficientes binomiales. Se verifica directamente que $\sum_{k=0}^n r_k(x)=1$, $\sum_{k=0}^n (k/n)r_k(x)=x$ y $\sum_{k=0}^n (k/n-x)^2 r_k(x)=x(1-x)/n\leq 1/(4n)$. Dado $\varepsilon>0$, sea $\delta>0$ tal que $|f(x)-f(y)|<\varepsilon$ si $|x-y|<\delta$ (continuidad uniforme). Partiendo la suma en $|k/n-x|<\delta$ y $|k/n-x|\geq\delta$ y usando la estimación de varianza $\sum (k/n-x)^2 r_k(x)=x(1-x)/n$, se obtiene
\[
  |B_n(f,x)-f(x)|\leq\varepsilon+2\|f\|_\infty\,\frac{x(1-x)}{n\delta^2}\leq\varepsilon+\frac{\|f\|_\infty}{2n\delta^2},
\]
lo cual es menor que $2\varepsilon$ para $n$ suficientemente grande, independientemente de $x$.
\end{proof}

\section{Teorema de Stone--Weierstrass}

El teorema de Weierstrass fue generalizado espectacularmente por Marshall Stone en 1937. Stone observó que la propiedad crucial de los polinomios no era su forma algebraica, sino su capacidad de separar puntos.

\begin{theorem}[Stone--Weierstrass]
\label{teo:stone-weierstrass}
Sea $X$ un espacio compacto Hausdorff y $\mathcal{A}\subseteq C(X,\mathbb{R})$ una subálgebra que:
\begin{enumerate}
  \item[(i)] separa puntos: para $x\neq y$ existe $f\in\mathcal{A}$ con $f(x)\neq f(y)$;
  \item[(ii)] contiene las constantes.
\end{enumerate}
Entonces $\mathcal{A}$ es densa en $C(X)$ con la métrica del supremo.
\end{theorem}
\begin{proof}[Bosquejo]
La demostración procede en tres etapas: (1) toda subálgebra cerrada es un retículo (cerrada bajo máximos y mínimos); (2) un retículo que separa puntos y contiene constantes es denso (por el teorema de Kakutani--Krein); (3) la clausura de $\mathcal{A}$ satisface estas condiciones.
\end{proof}

\begin{corollary}
\label{cor:stone-complejo}
Si $\mathcal{A}\subseteq C(X,\mathbb{C})$ es una subálgebra cerrada bajo conjugación, separa puntos y contiene constantes, entonces es densa.
\end{corollary}

\section{Funciones suaves}

\begin{theorem}
\label{teo:aprox-cinfinito}
Las funciones $C^{\infty}$ con soporte compacto son densas en $L^p(\mathbb{R}^n)$ para $1\leq p<\infty$.
\end{theorem}

\section{Densidad en espacios de funciones}

\begin{theorem}
\label{teo:densidad-lp}
Para $1\leq p<\infty$, los polinomios trigonométricos son densos en $L^p([0,2\pi])$.
\end{theorem}

\section{Aplicaciones}

\begin{example}[Métodos numéricos]
\label{ej:numericos}
La aproximación polinomial por trozos (splines) y los elementos finitos se basan en variantes del teorema de Weierstrass adaptadas a restricciones locales.
\end{example}

\begin{example}[Teoría de representación]
\label{ej:representacion}
En el análisis armónico abstracto, el teorema de Stone--Weierstrass garantiza que las funciones continuas sobre un grupo compacto pueden aproximarse por representaciones irreducibles (teorema de Peter--Weyl).
\end{example}

\begin{exercise}
Calcule el polinomio de Bernstein de grado 2 para $f(x)=x^2$ en $[0,1]$.
\end{exercise}

\begin{exercise}
Demuestre que las funciones lipschitzianas son densas en $C([a,b])$.
\end{exercise}

\begin{exercise}
Sea $X=\{1,2\}$ con la topología discreta. Describa explícitamente $C(X)$ y verifique el teorema de Stone--Weierstrass para este caso.
\end{exercise}

\begin{problem}
Use el teorema de Stone--Weierstrass para probar que las funciones trigonométricas $1,\cos x,\sin x,\cos 2x,\sin 2x,\dots$ generan un conjunto denso en $C([0,2\pi])$.
\end{problem}

\chapter{Introducción a las distribuciones}
\label{ch:distribuciones}
\index{Distribución}\index{Función de prueba}\index{Delta de Dirac}\index{Derivada débil}\index{Soporte}

\section{Motivación}

La física del siglo~XX obligó a los matemáticos a manipular objetos que no eran funciones en el sentido clásico. La ``delta de Dirac'', usada extensamente desde los trabajos de Paul Dirac en mecánica cuántica, era una ``función'' $\delta(x)$ nula en todo $x\neq 0$, infinita en $x=0$, y cuya integral sobre $\mathbb{R}$ valía $1$. Tales propiedades son incompatibles con la teoría de la integral de Riemann. La respuesta matemática, desarrollada por Laurent Schwartz a partir de 1945, fue la teoría de \emph{distribuciones}: una generalización del concepto de función en la que los objetos se definen no por sus valores puntuales, sino por su acción sobre un espacio de ``funciones de prueba''.

\begin{definition}
\label{def:funcion-prueba}
Una función $\varphi\colon\mathbb{R}\to\mathbb{R}$ es de \emph{prueba} (o \emph{test}) si es de clase $C^{\infty}$ y tiene \emph{soporte compacto}, es decir, existe $M>0$ tal que $\varphi(x)=0$ para $|x|>M$. Al espacio de tales funciones se le denota $\mathcal{D}(\mathbb{R})$.
\end{definition}

El espacio $\mathcal{D}(\mathbb{R})$ es no vacío: las funciones de la forma $\varphi(x)=e^{-1/(1-x^2)}$ para $|x|<1$ extendidas por $0$ fuera de $[-1,1]$ son ejemplo paradigmático.

\section{Funciones de prueba}

\begin{definition}
\label{def:convergencia-prueba}
Una sucesión $(\varphi_n)$ converge a $\varphi$ en $\mathcal{D}(\mathbb{R})$ si existe un compacto $K$ tal que $\operatorname{sop}(\varphi_n)\subseteq K$ para todo $n$, y $\varphi_n^{(k)}\to\varphi^{(k)}$ uniformemente para todo $k\geq 0$.
\end{definition}

\section{Distribuciones}

\begin{definition}
\label{def:distribucion}
Una \emph{distribución} es un funcional lineal $T\colon\mathcal{D}(\mathbb{R})\to\mathbb{R}$ que es \emph{continuo}: si $\varphi_n\to\varphi$ en $\mathcal{D}(\mathbb{R})$, entonces $T(\varphi_n)\to T(\varphi)$. Al conjunto de distribuciones se le denota $\mathcal{D}'(\mathbb{R})$.
\end{definition}

\begin{example}[Funciones localmente integrables]
\label{ej:distribucion-l1}
Si $f\in L^1_{\mathrm{loc}}(\mathbb{R})$, entonces
\[
  T_f(\varphi)=\int_{-\infty}^{\infty}f(x)\varphi(x)\,dx
\]
define una distribución. Así, toda función integrable da origen a una distribución.
\end{example}

\begin{example}[Medidas de Radon]
\label{ej:distribucion-medida}
Toda medida de Borel finita sobre compactos $\mu$ define una distribución $T_\mu(\varphi)=\int\varphi\,d\mu$.
\end{example}

\section{Delta de Dirac}

\begin{definition}
\label{def:delta-dirac}
La \emph{delta de Dirac} centrada en $a\in\mathbb{R}$ es la distribución definida por
\[
  \delta_a(\varphi)=\varphi(a).
\]
\end{definition}

\begin{remark}
La delta de Dirac no es una función en el sentido clásico: no existe $f\in L^1_{\mathrm{loc}}$ tal que $\int f\varphi=\varphi(0)$ para todo $\varphi$. Sin embargo, puede aproximarse por funciones: si $(f_n)$ es una sucesión de funciones no negativas con integral $1$ y soporte tendiendo a $\{0\}$, entonces $T_{f_n}\to\delta_0$ en $\mathcal{D}'(\mathbb{R})$.
\end{remark}

\section{Derivadas débiles}

La mayor ventaja de las distribuciones es que todas son infinitamente diferenciables.

\begin{definition}
\label{def:derivada-distribucion}
La \emph{derivada} de una distribución $T$ es la distribución $T'$ definida por
\[
  T'(\varphi)=-T(\varphi').
\]
\end{definition}

\begin{example}
\label{ej:derivada-heaviside}
Sea $H$ la función de Heaviside ($H(x)=0$ para $x<0$, $H(x)=1$ para $x\geq 0$). Como distribución, $H'=\delta_0$:
\[
  H'(\varphi)=-\int_0^{\infty}\varphi'(x)\,dx=\varphi(0)=\delta_0(\varphi).
\]
\end{example}

\section{Soporte de una distribución}

\begin{definition}
\label{def:soporte-distribucion}
El \emph{soporte} de $T\in\mathcal{D}'(\mathbb{R})$ es el complemento del mayor abierto $U$ tal que $T(\varphi)=0$ para toda $\varphi$ con soporte en $U$.
\end{definition}

\begin{example}
El soporte de $\delta_a$ es $\{a\}$.
\end{example}

\section{Aplicaciones elementales}

\begin{example}[Ecuaciones diferenciales]
\label{ej:edo-distribuciones}
La ecuación $y'=\delta_0$ tiene como solución general $y=H+C$ en el sentido de distribuciones, donde $H$ es la función de Heaviside.
\end{example}

\begin{example}[Ecuaciones en derivadas parciales]
\label{ej:edp-distribuciones}
La ecuación de Laplace $\Delta u=0$ admite soluciones fundamentales en el sentido de distribuciones, como la función de Green.
\end{example}

\begin{exercise}
Demuestre que si $f\in C^1(\mathbb{R})$, entonces la derivada distribucional de $T_f$ coincide con $T_{f'}$.
\end{exercise}

\begin{exercise}
Calcule la segunda derivada distribucional de $f(x)=|x|$.
\end{exercise}

\begin{exercise}
Sea $T_n(\varphi)=n\int_{-1/n}^{1/n}\varphi(x)\,dx$. Demuestre que $T_n\to 2\delta_0$ en $\mathcal{D}'(\mathbb{R})$.
\end{exercise}

\begin{problem}
Demuestre que toda distribución con soporte compacto puede extenderse a un funcional lineal continuo sobre $C^{\infty}(\mathbb{R})$.
\end{problem}

\chapter{Introducción a la teoría de la medida}
\label{ch:teoria-medida}
\index{Medida}\index{sigma-algebra@$\sigma$-álgebra}\index{Medida exterior}\index{Conjunto medible}\index{Medida de Lebesgue}\index{Función medible}\index{Integral de Lebesgue}

\section{Motivación}

La integral de Riemann, desarrollada en el siglo~XIX, era suficiente para la mayoría de las aplicaciones del análisis clásico. Sin embargo, presentaba limitaciones fundamentales: no podía integrar todos los límites de sucesiones de funciones integrables, no manejaba bien conjuntos de medida cero, y su extensión a dimensiones superiores resultaba torpe. La necesidad de una teoría de integración más robusta impulsó a Émile Borel, Henri Lebesgue y, posteriormente, Constantin Carathéodory a desarrollar la teoría de la medida.

La idea central de Lebesgue (1902) fue revolucionaria: en lugar de particionar el dominio de la función (como hace Riemann), se particiona el codominio. El área bajo la curva se aproxima por sumas de ``barras'' horizontales, lo cual permite integrar funciones mucho más generales y garantiza teoremas de convergencia poderosos.

\section{Álgebras y \texorpdfstring{$\sigma$}{sigma}-álgebras}

\begin{definition}
\label{def:sigma-algebra}
Una \emph{$\sigma$-álgebra} sobre un conjunto $X$ es una familia $\mathcal{A}\subseteq\mathcal{P}(X)$ tal que:
\begin{enumerate}
  \item[(i)] $X\in\mathcal{A}$.
  \item[(ii)] Si $A\in\mathcal{A}$, entonces $X\setminus A\in\mathcal{A}$.
  \item[(iii)] Si $(A_n)_{n\in\mathbb{N}}\subseteq\mathcal{A}$, entonces $\bigcup_{n=1}^{\infty}A_n\in\mathcal{A}$.
\end{enumerate}
\end{definition}

\begin{definition}
\label{def:algebra}
Una \emph{álgebra} de conjuntos satisface (i), (ii) y la unión finita (en lugar de numerable).
\end{definition}

\begin{example}
\label{ej:borel}
La $\sigma$-álgebra de \emph{Borel} $\mathcal{B}(\mathbb{R})$ es la menor $\sigma$-álgebra que contiene a todos los intervalos abiertos. Es generada también por los intervalos cerrados, los semicerrados, o las bolas abiertas.
\end{example}

\section{Medidas}

\begin{definition}
\label{def:medida}
Una \emph{medida} sobre una $\sigma$-álgebra $\mathcal{A}$ es una función $\mu\colon\mathcal{A}\to[0,+\infty]$ tal que:
\begin{enumerate}
  \item[(i)] $\mu(\emptyset)=0$.
  \item[(ii)] $\sigma$-aditividad: si $(A_n)$ son disjuntos dos a dos, $\mu(\bigcup_n A_n)=\sum_n\mu(A_n)$.
\end{enumerate}
Al terceto $(X,\mathcal{A},\mu)$ se le llama \emph{espacio de medida}.
\end{definition}

\begin{example}[Medida de conteo]
\label{ej:conteo}
Sobre cualquier conjunto $X$, $\mu(A)=|A|$ (número de elementos, o $+\infty$ si $A$ es infinito) es una medida.
\end{example}

\begin{example}[Medida de Dirac]
\label{ej:dirac-medida}
Fijado $x_0\in X$, $\delta_{x_0}(A)=1$ si $x_0\in A$, $0$ en otro caso, es una medida.
\end{example}

\section{Medida exterior}

\begin{definition}
\label{def:medida-exterior}
Una \emph{medida exterior} sobre $X$ es una función $\mu^*\colon\mathcal{P}(X)\to[0,+\infty]$ tal que:
\begin{enumerate}
  \item[(i)] $\mu^*(\emptyset)=0$.
  \item[(ii)] Monotonía: $A\subseteq B\Rightarrow\mu^*(A)\leq\mu^*(B)$.
  \item[(iii)] Subaditividad numerable: $\mu^*(\bigcup_n A_n)\leq\sum_n\mu^*(A_n)$.
\end{enumerate}
\end{definition}

\section{Conjuntos medibles}

\begin{theorem}[Carathéodory]
\label{teo:caratheodory}
Dada una medida exterior $\mu^*$, la familia
\[
  \mathcal{M}=\{A\subseteq X:\,\mu^*(E)=\mu^*(E\cap A)+\mu^*(E\setminus A)\text{ para todo }E\subseteq X\}
\]
es una $\sigma$-álgebra, y la restricción de $\mu^*$ a $\mathcal{M}$ es una medida completa.
\end{theorem}

\section{Medida de Lebesgue}

\begin{definition}
\label{def:lebesgue-exterior}
La \emph{medida exterior de Lebesgue} $\lambda^*$ en $\mathbb{R}$ se define por
\[
  \lambda^*(A)=\inf\Bigl\{\sum_{n=1}^{\infty}(b_n-a_n):\,A\subseteq\bigcup_{n=1}^{\infty}(a_n,b_n)\Bigr\}.
\]
\end{definition}

\begin{theorem}
\label{teo:lebesgue-medida}
Los conjuntos medibles en el sentido de Carathéodory para $\lambda^*$ forman una $\sigma$-álgebra $\mathcal{L}(\mathbb{R})$ que contiene a $\mathcal{B}(\mathbb{R})$. La restricción $\lambda=\lambda^*|_{\mathcal{L}}$ es la \emph{medida de Lebesgue}, invariante por traslaciones y que asigna longitud $b-a$ a cada intervalo $[a,b]$.
\end{theorem}

\section{Funciones medibles}

\begin{definition}
\label{def:medible}
Una función $f\colon(X,\mathcal{A})\to\mathbb{R}$ es \emph{medible} si $f^{-1}(U)\in\mathcal{A}$ para todo abierto $U\subseteq\mathbb{R}$.
\end{definition}

\begin{proposition}
\label{prop:limite-medible}
El límite puntual de una sucesión de funciones medibles es medible.
\end{proposition}

\section{Integral de Lebesgue (introducción)}

\begin{definition}
\label{def:integral-lebesgue}
Para una función simple no negativa $s=\sum_{i=1}^{n}a_i\mathbf{1}_{A_i}$, se define $\int s\,d\mu=\sum a_i\mu(A_i)$. Para $f\geq 0$ medible,
\[
  \int f\,d\mu=\sup\Bigl\{\int s\,d\mu:\,0\leq s\leq f,\,s\text{ simple}\Bigr\}.
\]
\end{definition}

\section{Principales teoremas de convergencia}

\begin{theorem}[de convergencia monótona]
\label{teo:monotona}
Si $0\leq f_n\uparrow f$ puntualmente, entonces $\int f_n\,d\mu\uparrow\int f\,d\mu$.
\end{theorem}

\begin{theorem}[de Fatou]
\label{teo:fatou}
Si $f_n\geq 0$, entonces $\int\liminf f_n\,d\mu\leq\liminf\int f_n\,d\mu$.
\end{theorem}

\begin{theorem}[de convergencia dominada]
\label{teo:dominada}
Si $f_n\to f$ puntualmente y existe $g$ integrable tal que $|f_n|\leq g$, entonces $\int f_n\,d\mu\to\int f\,d\mu$.
\end{theorem}

\begin{exercise}
Demuestre que la medida de Lebesgue de $\mathbb{Q}$ es cero.
\end{exercise}

\begin{exercise}
Sea $f_n(x)=n\mathbf{1}_{[0,1/n]}(x)$. Calcule $\lim_{n\to\infty}\int f_n\,d\lambda$ y compare con $\int\lim f_n\,d\lambda$. ¿Cuál teorema de convergencia falla aquí?
\end{exercise}

\begin{exercise}
Pruebe que toda función continua en $[a,b]$ es medible de Borel.
\end{exercise}

\begin{problem}
Demuestre que si $f\colon\mathbb{R}\to\mathbb{R}$ es derivable, entonces $f'$ es medible de Lebesgue (puede no ser continua).
\end{problem}

\chapter{Introducción al análisis funcional}
\label{ch:analisis-funcional}
\index{Espacio normado}\index{Espacio de Banach}\index{Espacio de Hilbert}\index{Operador lineal}\index{Funcional lineal}\index{Operador compacto}

\section{Espacios normados}

El análisis funcional nació de la necesidad de estudiar espacios cuyos elementos son funciones, y donde la noción de ``tamaño'' se generaliza mediante una norma. Stefan Banach, en su tesis de 1920, sistematizó el estudio de estos espacios completos, que hoy llevan su nombre. El salto conceptual es enorme: en lugar de estudiar una función aislada, se estudia la \emph{estructura} del espacio que la contiene, y las transformaciones lineales entre tales espacios.

\begin{definition}
\label{def:norma}
Una \emph{norma} sobre un espacio vectorial real $E$ es una función $\|\cdot\|\colon E\to\mathbb{R}$ tal que:
\begin{enumerate}
  \item[(i)] $\|x\|\geq 0$, y $\|x\|=0\Leftrightarrow x=0$.
  \item[(ii)] $\|\lambda x\|=|\lambda|\,\|x\|$ para todo $\lambda\in\mathbb{R}$.
  \item[(iii)] $\|x+y\|\leq\|x\|+\|y\|$ (desigualdad triangular).
\end{enumerate}
Al par $(E,\|\cdot\|)$ se le llama \emph{espacio normado}.
\end{definition}

\begin{example}
\label{ej:lp-normado}
Para $1\leq p<\infty$, el espacio $\ell^p$ de sucesiones $(x_n)$ con $\sum|x_n|^p<\infty$ es normado con $\|x\|_p=(\sum|x_n|^p)^{1/p}$.
\end{example}

\begin{example}
\label{ej:cb-normado}
$C([a,b])$ con $\|f\|_\infty=\sup_{t\in[a,b]}|f(t)|$ es un espacio normado.
\end{example}

\section{Espacios de Banach}

\begin{definition}
\label{def:banach}
Un espacio normado completo (en la métrica inducida por la norma) se llama \emph{espacio de Banach}.
\end{definition}

\begin{example}
\label{ej:banach-clasicos}
$\mathbb{R}^n$, $\ell^p$ para $1\leq p\leq\infty$, $C([a,b])$ con el supremo, $L^p([a,b])$ para $1\leq p\leq\infty$, son espacios de Banach.
\end{example}

\begin{theorem}
\label{teo:completitud-lp}
$L^p(X,\mu)$ es completo para $1\leq p\leq\infty$ (teorema de Riesz--Fischer).
\end{theorem}

\section{Espacios de Hilbert}

\begin{definition}
\label{def:producto-interno}
Un \emph{producto interno} sobre un espacio vectorial real $H$ es una forma bilineal simétrica $\langle\cdot,\cdot\rangle\colon H\times H\to\mathbb{R}$ tal que:
\begin{enumerate}
  \item[(i)] $\langle x,x\rangle\geq 0$, con igualdad solo si $x=0$.
  \item[(ii)] $\langle x,y\rangle=\langle y,x\rangle$.
  \item[(iii)] Es lineal en cada argumento.
\end{enumerate}
\end{definition}

\begin{definition}
\label{def:hilbert}
Un espacio con producto interno completo respecto a la norma $\|x\|=\sqrt{\langle x,x\rangle}$ se llama \emph{espacio de Hilbert}.
\end{definition}

\begin{example}
\label{ej:l2-hilbert}
$\ell^2$ y $L^2([a,b])$ son espacios de Hilbert con productos internos naturales.
\end{example}

\begin{theorem}[de representación de Riesz]
\label{teo:riesz}
Si $H$ es un espacio de Hilbert, todo funcional lineal continuo $\phi\colon H\to\mathbb{R}$ puede representarse como $\phi(x)=\langle x,y\rangle$ para un único $y\in H$.
\end{theorem}

\section{Operadores lineales}

\begin{definition}
\label{def:operador}
Un \emph{operador lineal} entre espacios normados $E$ y $F$ es una aplicación lineal $T\colon E\to F$. Es \emph{acotado} si existe $M>0$ tal que $\|Tx\|\leq M\|x\|$ para todo $x\in E$.
\end{definition}

\begin{theorem}
\label{teo:continuo-acotado}
Un operador lineal entre espacios normados es continuo si y solo si es acotado.
\end{theorem}

\section{Funcionales lineales}

\begin{definition}
\label{def:functional}
Un \emph{funcional lineal} es un operador lineal $E\to\mathbb{R}$. El espacio de funcionales lineales continuos se denomina \emph{dual topológico} $E'$.
\end{definition}

\begin{theorem}[Hahn--Banach, enunciado]
\label{teo:hahn-banach}
Todo funcional lineal definido en un subespacio y acotado por una seminorma puede extenderse a todo el espacio preservando la cota.
\end{theorem}

\section{Operadores compactos}

\begin{definition}
\label{def:operador-compacto}
Un operador $T\colon E\to F$ es \emph{compacto} si transforma conjuntos acotados en conjuntos relativamente compactos.
\end{definition}

\begin{theorem}
\label{teo:compacidad-rango}
Los operadores compactos de rango finito son densos en el espacio de operadores compactos (en espacios de Hilbert).
\end{theorem}

\section{Panorama de los teoremas fundamentales}

El análisis funcional posee tres grandes teoremas de estructura que aparecen una y otra vez:
\begin{enumerate}
  \item \textbf{Hahn--Banach}: existencia de suficientes funcionales para separar puntos.
  \item \textbf{Principio de acotación uniforme (Banach--Steinhaus)}: una familia de operadores acotada puntualmente es uniformemente acotada.
  \item \textbf{Teorema de la función abierta (Banach--Schauder)}: un operador lineal continuo y sobreyectivo entre Banach es abierto.
\end{enumerate}

\section{Aplicaciones}

\begin{example}[Ecuaciones diferenciales]
\label{ej:edos-af}
Los problemas de contorno lineales pueden formularse como $Lu=f$, donde $L$ es un operador entre espacios de Sobolev. La inversibilidad de $L$ equivale a la existencia de una solución única.
\end{example}

\begin{example}[Mecánica cuántica]
\label{ej:cuantica}
Los estados de un sistema cuántico son vectores de un espacio de Hilbert; los observables son operadores autoadjuntos; la evolución temporal está dada por un operador unitario.
\end{example}

\begin{exercise}
Demuestre que en un espacio normado, la bola unidad es compacta si y solo si el espacio tiene dimensión finita (teorema de Riesz).
\end{exercise}

\begin{exercise}
Pruebe la desigualdad de Cauchy--Schwarz en un espacio con producto interno.
\end{exercise}

\begin{exercise}
Demuestre que todo espacio de Hilbert es reflexivo.
\end{exercise}

\begin{problem}
Sea $T\colon\ell^2\to\ell^2$ definido por $T(x_1,x_2,\dots)=(x_1/1,x_2/2,\dots)$. Demuestre que $T$ es compacto y calcule su norma.
\end{problem}

\chapter{Diferenciación en espacios normados}
\label{ch:diferenciacion}
\index{Diferencial de Fréchet}\index{Derivada de Fréchet}\index{Teorema de la función inversa}\index{Regla de la cadena}

\section{El diferencial de Fréchet}

La noción de derivada se extiende de manera natural del cálculo en una variable
al contexto de espacios normados. La idea central sigue siendo la aproximación
lineal local de una función.

\begin{definition}\label{def:frechet}
Sean $E$ y $F$ espacios normados, $U\subseteq E$ abierto y $f\colon U\to F$.
Se dice que $f$ es \emph{diferenciable} (en el sentido de Fréchet) en
$x_0\in U$ si existe una aplicación lineal continua $T\colon E\to F$ tal que
\begin{equation}\label{eq:frechet}
  \lim_{h\to 0}\frac{\|f(x_0+h)-f(x_0)-T(h)\|}{\|h\|}=0.
\end{equation}
En tal caso, $T$ es única y se denomina \emph{diferencial} de $f$ en $x_0$,
notada $Df(x_0)$ o $f'(x_0)$.
\end{definition}

\begin{remark}
Cuando $E=F=\mathbb{R}$, la definición anterior coincide con la derivada usual:
$T(h)=f'(x_0)\,h$.
\end{remark}

\begin{example}\label{ej:derivada-lineal}
Si $f\colon E\to F$ es lineal y continua, entonces $f$ es diferenciable en
todo punto y $Df(x)=f$ para todo $x\in E$.
\end{example}

\begin{example}\label{ej:derivada-bilineal}
Sea $b\colon E_1\times E_2\to F$ una aplicación bilineal continua. Entonces $b$
es diferenciable en todo punto y su diferencial está dada por
\[
  Db(x_1,x_2)(h_1,h_2)=b(x_1,h_2)+b(h_1,x_2).
\]
\end{example}

\section{Reglas de diferenciación}

\begin{proposition}[Regla de la cadena]\label{prop:cadena}
Sean $U\subseteq E$ abierto, $V\subseteq F$ abierto, $f\colon U\to V$
diferenciable en $x_0$ y $g\colon V\to G$ diferenciable en $f(x_0)$. Entonces
$g\circ f$ es diferenciable en $x_0$ y
\begin{equation}\label{eq:cadena}
  D(g\circ f)(x_0)=Dg(f(x_0))\circ Df(x_0).
\end{equation}
\end{proposition}

\begin{theorem}[del valor medio]\label{teo:valor-medio}
Sea $f\colon U\to F$ diferenciable en el segmento $[x,y]\subseteq U$. Entonces
\begin{equation}\label{eq:valor-medio}
  \|f(y)-f(x)\|\leq\sup_{z\in[x,y]}\|Df(z)\|\,\|y-x\|.
\end{equation}
\end{theorem}

\section{El teorema de la función inversa}

\begin{theorem}[de la función inversa]\label{teo:inversa}
Sea $f\colon U\subseteq E\to F$ de clase $C^1$ (es decir, continuamente
diferenciable) y sea $x_0\in U$ tal que $Df(x_0)$ es un isomorfismo lineal
sobreyectivo. Entonces existen entornos abiertos $V\subseteq U$ de $x_0$ y
$W\subseteq F$ de $f(x_0)$ tales que $f\colon V\to W$ es biyectiva y su
inversa $f^{-1}\colon W\to V$ es de clase $C^1$. Además,
\begin{equation}\label{eq:inversa}
  D(f^{-1})(f(x_0))=\bigl(Df(x_0)\bigr)^{-1}.
\end{equation}
\end{theorem}

Este teorema constituye uno de los pilares del cálculo diferencial en varias
variables y en dimensiones infinitas, pues garantiza que la linealización local
de una función controla su comportamiento local no lineal siempre que la
linealización sea invertible.

\section{Ejercicios}

\begin{exercise}
Calcule el diferencial de Fréchet de la función $f\colon\mathbb{R}^2\to\mathbb{R}$
definida por $f(x,y)=x^2+xy+y^2$ en un punto genérico $(x_0,y_0)$.
\end{exercise}

\begin{exercise}
Demuestre que si $f$ es diferenciable en $x_0$, entonces $f$ es continua en
$x_0$.
\end{exercise}

\begin{exercise}
Sea $f\colon M_n(\mathbb{R})\to M_n(\mathbb{R})$ dada por $f(A)=A^2$. Calcule
$Df(A)(H)$ para matrices $A,H\in M_n(\mathbb{R})$.
\end{exercise}

\begin{exercise}
Use el \cref{teo:inversa} para probar que existe un entorno de la matriz
identidad $I$ en $M_n(\mathbb{R})$ en el cual toda matriz posee raíz cuadrada.
\end{exercise}

\begin{problem}
Formule y demuestre un teorema de la función implícita en espacios de Banach
usando el \cref{teo:inversa} como herramienta principal.
\end{problem}

\chapter{Introducción a la teoría espectral}
\label{ch:teoria-espectral}
\index{Valor propio}\index{Función propia}\index{Operador compacto}\index{Operador autoadjunto}\index{Espectro}\index{Descomposición espectral}

\section{Valores propios}

La teoría espectral busca generalizar, a operadores en dimensiones infinitas, la descomposición diagonal de matrices. David Hilbert, en sus trabajos sobre ecuaciones integrales (1904--1910), descubrió que muchos operadores importantes admiten una ``descomposición espectral'' análoga a la diagonalización, aunque los ``eigenvalores'' podían formar un espectro continuo. Esta teoría constituye uno de los puentes más profundos entre el análisis funcional y las aplicaciones físicas.

\begin{definition}
\label{def:espectro}
Sea $E$ un espacio de Banach y $T\colon E\to E$ un operador lineal acotado. El \emph{espectro} $\sigma(T)$ es el conjunto de $\lambda\in\mathbb{C}$ tales que $T-\lambda I$ no es invertible con inversa acotada. Los $\lambda$ para los cuales $T-\lambda I$ no es inyectivo son los \emph{valores propios} (o \emph{eigenvalores}).
\end{definition}

\begin{example}
\label{ej:multiplicacion}
Sea $T\colon L^2([0,1])\to L^2([0,1])$ definido por $Tf(x)=xf(x)$. Entonces $T$ no tiene valores propios, pero su espectro es $[0,1]$ (espectro continuo).
\end{example}

\section{Funciones propias}

\begin{definition}
\label{def:funcion-propia}
Si $\lambda$ es un valor propio de $T$, un vector no nulo $v$ tal que $Tv=\lambda v$ se denomina \emph{función propia} (o \emph{eigenvector}).
\end{definition}

\begin{example}
\label{ej:derivada}
El operador derivada $D\colon C^1([0,2\pi])\to C([0,2\pi])$ tiene valores propios $\lambda_n=in$ con funciones propias $e^{inx}$.
\end{example}

\section{Operadores compactos}

\begin{theorem}[Riesz--Schauder]
\label{teo:riesz-schauder}
Si $T$ es un operador compacto en un espacio de Banach, su espectro es a lo sumo numerable, $0$ es el único posible punto de acumulación, y todo $\lambda\neq 0$ en el espectro es un valor propio de multiplicidad finita.
\end{theorem}

\section{Operadores autoadjuntos (idea)}

\begin{definition}
\label{def:autoadjunto}
En un espacio de Hilbert, un operador $T$ es \emph{autoadjunto} si $\langle Tx,y\rangle=\langle x,Ty\rangle$ para todo $x,y$.
\end{definition}

\begin{theorem}[Espectral para autoadjuntos compactos]
\label{teo:espectral-compacto}
Si $T$ es autoadjunto y compacto, existe una base ortonormal de $H$ formada por funciones propias de $T$. Los valores propios son reales y tienden a $0$.
\end{theorem}

\section{Descomposición espectral (motivación)}

El teorema espectral general afirma que todo operador autoadjunto (no necesariamente compacto) admite una ``medida espectral'' $E$ tal que
\[
  T=\int_{\sigma(T)}\lambda\,dE(\lambda).
\]
Esta fórmula generaliza la diagonalización $A=PDP^{-1}$ de matrices simétricas.

\section{Aplicaciones a ecuaciones diferenciales}

\begin{example}[Problema de Sturm--Liouville]
\label{ej:sturm-liouville}
La ecuación $-(py')'+qy=\lambda w y$ en $[a,b]$ con condiciones de contorno define un operador autoadjunto cuyos valores propios forman una sucesión creciente $\lambda_n\to\infty$, y las funciones propias constituyen una base ortogonal de $L^2_w([a,b])$.
\end{example}

\begin{exercise}
Demuestre que los valores propios de un operador autoadjunto son reales.
\end{exercise}

\begin{exercise}
Sea $T\colon\ell^2\to\ell^2$ el operador shift derecha $T(x_1,x_2,\dots)=(0,x_1,x_2,\dots)$. Determine su espectro.
\end{exercise}

\begin{exercise}
Demuestre que si $T$ es compacto y autoadjunto, entonces $\|T\|$ es igual al máximo de los valores absolutos de sus valores propios.
\end{exercise}

\begin{problem}
Formule y demuestre un teorema de alternativa de Fredholm para operadores compactos: la ecuación $u-Tu=f$ tiene solución si y solo si $f$ es ortogonal al núcleo de $I-T^*$.
\end{problem}

\chapter{Introducción a las ecuaciones diferenciales}
\label{ch:ecuaciones-diferenciales}
\index{Ecuación diferencial}\index{Problema de valor inicial}\index{Existencia y unicidad}\index{Exponencial de matrices}\index{Estabilidad}\index{Ecuación en derivadas parciales}

\section{Problemas de valor inicial}

Las ecuaciones diferenciales fueron el motor original del cálculo: Newton las usó para describir la gravitación, y desde entonces han sido el lenguaje natural de la física y la ingeniería. Sin embargo, la pregunta fundamental ---¿tiene una ecuación dada una solución, y es única?--- no pudo responderse rigurosamente hasta que el análisis matemático desarrolló las herramientas de espacios métricos completos y contracciones.

\begin{definition}
\label{def:pvi}
Un \emph{problema de valor inicial} (PVI) de primer orden es una ecuación de la forma
\[
  y'(t)=F(t,y(t)),\qquad y(t_0)=y_0,
\]
donde $F\colon D\subseteq\mathbb{R}^2\to\mathbb{R}$ es una función dada.
\end{definition}

\section{Existencia y unicidad}

\begin{theorem}[Picard--Lindelöf]
\label{teo:picard}
Sea $F\colon[t_0-a,t_0+a]\times[y_0-b,y_0+b]\to\mathbb{R}$ continua y lipschitziana en la segunda variable. Entonces existe $\delta>0$ tal que el PVI tiene una única solución en $[t_0-\delta,t_0+\delta]$.
\end{theorem}
\begin{proof}[Bosquejo]
Se transforma el PVI en la ecuación integral
\[
  y(t)=y_0+\int_{t_0}^{t}F(s,y(s))\,ds.
\]
El operador de Picard $T(y)(t)=y_0+\int_{t_0}^t F(s,y(s))\,ds$ actúa sobre $C([t_0-\delta,t_0+\delta])$ y es contractivo para $\delta$ suficientemente pequeño. El teorema de Banach (\cref{teo:banach}) da la existencia y unicidad.
\end{proof}

\section{Métodos de punto fijo}

El método de punto fijo no es solo una prueba teórica; es también una herramienta numérica. Las iteraciones de Picard $y_{n+1}=T(y_n)$ convergen a la solución, y la estimación de error \eqref{eq:error-banach} permite controlar la precisión.

\begin{example}
\label{ej:picard-lineal}
Para $y'=ky$, $y(0)=1$, las iteradas de Picard son $y_n(t)=\sum_{j=0}^{n}(kt)^j/j!$, que convergen a $e^{kt}$.
\end{example}

\section{Sistemas lineales y exponencial de matrices}

\begin{definition}
\label{def:sistema-lineal}
Un \emph{sistema lineal} de ecuaciones diferenciales ordinarias es una ecuación vectorial
\[
  y'(t)=A\,y(t)+g(t),\qquad y(t_0)=y_0,
\]
donde $A\in M_n(\mathbb{R})$ es una matriz constante, $g$ es una función vectorial conocida y $y$ es la función incógnita con valores en $\mathbb{R}^n$.
\end{definition}

\begin{definition}
\label{def:exponencial-matriz-edo}
La \emph{exponencial de matriz} $e^{tA}$ se define como la serie
\[
  e^{tA}=\sum_{k=0}^{\infty}\frac{t^k A^k}{k!},
\]
que converge absolutamente para todo $t\in\mathbb{R}$ (véase el \cref{ap:algebra-lineal}).
\end{definition}

\begin{theorem}
\label{teo:solucion-lineal}
La solución del problema $y'=Ay$, $y(t_0)=y_0$ viene dada por $y(t)=e^{(t-t_0)A}y_0$.
\end{theorem}
\begin{proof}
Derivando término a término la serie de $e^{(t-t_0)A}$ se obtiene $\frac{d}{dt}e^{(t-t_0)A}=A\,e^{(t-t_0)A}$, luego $y'(t)=A\,y(t)$. Además $y(t_0)=e^{0}y_0=I y_0=y_0$.
\end{proof}

\begin{theorem}[Fórmula de variación de parámetros]
\label{teo:variacion-parametros}
La solución de $y'=Ay+g(t)$, $y(t_0)=y_0$ es
\[
  y(t)=e^{(t-t_0)A}y_0+\int_{t_0}^{t}e^{(t-s)A}g(s)\,ds.
\]
\end{theorem}

\section{Estabilidad}

\begin{definition}
\label{def:estabilidad}
Una solución $y(t)$ es \emph{estable} (en el sentido de Lyapunov) si para todo $\varepsilon>0$ existe $\delta>0$ tal que toda solución $\widetilde{y}(t)$ con $|\widetilde{y}(t_0)-y(t_0)|<\delta$ satisface $|\widetilde{y}(t)-y(t)|<\varepsilon$ para todo $t\geq t_0$.
\end{definition}

Recordemos que los \emph{valores propios} de una matriz $A$ son las raíces de $\det(A-\lambda I)=0$; los conceptos básicos de álgebra lineal necesarios están en el \cref{ap:algebra-lineal}.

\begin{theorem}
\label{teo:estabilidad-lineal}
Para sistemas lineales $y'=Ay$, la estabilidad de la solución nula está determinada por los valores propios de $A$: si todos tienen parte real negativa, la solución nula es asintóticamente estable. Si algún valor propio tiene parte real positiva, la solución nula es inestable.
\end{theorem}
\begin{proof}[Bosquejo]
Mediante la forma canónica de Jordan (o la descomposición espectral para matrices diagonalizables), se reduce el sistema a bloques donde cada bloque evoluciona como $e^{\lambda t}$. El comportamiento asintótico está gobernado por $\operatorname{Re}(\lambda)$.
\end{proof}

\section{Dependencia continua}

\begin{theorem}
\label{teo:dependencia-continua}
Bajo las hipótesis del teorema de Picard--Lindelöf, la solución depende continuamente de las condiciones iniciales $(t_0,y_0)$.
\end{theorem}

\section{Introducción a las ecuaciones en derivadas parciales}

Las ecuaciones en derivadas parciales (EDP) surgen cuando el estado de un sistema depende de múltiples variables independientes. Ejemplos clásicos son:
\begin{itemize}
  \item \textbf{Ecuación del calor} (parabólica): $u_t=\Delta u$.
  \item \textbf{Ecuación de ondas} (hiperbólica): $u_{tt}=\Delta u$.
  \item \textbf{Ecuación de Laplace} (elíptica): $\Delta u=0$.
\end{itemize}

La teoría de EDP requiere herramientas del análisis funcional (espacios de Sobolev), la teoría de la medida y la topología. Es uno de los campos más activos de las matemáticas contemporáneas.

\begin{exercise}
Aplique el método de Picard para $y'=y$, $y(0)=1$ y calcule las tres primeras iteradas.
\end{exercise}

\begin{exercise}
Demuestre que $y'=3y^{2/3}$, $y(0)=0$, tiene al menos dos soluciones. ¿Por qué falla Picard--Lindelöf?
\end{exercise}

\begin{exercise}
Analice la estabilidad del sistema $y'=Ay$ para $A=\begin{pmatrix}0&1\\-1&0\end{pmatrix}$.
\end{exercise}

\begin{problem}
Demuestre el teorema de Peano de existencia (sin unicidad) usando el teorema de Ascoli--Arzelà.
\end{problem}

\chapter{Introducción a la topología general}
\label{ch:topologia-general}
\index{Espacio topológico}\index{Subbase}\index{Axiomas de separación}\index{Producto de espacios}\index{Teorema de Tychonoff}

\section{Espacios topológicos}

La topología general nace de la abstracción de las propiedades de los espacios métricos que no dependen de la distancia explícita. Hausdorff, en su \emph{Grundzüge der Mengenlehre} (1914), formalizó esta idea definiendo un espacio topológico mediante una familia de conjuntos abiertos. Este salto conceptual permitió estudiar espacios que no provienen de ninguna métrica: topologías débiles en análisis funcional, topologías de Zariski en geometría algebraica, y topologías de orden en teoría de conjuntos.

\begin{definition}
\label{def:espacio-topologico}
Un \emph{espacio topológico} es un par $(X,\tau)$ donde $X$ es un conjunto y $\tau\subseteq\mathcal{P}(X)$ es una familia, llamada \emph{topología}, que satisface:
\begin{enumerate}
  \item[(T1)] $\emptyset,X\in\tau$.
  \item[(T2)] La unión arbitraria de elementos de $\tau$ pertenece a $\tau$.
  \item[(T3)] La intersección finita de elementos de $\tau$ pertenece a $\tau$.
\end{enumerate}
\end{definition}

\begin{example}
\label{ej:metrico-topologico}
Todo espacio métrico induce una topología métrica (\cref{teo:topologia-metrica}).
\end{example}

\begin{example}
\label{ej:cofinita}
La \emph{topología cofinita} sobre un conjunto infinito $X$, donde los abiertos son $\emptyset$ y los complementos de conjuntos finitos, no es metrizable.
\end{example}

\section{Bases y subbases}

\begin{definition}
\label{def:base-topologia}
Una \emph{base} de una topología $\tau$ es una subfamilia $\mathcal{B}\subseteq\tau$ tal que todo abierto no vacío es unión de elementos de $\mathcal{B}$.
\end{definition}

\begin{definition}
\label{def:subbase}
Una \emph{subbase} es una familia $\mathcal{S}$ tal que las intersecciones finitas de sus elementos forman una base.
\end{definition}

\section{Continuidad}

\begin{definition}
\label{def:continuidad-topologica}
Una función $f\colon(X,\tau_X)\to(Y,\tau_Y)$ es \emph{continua} si $f^{-1}(V)\in\tau_X$ para todo $V\in\tau_Y$.
\end{definition}

\begin{proposition}
\label{prop:comp-topologica}
La composición de funciones continuas es continua.
\end{proposition}

\section{Compacidad}

\begin{definition}
\label{def:compacto-topologia}
Un espacio topológico $X$ es \emph{compacto} si de todo recubrimiento abierto se puede extraer un subrecubrimiento finito.
\end{definition}

\begin{theorem}
\label{teo:compacto-cerrado}
En un espacio compacto Hausdorff, los compactos coinciden con los cerrados.
\end{theorem}

\section{Conexidad}

\begin{definition}
\label{def:conexo-topologia}
Un espacio topológico es \emph{conexo} si no puede escribirse como unión de dos abiertos disjuntos no vacíos.
\end{definition}

\begin{proposition}
\label{prop:imagen-conexa}
La imagen continua de un conexo es conexa.
\end{proposition}

\section{Axiomas de separación}

\begin{definition}
\label{def:separacion}
\begin{itemize}
  \item $T_0$: para puntos distintos, al menos uno tiene un entorno que no contiene al otro.
  \item $T_1$: para puntos distintos, cada uno tiene un entorno que no contiene al otro.
  \item $T_2$ (Hausdorff): puntos distintos tienen entornos disjuntos.
  \item $T_3$ (regular): cerrado y punto fuera de él tienen entornos disjuntos.
  \item $T_4$ (normal): dos cerrados disjuntos tienen entornos disjuntos.
\end{itemize}
\end{definition}

\begin{theorem}[Urysohn]
\label{teo:urysohn-general}
En un espacio normal, cualesquiera dos cerrados disjuntos pueden separarse por una función continua $f\colon X\to[0,1]$.
\end{theorem}

\begin{theorem}[Tietze]
\label{teo:tietze-general}
En un espacio normal, toda función continua definida en un cerrado y con valores en $\mathbb{R}$ puede extenderse continuamente a todo $X$.
\end{theorem}

\section{Productos de espacios}

\begin{definition}
\label{def:producto-topologico}
La \emph{topología producto} sobre $\prod_{i\in I}X_i$ es la generada por las proyecciones $\pi_j$.
\end{definition}

\section{Teorema de Tychonoff (panorama)}

\begin{theorem}[Tychonoff]
\label{teo:tychonoff}
El producto arbitrario de espacios compactos es compacto en la topología producto.
\end{theorem}

Este teorema, uno de los más profundos de la topología general, es equivalente al axioma de elección. Su demostración requiere lemas auxiliares (Alexander subbase theorem) y está fuera del alcance elemental de este capítulo.

\begin{exercise}
Demuestre que la topología cofinita sobre $\mathbb{N}$ es compacta pero no Hausdorff.
\end{exercise}

\begin{exercise}
Pruebe que un espacio es Hausdorff si y solo si la diagonal $\Delta=\{(x,x):\,x\in X\}$ es cerrada en $X\times X$.
\end{exercise}

\begin{exercise}
Demuestre el lema de Urysohn en espacios métricos construyendo $f(x)=d(x,A)/(d(x,A)+d(x,B))$.
\end{exercise}

\begin{problem}
Demuestre que todo espacio compacto Hausdorff es normal.
\end{problem}

\chapter{Introducción a la probabilidad moderna}
\label{ch:probabilidad-moderna}
\index{Espacio de probabilidad}\index{Variable aleatoria}\index{Esperanza}\index{Convergencia en probabilidad}\index{Varianza}\index{Leyes de los grandes números}\index{Teorema central del límite}

\section{Espacios de probabilidad}

La probabilidad moderna, fundada por Andrei Kolmogorov en 1933, no es más que la teoría de la medida aplicada a espacios donde la medida total es $1$. Este enfoque transformó la probabilidad de una colección de técnicas intuitivas en una rama del análisis matemático con definiciones rigurosas, teoremas generales y aplicaciones profundas.

\begin{definition}
\label{def:espacio-probabilidad}
Un \emph{espacio de probabilidad} es un espacio de medida $(\Omega,\mathcal{F},\mathbb{P})$ donde $\mathbb{P}(\Omega)=1$.
\end{definition}

\begin{example}
\label{ej:bernoulli-prob}
El espacio de Bernoulli: $\Omega=\{0,1\}^{\mathbb{N}}$, $\mathcal{F}$ la $\sigma$-álgebra producto, y $\mathbb{P}$ la medida producto con $\mathbb{P}(\omega_n=1)=p$.
\end{example}

\section{Variables aleatorias}

\begin{definition}
\label{def:variable-aleatoria}
Una \emph{variable aleatoria} es una función medible $X\colon\Omega\to\mathbb{R}$ (o a un espacio medible cualquiera).
\end{definition}

\begin{definition}
\label{def:distribucion-prob}
La \emph{distribución} de $X$ es la medida imagen $\mu_X(B)=\mathbb{P}(X^{-1}(B))$ sobre $\mathbb{R}$.
\end{definition}

\section{Esperanza matemática}

\begin{definition}
\label{def:esperanza}
La \emph{esperanza} de una variable aleatoria integrable es
\[
  \mathbb{E}[X]=\int_{\Omega}X(\omega)\,d\mathbb{P}(\omega).
\]
\end{definition}

\begin{theorem}[Propiedades]
\label{teo:prop-esperanza}
La esperanza es lineal, monótona, y $\mathbb{E}[1]=1$.
\end{theorem}

\section{Modos de convergencia}

\begin{definition}
\label{def:convergencias-probabilidad}
\begin{itemize}
  \item $X_n\to X$ \emph{casi seguramente} si $\mathbb{P}(\lim X_n=X)=1$.
  \item $X_n\to X$ \emph{en probabilidad} si $\mathbb{P}(|X_n-X|>\varepsilon)\to 0$ para todo $\varepsilon>0$.
  \item $X_n\to X$ en $L^p$ si $\mathbb{E}[|X_n-X|^p]\to 0$.
  \item $X_n\to X$ \emph{en distribución} si $\mathbb{E}[f(X_n)]\to\mathbb{E}[f(X)]$ para toda $f$ continua acotada.
\end{itemize}
\end{definition}

\begin{proposition}
\label{prop:implicaciones-convergencia}
Convergencia casi segura o en $L^p$ implica convergencia en probabilidad, la cual implica convergencia en distribución.
\end{proposition}

\section{Varianza y desigualdades básicas}

\begin{definition}
\label{def:varianza}
La \emph{varianza} de una variable aleatoria integrable $X$ es
\[
  \operatorname{Var}(X)=\mathbb{E}\bigl[(X-\mathbb{E}[X])^2\bigr]=\mathbb{E}[X^2]-(\mathbb{E}[X])^2,
\]
siempre que estas esperanzas existan.
\end{definition}

\begin{theorem}[Desigualdad de Chebyshev]
\label{teo:chebyshev}
Sea $X$ una variable aleatoria con varianza finita. Para todo $\varepsilon>0$,
\[
  \mathbb{P}(|X-\mathbb{E}[X]|\geq\varepsilon)\leq\frac{\operatorname{Var}(X)}{\varepsilon^2}.
\]
\end{theorem}
\begin{proof}
Aplicamos la desigualdad de Markov $\mathbb{P}(Y\geq a)\leq\mathbb{E}[Y]/a$ (válida para $Y\geq 0$, $a>0$) a $Y=(X-\mathbb{E}[X])^2$ y $a=\varepsilon^2$.
\end{proof}

\begin{corollary}[Desigualdad de Markov]
\label{cor:markov}
Si $X\geq 0$ casi seguramente y $a>0$, entonces $\mathbb{P}(X\geq a)\leq\mathbb{E}[X]/a$.
\end{corollary}
\begin{proof}
Notemos que $a\,\mathbf{1}_{\{X\geq a\}}\leq X$. Tomando esperanza en ambos lados se obtiene el resultado.
\end{proof}

\section{Leyes de los grandes números}

\begin{theorem}[Ley débil]
\label{teo:ley-debil}
Sea $(X_n)$ una sucesión de variables aleatorias independientes e idénticamente distribuidas (i.i.d.) con $\mathbb{E}[X_1]=\mu$ y varianza finita. Entonces
\[
  \frac{1}{n}\sum_{k=1}^{n}X_k\xrightarrow{\mathbb{P}}\mu.
\]
\end{theorem}
\begin{proof}[Bosquejo]
Por la desigualdad de Chebyshev (\cref{teo:chebyshev}),
\[
  \mathbb{P}\Bigl(\Bigl|\frac{S_n}{n}-\mu\Bigr|\geq\varepsilon\Bigr)\leq\frac{\operatorname{Var}(S_n/n)}{\varepsilon^2}=\frac{\sigma^2}{n\varepsilon^2}\to 0.
\]
\end{proof}

\begin{theorem}[Ley fuerte de Kolmogorov]
\label{teo:ley-fuerte}
Bajo las mismas hipótesis, $S_n/n\to\mu$ casi seguramente.
\end{theorem}

\section{Teorema central del límite (introducción)}

\begin{theorem}[Central del límite]
\label{teo:tcl}
Sea $(X_n)$ i.i.d. con $\mathbb{E}[X_1]=\mu$ y $\operatorname{Var}(X_1)=\sigma^2<\infty$. Entonces
\[
  \frac{S_n-n\mu}{\sigma\sqrt{n}}\xrightarrow{d}\mathcal{N}(0,1),
\]
donde $\mathcal{N}(0,1)$ es la distribución normal estándar.
\end{theorem}

\begin{remark}
El TCL explica por qué la distribución normal aparece universalmente en fenómenos que resultan de la suma de muchas contribuciones independientes: errores de medida, fluctuaciones termodinámicas, ruido en señales, etc.
\end{remark}

\begin{exercise}
Demuestre que si $X_n\to X$ en probabilidad, existe una subsucesión que converge casi seguramente.
\end{exercise}

\begin{exercise}
Calcule la esperanza y varianza de una variable aleatoria de Bernoulli.
\end{exercise}

\begin{exercise}
Use la ley débil para estimar la probabilidad de que la frecuencia de caras en $n$ lanzamientos de una moneda honesta difiera de $1/2$ en más de $0.01$.
\end{exercise}

\begin{problem}
Demuestre el teorema central del límite usando funciones características y el teorema de continuidad de Lévy.
\end{problem}

\chapter{Introducción al análisis convexo y la optimización}
\label{ch:analisis-convexo}
\index{Conjunto convexo}\index{Función convexa}\index{Teorema de separación}\index{Optimización}\index{Punto de silla}

\section{Conjuntos convexos}

La convexidad es una propiedad geométrica que permea gran parte del análisis moderno: desigualdades, optimización, análisis funcional, probabilidad y, más recientemente, aprendizaje automático. Hermann Minkowski, a finales del siglo~XIX, inició el estudio sistemático de los cuerpos convexos, y las herramientas que desarrolló ---teoremas de separación, funciones de soporte, poliedros--- son hoy fundamentales.

\begin{definition}
\label{def:convexo}
Un conjunto $C\subseteq V$ en un espacio vectorial es \emph{convexo} si para todo $x,y\in C$ y todo $\lambda\in[0,1]$,
\[
  \lambda x+(1-\lambda)y\in C.
\]
\end{definition}

\begin{example}
\label{ej:convexos-basicos}
Las bolas en espacios normados, los semiespacios, los simplex, y las intersecciones arbitrarias de convexos, son convexos.
\end{example}

\begin{theorem}[Minkowski]
\label{teo:minkowski}
En $\mathbb{R}^n$, todo conjunto convexo compacto es la envolvente convexa de sus puntos extremos.
\end{theorem}

\section{Funciones convexas}

\begin{definition}
\label{def:funcion-convexa}
Una función $f\colon C\to\mathbb{R}$ definida en un convexo es \emph{convexa} si
\[
  f(\lambda x+(1-\lambda)y)\leq\lambda f(x)+(1-\lambda)f(y)
\]
para todo $x,y\in C$ y $\lambda\in[0,1]$.
\end{definition}

\begin{proposition}
\label{prop:convexa-continua}
Toda función convexa en un abierto convexo de $\mathbb{R}^n$ es continua.
\end{proposition}

\begin{theorem}[Desigualdad de Jensen]
\label{teo:jensen}
Si $f$ es convexa y $X$ es una variable aleatoria integrable, entonces $f(\mathbb{E}[X])\leq\mathbb{E}[f(X)]$.
\end{theorem}

\section{Separación de conjuntos convexos}

\begin{theorem}[Hahn--Banach geométrico]
\label{teo:hahn-banach-convexo}
Sean $A,B\subseteq V$ convexos disjuntos no vacíos en un espacio normado, con $A$ abierto. Entonces existe un funcional lineal continuo $\phi$ y $\alpha\in\mathbb{R}$ tales que
\[
  \phi(a)<\alpha\leq\phi(b)\quad\text{para todo }a\in A,\,b\in B.
\]
\end{theorem}

\section{Optimización en espacios normados}

\begin{definition}
\label{def:problema-optimizacion}
Un \emph{problema de optimización convexo} es minimizar una función convexa $f$ sobre un conjunto convexo $C$, es decir, hallar $x^*\in C$ tal que $f(x^*)=\inf_{x\in C}f(x)$.
\end{definition}

\begin{theorem}
\label{teo:minimo-convexo}
Si $f$ es convexa y $C$ es convexo, todo mínimo local es global. Si $f$ es estrictamente convexa, el mínimo es único.
\end{theorem}

\begin{theorem}[Condiciones de optimalidad]
\label{teo:kkt}
Bajo ciertas condiciones de regularidad (qualifications), un punto $x^*$ es mínimo de un problema convexo con restricciones $g_i(x)\leq 0$ si y solo si existen multiplicadores de Lagrange $\lambda_i\geq 0$ tales que
\[
  \nabla f(x^*)+\sum\lambda_i\nabla g_i(x^*)=0,\qquad \lambda_i g_i(x^*)=0.
\]
\end{theorem}

\section{Aplicaciones al aprendizaje automático}

\begin{example}[Regresión logística]
\label{ej:regresion-logistica}
La función de pérdida de la regresión logística es convexa en los parámetros, lo cual garantiza que los métodos de descenso por gradiente converjan a un mínimo global.
\end{example}

\begin{example}[Máquinas de vectores de soporte]
\label{ej:svm}
El problema primal de una SVM es minimizar $\|w\|^2$ sujeto a restricciones lineales, un problema cuadrático convexo que admite solución eficiente y única.
\end{example}

\begin{example}[Descenso por gradiente estocástico]
\label{ej:sgd}
En optimización de funciones convexas, el método SGD converge en media al mínimo global bajo hipótesis de acotación del gradiente y tasas de aprendizaje decrecientes.
\end{example}

\begin{exercise}
Demuestre que la intersección de convexos es convexa, y que la imagen de un convexo por una aplicación afín es convexa.
\end{exercise}

\begin{exercise}
Pruebe que $f(x)=\|x\|^2$ es estrictamente convexa en un espacio de Hilbert.
\end{exercise}

\begin{exercise}
Formule el problema de mínimos cuadrados $\min_x\|Ax-b\|^2$ como un problema convexo y derive la ecuación normal $A^\top Ax=A^\top b$.
\end{exercise}

\begin{problem}
Demuestre que el subdiferencial $\partial f(x)=\{v:\,f(y)\geq f(x)+\langle v,y-x\rangle\text{ para todo }y\}$ de una función convexa es no vacío en todo punto interior del dominio.
\end{problem}


% Apéndices
\appendix
\chapter{Teoría de conjuntos}
\label{ap:teoria-conjuntos}
\index{Conjunto}\index{Pertenencia}\index{Relación de equivalencia}\index{Relación de orden}\index{Función}\index{Familia}\index{Equipotencia}\index{Conjunto numerable}\index{Cardinalidad}

\section{Conjuntos y la crisis de los fundamentos}

La teoría de conjuntos constituye el lenguaje universal de las matemáticas modernas. Georg Cantor, su creador, la describió como ``la ciencia de lo infinito''. Sin embargo, su desarrollo inicial fue turbulenta: la aparición de paradojas lógicas obligó a los matemáticos del siglo~XX a axiomatizarla cuidadosamente. En este capítulo presentamos los elementos esenciales de la teoría ``ingenua'' de conjuntos, suficientes para las necesidades del análisis matemático, y damos un vistazo a la axiomatización de Zermelo--Fraenkel que subyace a la práctica contemporánea.

Un \emph{conjunto} es una colección bien definida de objetos, llamados \emph{elementos}. La pertenencia de un elemento $x$ a un conjunto $A$ se denota $x\in A$; su no pertenencia, $x\notin A$.

\begin{definition}
\label{def:igualdad-conjuntos}
Dos conjuntos $A$ y $B$ son \emph{iguales}, denotado $A=B$, si tienen exactamente los mismos elementos:
\[
  A=B\quad\Longleftrightarrow\quad\forall x:\,(x\in A\Leftrightarrow x\in B).
\]
\end{definition}

La \cref{def:igualdad-conjuntos} encapsula el \emph{principio de extensión}: un conjunto queda completamente determinado por sus elementos, no por el modo en que se describa.

\begin{example}
\label{ej:descripciones}
Los siguientes conjuntos son iguales:
\[
  A=\{1,2,3\},\qquad B=\{x\in\mathbb{Z}:\,1\leq x\leq 3\},\qquad C=\{x\in\mathbb{N}:\,x^2<10\}.
\]
\end{example}

\begin{definition}
\label{def:conjuntos-especiales}
El \emph{conjunto vacío}, denotado $\emptyset$, es el único conjunto que carece de elementos. El \emph{conjunto potencia} de $A$, denotado $\mathcal{P}(A)$, es el conjunto de todos los subconjuntos de $A$:
\[
  \mathcal{P}(A)=\{B:\,B\subseteq A\}.
\]
\end{definition}

\begin{remark}
Si $A$ tiene $n$ elementos, entonces $\mathcal{P}(A)$ tiene $2^n$ elementos. Esta observación será fundamental cuando estudiemos cardinalidad.
\end{remark}

\subsection{La paradoja de Russell}

A finales del siglo~XIX, la concepción ``ingenua'' de conjunto ---cualquier propiedad determina un conjunto--- parecía inocua. Sin embargo, en 1901 Bertrand Russell descubrió la siguiente contradicción. Consideremos el ``conjunto''
\[
  R=\{x : x\notin x\}.
\]
Si $R\in R$, entonces por definición $R\notin R$. Si $R\notin R$, entonces satisface la condición de pertenencia, luego $R\in R$. En ambos casos se llega a una contradicción.

Esta \emph{paradoja de Russell} demostró que no toda descripción define un conjunto. La solución, desarrollada por Ernst Zermelo y Abraham Fraenkel, fue restringir la formación de conjuntos mediante un sistema axiomático (ZF) en el que la existencia de un conjunto requiere justificación explícita a partir de conjuntos ya dados.

\begin{remark}
Los axiomas de Zermelo--Fraenkel (con el axioma de elección, ZFC) son la base de la matemática moderna. Entre ellos destacan: el axioma de extensión (igualdad por elementos), el axioma del conjunto vacío, el axioma de par y unión, el axioma de potencia, el axioma de separación (que restringe la formación de subconjuntos a partir de uno dado, evitando la paradoja de Russell), el axioma de infinito (que garantiza la existencia de $\mathbb{N}$), el axioma de reemplazo y el axioma de regularidad. El lector interesado puede consultar textos especializados de teoría de conjuntos para un tratamiento formal.
\end{remark}

\section{Operaciones entre conjuntos}

\begin{definition}
\label{def:operaciones}
Sean $A$ y $B$ conjuntos contenidos en un universo $U$.
\begin{itemize}
  \item \emph{Unión}: $A\cup B=\{x\in U:\,x\in A\text{ o }x\in B\}$.
  \item \emph{Intersección}: $A\cap B=\{x\in U:\,x\in A\text{ y }x\in B\}$.
  \item \emph{Diferencia}: $A\setminus B=\{x\in U:\,x\in A\text{ y }x\notin B\}$.
  \item \emph{Complemento}: $A^c=U\setminus A$.
  \item \emph{Diferencia simétrica}: $A\triangle B=(A\setminus B)\cup(B\setminus A)$.
\end{itemize}
\end{definition}

\begin{center}
\begin{tikzpicture}[scale=0.75]
  % Universo
  \draw[thick,gray] (-0.5,-0.5) rectangle (8.5,3.5) node[below left] {$U$};

  % Union A cup B
  \begin{scope}[shift={(0,0)}]
    \draw[fill=blue!20,thick] (1,1.5) circle (1.2cm);
    \draw[fill=blue!20,thick] (2.2,1.5) circle (1.2cm);
    \node at (1.6,1.5) {$A\cup B$};
    \node at (1,2.9) {$A$};
    \node at (2.2,2.9) {$B$};
  \end{scope}

  % Intersection A cap B
  \begin{scope}[shift={(3.5,0)}]
    \draw[thick] (1,1.5) circle (1.2cm);
    \draw[thick] (2.2,1.5) circle (1.2cm);
    \begin{scope}
      \clip (1,1.5) circle (1.2cm);
      \fill[red!30] (2.2,1.5) circle (1.2cm);
    \end{scope}
    \node at (1.6,1.5) {$A\cap B$};
    \node at (1,2.9) {$A$};
    \node at (2.2,2.9) {$B$};
  \end{scope}

  % Complement
  \begin{scope}[shift={(7,0)}]
    \fill[yellow!20] (0.3,0.3) rectangle (2.9,2.7);
    \draw[fill=white,thick] (1.6,1.5) circle (0.9cm);
    \node at (1.6,1.5) {$A$};
    \node at (2.4,2.4) {$A^c$};
  \end{scope}
\end{tikzpicture}
\end{center}

\begin{proposition}[Propiedades algebraicas]
\label{prop:algebra-conjuntos}
Para cualesquiera $A,B,C\subseteq U$:
\begin{enumerate}
  \item[(i)] \emph{Conmutatividad}: $A\cup B=B\cup A$, $A\cap B=B\cap A$.
  \item[(ii)] \emph{Asociatividad}: $(A\cup B)\cup C=A\cup(B\cup C)$, $(A\cap B)\cap C=A\cap(B\cap C)$.
  \item[(iii)] \emph{Distributividad}: $A\cup(B\cap C)=(A\cup B)\cap(A\cup C)$, $A\cap(B\cup C)=(A\cap B)\cup(A\cap C)$.
  \item[(iv)] \emph{Leyes de De Morgan}: $(A\cup B)^c=A^c\cap B^c$, $(A\cap B)^c=A^c\cup B^c$.
  \item[(v)] \emph{Absorción}: $A\cup(A\cap B)=A$, $A\cap(A\cup B)=A$.
  \item[(vi)] \emph{Idempotencia}: $A\cup A=A$, $A\cap A=A$.
\end{enumerate}
\end{proposition}

\begin{proof}
Cada igualdad se demuestra verificando la doble inclusión. Ilustramos con (iv): si $x\in(A\cup B)^c$, entonces $x\notin A\cup B$, luego $x\notin A$ y $x\notin B$, es decir, $x\in A^c\cap B^c$. El recíproco es análogo.
\end{proof}

\begin{center}
\begin{tikzpicture}[scale=0.8]
  % Ley de De Morgan: (A cup B)^c = A^c cap B^c
  \begin{scope}[shift={(0,0)}]
    \fill[gray!20] (-0.3,-0.3) rectangle (3.3,2.7);
    \draw[thick] (-0.3,-0.3) rectangle (3.3,2.7) node[below left] {$U$};
    \draw[fill=white,thick] (1.2,1.2) circle (0.9cm);
    \draw[fill=white,thick] (1.8,1.2) circle (0.9cm);
    \node at (1.5,1.2) {$A\cup B$};
    \fill[blue!30] (-0.3,-0.3) rectangle (3.3,2.7);
    \begin{scope}
      \clip (1.2,1.2) circle (0.9cm);
      \fill[gray!20] (1.8,1.2) circle (0.9cm);
    \end{scope}
    \begin{scope}
      \clip (1.8,1.2) circle (0.9cm);
      \fill[gray!20] (1.2,1.2) circle (0.9cm);
    \end{scope}
    \draw[thick] (1.2,1.2) circle (0.9cm);
    \draw[thick] (1.8,1.2) circle (0.9cm);
    \node at (1.5,-0.1) {$(A\cup B)^c$};
  \end{scope}

  \node at (4.5,1.2) {$=$};

  \begin{scope}[shift={(6,0)}]
    \fill[gray!20] (-0.3,-0.3) rectangle (3.3,2.7);
    \draw[thick] (-0.3,-0.3) rectangle (3.3,2.7) node[below left] {$U$};
    \draw[thick] (1.2,1.2) circle (0.9cm);
    \draw[thick] (1.8,1.2) circle (0.9cm);
    \fill[blue!30] (-0.3,-0.3) rectangle (3.3,2.7);
    \begin{scope}
      \clip (1.2,1.2) circle (0.9cm);
      \fill[gray!20] (1.8,1.2) circle (0.9cm);
    \end{scope}
    \begin{scope}
      \clip (1.8,1.2) circle (0.9cm);
      \fill[gray!20] (1.2,1.2) circle (0.9cm);
    \end{scope}
    \draw[thick] (1.2,1.2) circle (0.9cm);
    \draw[thick] (1.8,1.2) circle (0.9cm);
    \node at (0.3,1.2) {$A^c$};
    \node at (2.7,1.2) {$B^c$};
    \node at (1.5,-0.1) {$A^c\cap B^c$};
  \end{scope}
\end{tikzpicture}
\end{center}

\begin{definition}
\label{def:producto-cartesiano}
El \emph{producto cartesiano} de $A$ y $B$ es
\[
  A\times B=\{(a,b):\,a\in A,\,b\in B\}.
\]
\end{definition}

El producto cartesiano se extiende a familias finitas $A_1\times\cdots\times A_n$, y a familias infinitas en el contexto de la axiomática de conjuntos. Dos pares ordenados $(a,b)$ y $(c,d)$ son iguales si y solo si $a=c$ y $b=d$.

\section{Relaciones}

\begin{definition}
\label{def:relacion}
Una \emph{relación} de $A$ en $B$ es un subconjunto $R\subseteq A\times B$. Si $(a,b)\in R$, escribimos $a\,R\,b$.
\end{definition}

\begin{definition}
\label{def:relacion-equivalencia}
Una relación $R$ en $A$ es una \emph{relación de equivalencia} si es:
\begin{enumerate}
  \item[(i)] \emph{Reflexiva}: $\forall a\in A:\,a\,R\,a$.
  \item[(ii)] \emph{Simétrica}: $\forall a,b\in A:\,a\,R\,b\Rightarrow b\,R\,a$.
  \item[(iii)] \emph{Transitiva}: $\forall a,b,c\in A:\,a\,R\,b\text{ y }b\,R\,c\Rightarrow a\,R\,c$.
\end{enumerate}
\end{definition}

Las relaciones de equivalencia particionan un conjunto en \emph{clases de equivalencia} disjuntas. Dado $a\in A$, su clase de equivalencia es $[a]=\{x\in A:\,x\,R\,a\}$. El conjunto de todas las clases se denomina \emph{conjunto cociente} $A/R$.

\begin{example}
La relación ``tener la misma paridad'' en $\mathbb{Z}$ es de equivalencia. Las clases son los números pares y los impares. El conjunto cociente tiene dos elementos: $\mathbb{Z}/\text{paridad}\cong\{0,1\}$.
\end{example}

\begin{example}
En la construcción de Cantor de $\mathbb{R}$, dos sucesiones de Cauchy de racionales se declaran equivalentes si su diferencia tiende a cero. El conjunto cociente es precisamente $\mathbb{R}$. Este ejemplo muestra que las relaciones de equivalencia no son un mero formalismo: son la herramienta constructiva para ``completar'' $\mathbb{Q}$.
\end{example}

\begin{definition}
\label{def:relacion-orden}
Una relación $R$ en $A$ es un \emph{orden parcial} si es reflexiva, transitiva y \emph{antisimétrica} ($a\,R\,b$ y $b\,R\,a\Rightarrow a=b$). Si además cualesquiera dos elementos son comparables, se dice que es un \emph{orden total}.
\end{definition}

\begin{example}
La inclusión $\subseteq$ entre subconjuntos de un conjunto dado es un orden parcial, pero no total (si $A$ tiene al menos dos elementos, existen subconjuntos no comparables).
\end{example}

\begin{example}
El orden usual $\leq$ en $\mathbb{R}$ es un orden total. Este hecho no es trivial: depende de la construcción de $\mathbb{R}$ y de la definición de $<$ en términos de cortaduras o de sucesiones.
\end{example}

\begin{definition}
Sea $\leq$ un orden parcial en $A$ y sea $B\subseteq A$.
\begin{itemize}
  \item Un elemento $m\in B$ es \emph{minimal} en $B$ si no existe $b\in B$ con $b<m$.
  \item Un elemento $M\in A$ es \emph{cota superior} de $B$ si $b\leq M$ para todo $b\in B$.
  \item Un elemento $s\in A$ es \emph{supremo} de $B$ si es la mínima cota superior.
\end{itemize}
\end{definition}

\begin{remark}[Lema de Zorn]
Si todo subconjunto totalmente ordenado de un conjunto parcialmente ordenado $A$ tiene una cota superior en $A$, entonces $A$ posee al menos un elemento maximal. Este resultado, equivalente al axioma de elección, se utiliza en el análisis funcional para probar la existencia de bases de Hamel y de extensiones de funcionales lineales (teorema de Hahn--Banach).
\end{remark}

\section{Funciones}

Las funciones son el objeto central del análisis matemático. Desde el punto de vista de la teoría de conjuntos, una función no es más que una relación con una propiedad adicional.

\begin{definition}
\label{def:funcion}
Una \emph{función} $f\colon A\to B$ es una relación $f\subseteq A\times B$ tal que para cada $a\in A$ existe un único $b\in B$ con $(a,b)\in f$. Este único $b$ se denota $f(a)$.
\end{definition}

\begin{definition}
\label{def:inyectiva-sobreyectiva}
Sea $f\colon A\to B$.
\begin{itemize}
  \item $f$ es \emph{inyectiva} si $f(a_1)=f(a_2)\Rightarrow a_1=a_2$.
  \item $f$ es \emph{sobreyectiva} si $\forall b\in B\ \exists a\in A:\,f(a)=b$.
  \item $f$ es \emph{biyectiva} si es inyectiva y sobreyectiva.
\end{itemize}
\end{definition}

\begin{definition}
\label{def:imagen-preimagen}
Sea $f\colon A\to B$ y sean $X\subseteq A$, $Y\subseteq B$.
\begin{itemize}
  \item La \emph{imagen directa} de $X$ es $f(X)=\{f(x):\,x\in X\}$.
  \item La \emph{imagen inversa} de $Y$ es $f^{-1}(Y)=\{x\in A:\,f(x)\in Y\}$.
\end{itemize}
\end{definition}

\begin{proposition}
\label{prop:imagen-inversa}
La imagen inversa preserva uniones, intersecciones y complementos:
\begin{enumerate}
  \item[(i)] $f^{-1}(Y_1\cup Y_2)=f^{-1}(Y_1)\cup f^{-1}(Y_2)$.
  \item[(ii)] $f^{-1}(Y_1\cap Y_2)=f^{-1}(Y_1)\cap f^{-1}(Y_2)$.
  \item[(iii)] $f^{-1}(Y^c)=(f^{-1}(Y))^c$.
\end{enumerate}
\end{proposition}

\begin{proof}
La verificación es elemental usando la definición de imagen inversa. Por ejemplo, para (ii): $x\in f^{-1}(Y_1\cap Y_2)\Leftrightarrow f(x)\in Y_1\cap Y_2\Leftrightarrow f(x)\in Y_1\text{ y }f(x)\in Y_2\Leftrightarrow x\in f^{-1}(Y_1)\cap f^{-1}(Y_2)$.
\end{proof}

\begin{remark}
La imagen directa, en cambio, no preserva intersecciones en general: puede ocurrir que $f(X_1\cap X_2)\subsetneq f(X_1)\cap f(X_2)$. La igualdad sí se da si $f$ es inyectiva.
\end{remark}

\section{Familias de conjuntos}

Cuando necesitamos considerar una colección indexada de conjuntos, utilizamos la noción de \emph{familia}.

\begin{definition}
\label{def:familia}
Una \emph{familia de conjuntos} indexada por un conjunto $I$ es una función $A\colon I\to\mathcal{P}(U)$, habitualmente denotada $(A_i)_{i\in I}$, donde $A_i=A(i)$.
\end{definition}

\begin{definition}
\label{def:union-interseccion-arbitraria}
Las uniones e intersecciones arbitrarias se definen por:
\[
  \bigcup_{i\in I} A_i=\{x\in U:\,\exists i\in I,\,x\in A_i\},\qquad
  \bigcap_{i\in I} A_i=\{x\in U:\,\forall i\in I,\,x\in A_i\}.
\]
\end{definition}

\begin{theorem}[Leyes de De Morgan generalizadas]
\label{teo:demorgan-general}
\[
  \left(\bigcup_{i\in I} A_i\right)^c=\bigcap_{i\in I} A_i^c,\qquad
  \left(\bigcap_{i\in I} A_i\right)^c=\bigcup_{i\in I} A_i^c.
\]
\end{theorem}

\begin{proof}
Para la primera igualdad: $x\in(\bigcup A_i)^c\Leftrightarrow x\notin\bigcup A_i\Leftrightarrow\forall i,\,x\notin A_i\Leftrightarrow\forall i,\,x\in A_i^c\Leftrightarrow x\in\bigcap A_i^c$. La segunda es análoga.
\end{proof}

\begin{definition}
Una familia $(A_i)_{i\in I}$ es \emph{disjunta dos a dos} si $A_i\cap A_j=\emptyset$ para todo $i\neq j$. Una \emph{partición} de $A$ es una familia de subconjuntos no vacíos, disjunta dos a dos, cuya unión es $A$.
\end{definition}

\begin{proposition}
Toda relación de equivalencia en $A$ determina una partición de $A$ (las clases de equivalencia), y toda partición de $A$ determina una relación de equivalencia (``pertenecer al mismo bloque de la partición'').
\end{proposition}

\section{Cardinalidad}

La cardinalidad es la medida del ``tamaño'' de un conjunto en el sentido de la teoría de conjuntos. Dos conjuntos tienen el mismo cardinal si existe una biyección entre ellos.

\begin{definition}
\label{def:equipotencia}
Dos conjuntos $A$ y $B$ son \emph{equipotentes}, denotado $A\sim B$, si existe una biyección $f\colon A\to B$.
\end{definition}

\begin{theorem}[Cantor--Bernstein--Schröder]
\label{teo:cbs}
Si existen inyecciones $f\colon A\to B$ y $g\colon B\to A$, entonces existe una biyección entre $A$ y $B$.
\end{theorem}

\begin{proof}[Bosquejo]
Se define una sucesión de conjuntos $C_0=A\setminus g(B)$, $C_{n+1}=g(f(C_n))$, y se descompone $A$ en $C=\bigcup C_n$ y su complemento. La biyección se construye usando $f$ en $C$ y $g^{-1}$ fuera de $C$.
\end{proof}

\begin{definition}
El \emph{cardinal} de un conjunto $A$, denotado $|A|$ o $\#A$, es la clase de equivalencia de $A$ bajo la relación de equipotencia. En la práctica, se identifica con un \emph{número cardinal} representativo:
\begin{itemize}
  \item $0=|\emptyset|$, $1=|\{0\}|$, $2=|\{0,1\}|$, etc.
  \item $\aleph_0=|\mathbb{N}|$ (el cardinal de todo conjunto infinito numerable).
  \item $\mathfrak{c}=|\mathbb{R}|$ (el cardinal del continuo).
\end{itemize}
\end{definition}

\begin{theorem}[Cantor]
\label{teo:cantor-potencia}
Para todo conjunto $A$, no existe ninguna biyección entre $A$ y $\mathcal{P}(A)$. En particular, $\mathcal{P}(A)$ tiene cardinalidad estrictamente mayor que $A$.
\end{theorem}

\begin{proof}
Si existiera $f\colon A\to\mathcal{P}(A)$ biyectiva, consideremos $B=\{x\in A:\,x\notin f(x)\}$. Si $B=f(y)$ para algún $y$, entonces $y\in B\Leftrightarrow y\notin f(y)=B$, contradicción.
\end{proof}

\begin{corollary}
La jerarquía de cardinales es infinita: $\aleph_0 < 2^{\aleph_0} < 2^{2^{\aleph_0}} < \cdots$. En particular, no hay un ``conjunto de todos los conjuntos'' (pues si existiera, su cardinalidad sería la máxima, pero su potencia sería mayor).
\end{corollary}

\begin{remark}
El teorema de Cantor implica que $|\mathcal{P}(\mathbb{N})|=2^{\aleph_0}=\mathfrak{c}$. Una pregunta natural es si existen cardinales estrictamente entre $\aleph_0$ y $\mathfrak{c}$. La \emph{hipótesis del continuo} (HC) afirma que no existe tal cardinal. Gödel (1938) demostró que HC no puede refutarse a partir de ZFC; Cohen (1963) demostró que tampoco puede demostrarse a partir de ZFC. La HC es, por tanto, \emph{independiente} de los axiomas estándar.
\end{remark}

\section{Conjuntos numerables}

\begin{definition}
\label{def:numerable}
Un conjunto $A$ es \emph{numerable} (o \emph{contable}) si es finito o equipotente a $\mathbb{N}$. Equivalentemente, si sus elementos pueden enumerarse en una sucesión.
\end{definition}

\begin{theorem}
\label{teo:q-numerable}
El conjunto $\mathbb{Q}$ de los números racionales es numerable.
\end{theorem}
\begin{proof}
Se puede enumerar $\mathbb{Q}$ diagonalmente: $\frac{1}{1},\frac{2}{1},\frac{1}{2},\frac{3}{1},\frac{2}{2},\frac{1}{3},\dots$, eliminando las fracciones no reducidas.
\end{proof}

\begin{theorem}
\label{teo:union-numerable}
La unión numerable de conjuntos numerables es numerable.
\end{theorem}

\begin{proof}
Sea $A_n$ numerable para cada $n\in\mathbb{N}$. Escribamos $A_n=\{a_{n1},a_{n2},a_{n3},\dots\}$. La aplicación $(n,k)\mapsto a_{nk}$ es sobreyectiva sobre $\bigcup A_n$. Como $\mathbb{N}\times\mathbb{N}$ es numerable (por diagonalización), la unión lo es.
\end{proof}

\begin{corollary}
El conjunto de los números algebraicos (raíces de polinomios con coeficientes enteros) es numerable.
\end{corollary}

\begin{proof}
Para cada grado $n$, el conjunto de polinomios de grado $n$ con coeficientes enteros es numerable (pues $\mathbb{Z}^{n+1}$ es numerable). Cada polinomio tiene a lo sumo $n$ raíces complejas. La unión numerable de conjuntos finitos es numerable.
\end{proof}

\section{Conjuntos no numerables}

\begin{theorem}[Cantor]
\label{teo:cantor-uncountable}
El conjunto $\mathbb{R}$ de los números reales no es numerable.
\end{theorem}
\begin{proof}
Supongamos que $\mathbb{R}$ fuera numerable y sea $(x_n)$ una enumeración de $[0,1]$. Construimos por inducción una sucesión de intervalos cerrados encajados $I_n$ tales que $x_n\notin I_n$. La intersección $\bigcap I_n$ es no vacía por completitud, pero no contiene ningún $x_n$, contradicción.
\end{proof}

\begin{theorem}[Argumento diagonal de Cantor]
El conjunto de todas las sucesiones de ceros y unos, $\{0,1\}^{\mathbb{N}}$, no es numerable.
\end{theorem}
\begin{proof}
Si $(s_n)$ fuera una enumeración, construimos la sucesión $d$ tal que $d_n\neq s_{n}(n)$. Entonces $d$ difiere de cada $s_n$ en la posición $n$, luego no está en la enumeración.
\end{proof}

\begin{remark}
El argumento diagonal es una de las técnicas más profundas de las matemáticas. Gödel lo adaptó para demostrar sus teoremas de incompletitud, y Turing lo utilizó para probar la indecidibilidad del problema de la parada.
\end{remark}

\section{El axioma de elección}

\begin{axiom}[Axioma de elección]
Dada cualquier familia $(A_i)_{i\in I}$ de conjuntos no vacíos, existe una función $f\colon I\to\bigcup A_i$ tal que $f(i)\in A_i$ para todo $i\in I$.
\end{axiom}

El axioma de elección (AE) afirma que podemos elegir simultáneamente un elemento de cada conjunto de una familia arbitraria. Para familias finitas, esto es trivial; para familias infinitas, no puede deducirse de los demás axiomas de ZF. Su aceptación tiene consecuencias sorprendentes:

\begin{itemize}
  \item \emph{Teorema de la buena ordenación} (Zermelo, 1904): todo conjunto puede bien ordenarse.
  \item \emph{Lema de Zorn}: todo conjunto parcialmente ordenado en el que toda cadena tiene cota superior posee un elemento maximal.
  \item \emph{Existencia de conjuntos no medibles} (Vitali, 1905): en $\mathbb{R}$ existen conjuntos que no son Lebesgue-medibles.
\end{itemize}

El axioma de elección es, por tanto, una ``espada de doble filo'': permite demostrar existencia de objetos que no pueden construirse explícitamente. En el análisis matemático, la mayoría de los textos adoptan ZFC (ZF + AE) como marco de trabajo.

\section{Ejercicios}

\begin{exercise}
Demuestre que $A\times B\sim B\times A$ y que $(A\times B)\times C\sim A\times(B\times C)$.
\end{exercise}

\begin{exercise}
Pruebe que el conjunto de los números algebraicos es numerable, y deduzca que existen números trascendentes.
\end{exercise}

\begin{exercise}
Muestre que la familia de los intervalos abiertos disjuntos de $\mathbb{R}$ es a lo sumo numerable.
\end{exercise}

\begin{exercise}
Sea $R$ una relación de equivalencia en $A$. Demuestre que las clases de equivalencia forman una partición de $A$, y recíprocamente, que toda partición define una relación de equivalencia cuyas clases son los bloques de la partición.
\end{exercise}

\begin{exercise}
Demuestre que si $f\colon A\to B$ es inyectiva, entonces $|A|\leq|B|$. Si es sobreyectiva, entonces $|B|\leq|A|$. Use el teorema de Cantor--Bernstein--Schröder para concluir que si ambas cosas ocurren, $|A|=|B|$.
\end{exercise}

\begin{exercise}
Demuestre que $\mathcal{P}(\mathbb{N})\sim\mathbb{R}$. Sugerencia: asocie a cada subconjunto $A\subseteq\mathbb{N}$ la serie $\sum_{n\in A}2^{-n}$.
\end{exercise}

\begin{exercise}
Sea $A$ un conjunto infinito. Demuestre que $A$ contiene un subconjunto numerable. Deduzca que $\aleph_0$ es el cardinal infinito más pequeño.
\end{exercise}

\begin{exercise}
Pruebe que la composición de funciones inyectivas es inyectiva, y que la composición de funciones sobreyectivas es sobreyectiva.
\end{exercise}

\begin{exercise}
Sea $f\colon A\to B$. Demuestre que $f$ es inyectiva si y solo si existe $g\colon B\to A$ tal que $g\circ f=\operatorname{id}_A$.
\end{exercise}

\begin{exercise}
Demuestre que $|\mathbb{R}^2|=|\mathbb{R}|$. Sugerencia: use la representación decimal y un argumento de intercalación de dígitos.
\end{exercise}

\begin{exercise}
Formalice la paradoja de Russell en el lenguaje de la teoría de conjuntos, y explique por qué el axioma de separación de ZF la evita.
\end{exercise}

\begin{problem}
Sea $\mathcal{F}$ una familia de subconjuntos de $\mathbb{N}$ tal que cualesquiera dos elementos de $\mathcal{F}$ tienen intersección finita. ¿Puede $\mathcal{F}$ ser no numerable? Construya un ejemplo o demuestre que es imposible.
\end{problem}

\begin{problem}
Demuestre que todo conjunto infinito $A$ satisface $A\times A\sim A$. (Este resultado requiere el axioma de elección.)
\end{problem}

\begin{problem}
Sea $\mathcal{C}$ el conjunto de todas las funciones continuas $f\colon[0,1]\to\mathbb{R}$. Demuestre que $|\mathcal{C}|=\mathfrak{c}$. (Sugerencia: una función continua queda determinada por sus valores en $\mathbb{Q}\cap[0,1]$.)
\end{problem}

\chapter{Los números reales}
\label{ap:numeros-reales}
\index{Números reales}\index{Cortadura de Dedekind}\index{Sucesión de Cauchy}\index{Supremo}\index{Ínfimo}\index{Valor absoluto}\index{Propiedad arquimediana}\index{Densidad de los racionales}

\section{Construcción de los sistemas numéricos}

Antes de estudiar las propiedades de $\mathbb{R}$, conviene comprender cómo se construyen sistemáticamente los números a partir de los naturales. Cada extensión resuelve una ecuación que el sistema anterior no podía manejar.

\begin{center}
\begin{tikzpicture}[node distance=2.2cm, auto]
  \node[draw,rounded corners,fill=green!15] (N) {$\mathbb{N}$};
  \node[draw,rounded corners,fill=blue!15,right of=N] (Z) {$\mathbb{Z}$};
  \node[draw,rounded corners,fill=orange!15,right of=Z] (Q) {$\mathbb{Q}$};
  \node[draw,rounded corners,fill=red!15,right of=Q] (R) {$\mathbb{R}$};
  \draw[->,thick] (N) -- (Z) node[midway,above] {$x+a=b$} node[midway,below] {resta};
  \draw[->,thick] (Z) -- (Q) node[midway,above] {$ax=b$} node[midway,below] {división};
  \draw[->,thick] (Q) -- (R) node[midway,above] {huecos} node[midway,below] {completitud};
\end{tikzpicture}
\end{center}

\subsection{De \texorpdfstring{$\mathbb{N}$}{N} a \texorpdfstring{$\mathbb{Z}$}{Z}}

Los números naturales $\mathbb{N}$ se axiomatizan mediante los axiomas de Peano. En $\mathbb{N}$, la ecuación $x+a=b$ solo tiene solución cuando $a\leq b$. Para resolverla en general, se define $\mathbb{Z}$ como el conjunto de pares $(a,b)\in\mathbb{N}\times\mathbb{N}$ bajo la relación de equivalencia
\[
  (a,b)\sim(c,d)\quad\Longleftrightarrow\quad a+d=b+c.
\]
Intuitivamente, el par $(a,b)$ representa la ``diferencia'' $a-b$. La clase de $(a,b)$ se identifica con $a-b$, y las operaciones se definen representante a representante:
\[
  [(a,b)]+[(c,d)]=[(a+c,b+d)],\qquad [(a,b)]\cdot[(c,d)]=[(ac+bd,ad+bc)].
\]

\subsection{De \texorpdfstring{$\mathbb{Z}$}{Z} a \texorpdfstring{$\mathbb{Q}$}{Q}}

En $\mathbb{Z}$, la ecuación $ax=b$ ($a\neq 0$) no siempre tiene solución. Se define $\mathbb{Q}$ como el conjunto de pares $(p,q)\in\mathbb{Z}\times\mathbb{Z}$ con $q\neq 0$, bajo la relación
\[
  (p,q)\sim(r,s)\quad\Longleftrightarrow\quad ps=qr.
\]
La clase de $(p,q)$ se denota $\frac{p}{q}$, y las operaciones de cuerpo se definen de manera natural.

\subsection{De \texorpdfstring{$\mathbb{Q}$}{Q} a \texorpdfstring{$\mathbb{R}$}{R}}

Los números reales constituyen el escenario fundamental del análisis matemático. A diferencia de los números racionales, que surgen naturalmente de las operaciones aritméticas, los reales son una construcción intelectual cuyo objetivo es ``llenar los huecos'' de la recta racional. Presentamos dos construcciones clásicas ---las de Dedekind y Cantor--- y desarrollamos las propiedades que hacen de $\mathbb{R}$ un cuerpo ordenado completo.

\subsubsection{Cortaduras de Dedekind}

\begin{definition}
\label{def:cortadura}
Una \emph{cortadura} de Dedekind es un par $(A,B)$ de subconjuntos no vacíos de $\mathbb{Q}$ tales que:
\begin{enumerate}
  \item[(i)] $A\cup B=\mathbb{Q}$ y $A\cap B=\emptyset$.
  \item[(ii)] Si $a\in A$ y $q\in\mathbb{Q}$ con $q<a$, entonces $q\in A$.
  \item[(iii)] $A$ no tiene máximo.
\end{enumerate}
\end{definition}

La condición (iii) es la que distingue una cortadura de una simple partición: el ``corte'' ocurre donde debería estar un número real que no es racional.

\begin{example}
La cortadura que define $\sqrt{2}$ está dada por $A=\{q\in\mathbb{Q}:\,q^2<2\text{ o }q<0\}$ y $B=\mathbb{Q}\setminus A$.
\end{example}

\begin{definition}
\label{def:reales-dedekind}
El conjunto $\mathbb{R}$ de los números reales es el conjunto de todas las cortaduras de Dedekind.
\end{definition}

\subsubsection{Sucesiones de Cauchy}

\begin{definition}
\label{def:cauchy-racionales}
Una sucesión $(q_n)$ de números racionales es de \emph{Cauchy} si para todo $\varepsilon\in\mathbb{Q}^+$ existe $N\in\mathbb{N}$ tal que $|q_m-q_n|<\varepsilon$ siempre que $m,n\geq N$.
\end{definition}

\begin{definition}
\label{def:reales-cantor}
Definimos $\mathbb{R}$ como el conjunto de clases de equivalencia de sucesiones de Cauchy de racionales, donde $(q_n)\sim (r_n)$ si $\lim_{n\to\infty}|q_n-r_n|=0$.
\end{definition}

Ambas construcciones producen cuerpos ordenados isomorfos. La elección entre una u otra es cuestión de gusto; en el análisis funcional, la perspectiva de Cauchy (completitud métrica) es la que generaliza más naturalmente.

\section{Propiedades algebraicas y de orden}

\begin{theorem}
\label{teo:cuerpo-ordenado}
$\mathbb{R}$ es un \emph{cuerpo ordenado}: satisface los axiomas de cuerpo (asociatividad, conmutatividad, distributividad, elementos neutros e inversos) y los de orden (compatibilidad de la suma y el producto con el orden).
\end{theorem}

\begin{proof}[Bosquejo]
Las operaciones se definen término a término en las sucesiones de Cauchy (o mediante las cortaduras), y se verifica que las clases de equivalencia heredan las propiedades de $\mathbb{Q}$.
\end{proof}

\section{Suprema, ínfimos y completitud}

\begin{definition}
\label{def:supremo}
Sea $A\subseteq\mathbb{R}$ no vacío. Un número $M\in\mathbb{R}$ es una \emph{cota superior} de $A$ si $a\leq M$ para todo $a\in A$. El \emph{supremo} de $A$, denotado $\sup A$, es la mínima cota superior de $A$, si existe. Análogamente se define el \emph{ínfimo} $\inf A$ como la máxima cota inferior.
\end{definition}

\begin{theorem}[Completitud de Dedekind]
\label{teo:completitud-dedekind}
Todo subconjunto no vacío de $\mathbb{R}$ acotado superiormente posee supremo.
\end{theorem}

\begin{proof}
Dado $A\subseteq\mathbb{R}$ acotado superiormente, consideremos el conjunto $S$ de todas las cotas superiores de $A$. El par $(\mathbb{R}\setminus S, S)$ constituye una cortadura de Dedekind; el número real que representa es, por construcción, $\sup A$.
\end{proof}

\begin{theorem}[Caracterización del supremo]
\label{teo:carac-sup}
Sea $A\subseteq\mathbb{R}$ no vacío y $s\in\mathbb{R}$. Entonces $s=\sup A$ si y solo si:
\begin{enumerate}
  \item[(i)] $s$ es cota superior de $A$;
  \item[(ii)] para todo $\varepsilon>0$, existe $a\in A$ tal que $s-\varepsilon<a\leq s$.
\end{enumerate}
\end{theorem}

\begin{proof}
Si $s=\sup A$, entonces (i) es inmediata. Si (ii) fallara para algún $\varepsilon>0$, entonces $s-\varepsilon$ sería cota superior, contradiciendo la minimalidad de $s$. Recíprocamente, si $s$ satisface (i) y (ii), cualquier cota superior $M<s$ no puede cumplir (ii) para $\varepsilon=s-M$.
\end{proof}

\section{Valor absoluto y desigualdades fundamentales}

\begin{definition}
\label{def:valor-absoluto}
El \emph{valor absoluto} de $x\in\mathbb{R}$ se define por
\[
  |x|=\begin{cases}x & \text{si }x\geq 0,\\ -x & \text{si }x<0.\end{cases}
\]
\end{definition}

\begin{proposition}
\label{prop:prop-abs}
Para todo $x,y\in\mathbb{R}$:
\begin{enumerate}
  \item[(i)] $|x|\geq 0$, y $|x|=0\Leftrightarrow x=0$.
  \item[(ii)] $|xy|=|x|\,|y|$.
  \item[(iii)] $|x+y|\leq |x|+|y|$ (desigualdad triangular).
  \item[(iv)] $||x|-|y||\leq |x-y|$.
\end{enumerate}
\end{proposition}

\begin{proof}
La propiedad (iii) es la más relevante. Si $x+y\geq 0$, entonces $|x+y|=x+y\leq|x|+|y|$. Si $x+y<0$, entonces $|x+y|=-(x+y)=(-x)+(-y)\leq|x|+|y|$.
\end{proof}

\section{Propiedad arquimediana}

\begin{theorem}[Propiedad arquimediana]
\label{teo:arquimediana}
Para cualesquiera $x,y\in\mathbb{R}$ con $x>0$, existe $n\in\mathbb{N}$ tal que $nx>y$.
\end{theorem}

\begin{proof}
Si no existiera tal $n$, entonces $y$ sería cota superior de $\{nx:\,n\in\mathbb{N}\}$. Sea $s=\sup\{nx:\,n\in\mathbb{N}\}$. Entonces $s-x<s$, luego existe $n$ tal que $nx>s-x$, es decir, $(n+1)x>s$, contradicción.
\end{proof}

\begin{center}
\begin{tikzpicture}[scale=1.0]
  \draw[->] (-0.5,0) -- (7.5,0) node[right] {$\mathbb{R}$};
  \foreach \x/\n in {0/0,1/x,2/2x,3/3x}
    \draw (\x,0.15) -- (\x,-0.15) node[below] {\n};
  \draw[dashed] (5,0.15) -- (5,-0.15) node[below] {$nx$};
  \draw[thick,red] (6,0.15) -- (6,-0.15) node[below] {$y$};
  \draw[decorate,decoration={brace,amplitude=5pt}] (0,0.4) -- (1,0.4) node[midway,above=5pt] {$x$};
  \draw[->,thick,blue] (4.8,0.3) -- (5.2,0.3) node[above,midway] {$nx>y$};
\end{tikzpicture}
\end{center}

\begin{corollary}
\label{cor:racionales-densos-orden}
Para todo $\varepsilon>0$ existe $n\in\mathbb{N}$ tal que $1/n<\varepsilon$.
\end{corollary}

\section{Densidad de \texorpdfstring{$\mathbb{Q}$}{Q} y de los irracionales}

\begin{theorem}[Densidad de los racionales]
\label{teo:densidad-q}
Entre dos números reales distintos existe al menos un número racional. Equivalentemente, para todo $x\in\mathbb{R}$ y todo $\varepsilon>0$ existe $q\in\mathbb{Q}$ tal que $|x-q|<\varepsilon$.
\end{theorem}

\begin{proof}
Sea $x<y$. Por la propiedad arquimediana, existe $n\in\mathbb{N}$ tal que $n(y-x)>1$. Sea $m$ el menor entero tal que $m>nx$. Entonces $nx<m\leq nx+1<ny$, luego $x<m/n<y$.
\end{proof}

\begin{center}
\begin{tikzpicture}[scale=1.2]
  \draw[->] (-0.5,0) -- (5.5,0) node[right] {$\mathbb{R}$};
  \foreach \x in {0,1,2,3,4,5}
    \draw (\x,0.1) -- (\x,-0.1) node[below] {\x};
  \draw[thick,blue] (1.5,0.15) -- (1.5,-0.15) node[below] {$x$};
  \draw[thick,red] (3.5,0.15) -- (3.5,-0.15) node[below] {$y$};
  \draw[thick,green!60!black] (2.333,0.15) -- (2.333,-0.15) node[below=0.2cm] {$\frac{m}{n}$};
  \draw[<->,dashed] (1.5,0.4) -- (3.5,0.4) node[midway,above] {$y-x$};
\end{tikzpicture}
\end{center}

\begin{corollary}[Densidad de los irracionales]
\label{cor:densidad-irracionales}
Entre dos números reales distintos existe al menos un número irracional.
\end{corollary}
\begin{proof}
Si $x<y$, entonces $x-\sqrt{2}<y-\sqrt{2}$. Por densidad de $\mathbb{Q}$, existe $q\in\mathbb{Q}\cap(x-\sqrt{2},y-\sqrt{2})$. Luego $q+\sqrt{2}$ es irracional y pertenece a $(x,y)$.
\end{proof}

\begin{remark}
La densidad implica que cualquier número real puede aproximarse arbitrariamente por racionales e irracionales. Esta propiedad es esencial para definir funciones mediante prolongación por continuidad y para la teoría de la integración.
\end{remark}

\section{Intervalos encajados y completitud}

La completitud de $\mathbb{R}$ puede formularse de múltiples maneras equivalentes. Una de las más intuitivas es el principio de los intervalos encajados.

\begin{theorem}[Principio de los intervalos encajados de Cantor]
\label{teo:intervalos-encajados}
Sea $(I_n)$ una sucesión de intervalos cerrados y acotados $I_n=[a_n,b_n]$ tales que $I_{n+1}\subseteq I_n$ para todo $n$. Entonces $\bigcap_{n=1}^{\infty} I_n\neq\emptyset$.
\end{theorem}

\begin{proof}
La sucesión $(a_n)$ es creciente y acotada superiormente por $b_1$. Por completitud de Dedekind, existe $a=\sup a_n$. Análogamente, $(b_n)$ es decreciente y acotada inferiormente, luego existe $b=\inf b_n$. Como $a_n\leq b_n$ para todo $n$, se tiene $a\leq b$, y por tanto $[a,b]\subseteq\bigcap I_n$.
\end{proof}

\begin{center}
\begin{tikzpicture}[scale=0.9]
  \draw[->] (-0.5,0) -- (8.5,0);
  \draw[thick] (0,0.3) rectangle (8,-0.3) node[midway,above=0.4cm] {$I_1$};
  \draw[thick] (1,0.6) rectangle (6,-0.6) node[midway,above=0.4cm] {$I_2$};
  \draw[thick] (2,0.9) rectangle (4.5,-0.9) node[midway,above=0.4cm] {$I_3$};
  \draw[thick,red] (2.8,1.2) rectangle (3.6,-1.2) node[midway,above=0.5cm] {$\bigcap I_n$};
  \foreach \x/\n in {0/a_1,8/b_1,1/a_2,6/b_2,2/a_3,4.5/b_3}
    \draw (\x,0.15) -- (\x,-0.15) node[below] {$\n$};
\end{tikzpicture}
\end{center}

\begin{remark}
Si además la longitud $b_n-a_n\to 0$, entonces $\bigcap I_n$ consiste en un único punto. Esta versión fuerte se utiliza para demostrar la no numerabilidad de $\mathbb{R}$ y para construir funciones continuas nowhere-differentiable.
\end{remark}

\section{Completitud de Cauchy}

La completitud de $\mathbb{R}$ es la propiedad que distingue al análisis real del análisis sobre los racionales. Garantiza que los ``huecos'' de $\mathbb{Q}$ han sido colmatados.

\begin{theorem}[Completitud de $\mathbb{R}$]
\label{teo:completitud-r}
Toda sucesión de Cauchy en $\mathbb{R}$ converge a un límite en $\mathbb{R}$.
\end{theorem}

\begin{proof}[Bosquejo]
Sea $(x_n)$ de Cauchy. La sucesión está acotada, luego posee una subsucesión convergente (Bolzano--Weierstrass, que se prueba a partir de la completitud de Dedekind). Una sucesión de Cauchy con una subsucesión convergente es ella misma convergente.
\end{proof}

\begin{remark}
La completitud de $\mathbb{R}$ puede formularse de múltiples maneras equivalentes: completitud de Dedekind, completitud de Cauchy, principio de los intervalos encajados de Cantor, teorema de Bolzano--Weierstrass, y teorema de Heine--Borel para intervalos cerrados. Esta redundancia de formulaciones equivalentes es característica de las propiedades estructurales profundas.
\end{remark}

\section{La recta real extendida}

Es útil adjuntar a $\mathbb{R}$ dos símbolos $+\infty$ y $-\infty$ que representan los extremos ideales de la recta.

\begin{definition}
La \emph{recta real extendida} es $\overline{\mathbb{R}}=\mathbb{R}\cup\{-\infty,+\infty\}$. Se extiende el orden declarando $-\infty<x<+\infty$ para todo $x\in\mathbb{R}$. Las operaciones aritméticas con $\pm\infty$ se definen con las convenciones naturales, salvo las formas indeterminadas $\infty-\infty$, $0\cdot\infty$, $\infty/\infty$ y $0/0$.
\end{definition}

\begin{example}
En $\overline{\mathbb{R}}$, todo subconjunto tiene supremo e ínfimo. En particular, $\sup\emptyset=-\infty$ y $\inf\emptyset=+\infty$.
\end{example}

\section{Sucesiones reales elementales}

Aunque el estudio sistemático de las sucesiones corresponde al \cref{ch:sucesiones}, introducimos aquí algunos ejemplos paradigmáticos que ilustran las propiedades de $\mathbb{R}$.

\begin{example}[Serie armónica]
La serie $\sum_{n=1}^{\infty}\frac{1}{n}$ diverge. En efecto, agrupando términos
\[
  1+\frac{1}{2}+\Bigl(\frac{1}{3}+\frac{1}{4}\Bigr)+\Bigl(\frac{1}{5}+\cdots+\frac{1}{8}\Bigr)+\cdots
  \geq 1+\frac{1}{2}+\frac{1}{2}+\frac{1}{2}+\cdots
\]
se obtienen sumas parciales arbitrariamente grandes. Este hecho muestra que en $\mathbb{R}$, una sucesión creciente no siempre converge: necesita estar acotada.
\end{example}

\begin{example}[Serie armónica alternada]
La serie $\sum_{n=1}^{\infty}\frac{(-1)^{n+1}}{n}=1-\frac{1}{2}+\frac{1}{3}-\frac{1}{4}+\cdots$ converge a $\log 2$. Esta convergencia condicional (no absoluta) será demostrada en el \cref{ch:series} mediante el criterio de Leibniz.
\end{example}

\begin{example}[Sucesión de Fibonacci y razón áurea]
Definida por $F_1=F_2=1$ y $F_{n+2}=F_{n+1}+F_n$, la sucesión de Fibonacci satisface
\[
  \lim_{n\to\infty}\frac{F_{n+1}}{F_n}=\varphi=\frac{1+\sqrt{5}}{2}\approx 1.618\dots
\]
El número $\varphi$, llamado \emph{razón áurea}, es irracional. Este límite puede probarse resolviendo la ecuación característica $r^2=r+1$.
\end{example}

\section{Ejercicios}

\begin{exercise}
Demuestre que si $A\subseteq\mathbb{R}$ es no vacío y acotado inferiormente, entonces $\inf A$ existe.
\end{exercise}

\begin{exercise}
Pruebe que todo número real es límite de una sucesión de números racionales.
\end{exercise}

\begin{exercise}
Sea $A\subseteq\mathbb{R}$ no vacío, acotado superiormente, y tal que $\sup A\notin A$. Demuestre que existe una sucesión estrictamente creciente en $A$ que converge a $\sup A$.
\end{exercise}

\begin{exercise}
Use el principio de los intervalos encajados para probar que $[0,1]$ no es numerable.
\end{exercise}

\begin{exercise}
Demuestre que la sucesión $x_n=(1+\frac{1}{n})^n$ es creciente y acotada superiormente (por ejemplo, por 3), luego converge. El límite es el número $e$. (Sugerencia: use el binomio de Newton y compare término a término con $x_{n+1}$.)
\end{exercise}

\begin{exercise}
Sea $a_1=1$ y $a_{n+1}=\sqrt{2+a_n}$. Demuestre que $(a_n)$ converge y calcule su límite. Sugerencia: pruebe por inducción que $a_n\leq 2$ y que $(a_n)$ es creciente.
\end{exercise}

\begin{exercise}
Demuestre que entre dos números reales distintos hay infinitos racionales e infinitos irracionales.
\end{exercise}

\begin{exercise}
Sea $x\in\mathbb{R}$. Demuestre que existe una única sucesión $(d_n)$ con $d_n\in\{0,1,\dots,9\}$ tal que $x=\sum_{n=1}^{\infty}d_n10^{-n}$ y que la sucesión no termina en una cola de nueves. Esto justifica rigurosamente la expansión decimal.
\end{exercise}

\begin{exercise}
Demuestre que $\sqrt{2}+\sqrt{3}$ es irracional. Sugerencia: si fuera racional, entonces $(\sqrt{2}+\sqrt{3})^2$ también lo sería, lo que lleva a una contradicción.
\end{exercise}

\begin{exercise}
Sea $(x_n)$ una sucesión de Cauchy. Demuestre que $(x_n)$ está acotada. Use esto para dar una demostración directa (sin Bolzano--Weierstrass) de que toda sucesión de Cauchy en $\mathbb{R}$ converge. (Sugerencia: tome una subsucesión que converja a $\limsup x_n$.)
\end{exercise}

\begin{problem}
Un conjunto $A\subseteq\mathbb{R}$ se dice \emph{inductivo} si $1\in A$ y $x\in A\Rightarrow x+1\in A$. Sea $\mathbb{N}'=\bigcap\{A\subseteq\mathbb{R}:\,A\text{ inductivo}\}$. Demuestre que $\mathbb{N}'$ satisface el principio de inducción matemática y que coincide con el conjunto de los números naturales usuales.
\end{problem}

\begin{problem}
Demuestre que el conjunto de los números de Liouville (reales que pueden aproximarse ``demasiado bien'' por racionales) es denso en $\mathbb{R}$ pero de medida nula. Este resultado conecta la teoría de números con la teoría de la medida.
\end{problem}

\chapter{Números complejos}
\label{ap:numeros-complejos}
\index{Número complejo}\index{Conjugado}\index{Módulo}\index{Forma polar}\index{Raíz de la unidad}\index{Teorema fundamental del álgebra}\index{Exponencial compleja}

Este apéndice recoge los conceptos básicos de los números complejos que se emplean a lo largo del libro: la construcción formal del cuerpo $\mathbb{C}$, la representación binómica y polar, las raíces de la unidad, el teorema fundamental del álgebra y las funciones elementales de variable compleja. Los números complejos aparecen de manera natural en el análisis, la teoría espectral y las ecuaciones diferenciales.

\section{Construcción del cuerpo \texorpdfstring{$\mathbb{C}$}{C}}

Los números reales no resuelven todas las ecuaciones polinómicas: la ecuación $x^2+1=0$ no tiene solución en $\mathbb{R}$. Para remediarlo se amplía $\mathbb{R}$ añadiendo una raíz de $-1$.

\begin{definition}
\label{def:complejo}
El conjunto $\mathbb{C}$ de los \emph{números complejos} es $\mathbb{R}^2$ dotado de las operaciones
\begin{equation}\label{eq:suma-compleja}
  (a,b)+(c,d)=(a+c,b+d),
\end{equation}
\begin{equation}\label{eq:producto-complejo}
  (a,b)\cdot(c,d)=(ac-bd,ad+bc).
\end{equation}
\end{definition}

\begin{theorem}[$\mathbb{C}$ es un cuerpo]
\label{teo:c-cuerpo}
$(\mathbb{C},+,\cdot)$ es un cuerpo. El elemento neutro de la suma es $(0,0)$, el del producto es $(1,0)$, y el inverso de $z=(a,b)\neq 0$ es
\[
  z^{-1}=\Bigl(\frac{a}{a^2+b^2},\frac{-b}{a^2+b^2}\Bigr).
\]
\end{theorem}

Identificando $(a,0)$ con $a\in\mathbb{R}$, el cuerpo $\mathbb{R}$ se sumerge en $\mathbb{C}$. Escribiendo $i=(0,1)$, se tiene $i^2=(-1,0)=-1$, y todo número complejo se expresa de manera única como
\[
  z=a+bi,\qquad a,b\in\mathbb{R},
\]
la forma \emph{binómica} o \emph{cartesiana} de $z$. El número real $a=\operatorname{Re}(z)$ es la \emph{parte real} y $b=\operatorname{Im}(z)$ la \emph{parte imaginaria}.

\section{Conjugación y módulo}

\begin{definition}
\label{def:conjugado-modulo}
Para $z=a+bi\in\mathbb{C}$:
\begin{itemize}
  \item El \emph{conjugado} de $z$ es $\overline{z}=a-bi$.
  \item El \emph{módulo} de $z$ es $|z|=\sqrt{a^2+b^2}$.
\end{itemize}
\end{definition}

\begin{theorem}
\label{teo:prop-conjugacion}
Para cualesquiera $z,w\in\mathbb{C}$:
\begin{enumerate}
  \item[(i)] $\overline{z+w}=\overline{z}+\overline{w}$ y $\overline{z\,w}=\overline{z}\,\overline{w}$;
  \item[(ii)] $z\overline{z}=|z|^2$, de modo que $z^{-1}=\overline{z}/|z|^2$ si $z\neq 0$;
  \item[(iii)] $|zw|=|z|\,|w|$ y $|z+w|\leq|z|+|w|$ (desigualdad triangular);
  \item[(iv)] $\operatorname{Re}(z)=\frac{z+\overline{z}}{2}$ e $\operatorname{Im}(z)=\frac{z-\overline{z}}{2i}$.
\end{enumerate}
\end{theorem}

\begin{proof}
La parte (i) se comprueba directamente con las definiciones \eqref{eq:suma-compleja} y \eqref{eq:producto-complejo}. Para (iii), nótese que $|zw|^2=zw\,\overline{zw}=z\overline{z}\,w\overline{w}=|z|^2|w|^2$. La desigualdad triangular sigue de $|z+w|^2=|z|^2+|w|^2+2\operatorname{Re}(z\overline{w})\leq(|z|+|w|)^2$.
\end{proof}

\section{Forma polar}

\begin{definition}
\label{def:forma-polar}
Todo número complejo $z\neq 0$ admite la representación \emph{polar}
\[
  z=r(\cos\theta+i\sin\theta),
\]
donde $r=|z|$ y $\theta$ es un \emph{argumento} de $z$, es decir, un ángulo tal que $\cos\theta=a/r$ y $\sin\theta=b/r$.
\end{definition}

El argumento no es único: si $\theta$ es un argumento, también lo es $\theta+2k\pi$ para todo $k\in\mathbb{Z}$. Se denota $\operatorname{Arg}(z)$ al argumento principal, tomado en $(-\pi,\pi]$.

\begin{theorem}[Fórmula de De Moivre]
\label{teo:demoivre}
Para $z=r(\cos\theta+i\sin\theta)$ y $n\in\mathbb{N}$,
\[
  z^n=r^n(\cos(n\theta)+i\sin(n\theta)).
\]
\end{theorem}

En particular, la fórmula de Euler $e^{i\theta}=\cos\theta+i\sin\theta$ permite escribir $z=re^{i\theta}$ y recuperar las leyes del producto $|zw|=|z|\,|w|$ y $\operatorname{Arg}(zw)=\operatorname{Arg}(z)+\operatorname{Arg}(w)$.

\section{Raíces de la unidad}

\begin{theorem}
\label{teo:raices-unidad}
Para cada $n\in\mathbb{N}$, la ecuación $z^n=1$ tiene exactamente $n$ soluciones en $\mathbb{C}$, a saber
\[
  \zeta_k=e^{2k\pi i/n}=\cos\frac{2k\pi}{n}+i\sin\frac{2k\pi}{n},\qquad k=0,1,\dots,n-1.
\]
\end{theorem}

Estas $n$ raíces forman un grupo cíclico bajo el producto, cuyos generadores se denominan \emph{raíces primitivas} de la unidad.

\begin{example}
Las raíces cúbicas de la unidad son
\[
  1,\qquad \omega=e^{2\pi i/3}=-\frac12+\frac{\sqrt3}{2}i,\qquad \omega^2=e^{4\pi i/3}=-\frac12-\frac{\sqrt3}{2}i,
\]
y satisfacen $1+\omega+\omega^2=0$.
\end{example}

\section{El teorema fundamental del álgebra}

\begin{theorem}[Fundamental del álgebra]
\label{teo:fundamental-algebra}
Todo polinomio $p(z)=a_nz^n+a_{n-1}z^{n-1}+\cdots+a_0$ de grado $n\geq 1$ con coeficientes en $\mathbb{C}$ posee al menos una raíz compleja. En consecuencia, $p$ factoriza como
\[
  p(z)=a_n\prod_{k=1}^{n}(z-z_k),
\]
donde $z_1,\dots,z_n\in\mathbb{C}$ son sus raíces, contadas con multiplicidad.
\end{theorem}

\begin{proof}[Bosquejo]
Como $|p(z)|\to\infty$ cuando $|z|\to\infty$, la función $|p|$ alcanza su mínimo en algún punto $z_0$. Si $p(z_0)\neq 0$, un argumento de decrecimiento por el término de menor grado de $p(z_0+h)$ muestra que en la vecindad de $z_0$ el módulo puede reducirse, contradicción. Luego $p(z_0)=0$. La factorización se obtiene por división sucesiva.
\end{proof}

\begin{corollary}
Todo polinomio de grado $n$ con coeficientes reales factoriza en $\mathbb{C}$ como producto de factores lineales; sus raíces no reales aparecen por pares conjugados.
\end{corollary}

\section{Funciones de variable compleja}

\begin{definition}
\label{def:exp-compleja}
La \emph{exponencial compleja} se define por
\[
  e^{z}=e^{a+bi}=e^{a}(\cos b+i\sin b).
\]
Las funciones \emph{seno} y \emph{coseno} complejos se definen por
\[
  \sin z=\frac{e^{iz}-e^{-iz}}{2i},\qquad
  \cos z=\frac{e^{iz}+e^{-iz}}{2}.
\]
\end{definition}

La exponencial compleja satisface $e^{z+w}=e^{z}e^{w}$ y es periódica de período $2\pi i$. Las funciones trigonométricas complejas heredan las identidades clásicas, y el logaritmo complejo se define como inversa local de la exponencial.

\section{Ejercicios}

\begin{exercise}
Demuestre que el inverso de $z=a+bi\neq 0$ dado en el \cref{teo:c-cuerpo} es correcto verificando $z z^{-1}=1$.
\end{exercise}

\begin{exercise}
Calcule $(1+i)^n$ para $n=1,2,\dots,8$ usando la forma polar. ¿Qué patrón se observa?
\end{exercise}

\begin{exercise}
Resuelva la ecuación $z^2-2z+5=0$ y compruebe que las raíces son conjugadas.
\end{exercise}

\begin{exercise}
Demuestre que $\sin(2z)=2\sin z\cos z$ para $z\in\mathbb{C}$ usando las definiciones anteriores.
\end{exercise}

\begin{exercise}
Pruebe que $|z+w|^2+|z-w|^2=2(|z|^2+|w|^2)$ (identidad del paralelogramo).
\end{exercise}

\begin{exercise}
Sea $z=re^{i\theta}$. Exprese $z^n+\overline{z}^n$ en términos de $r$ y $\cos(n\theta)$.
\end{exercise}

\begin{exercise}
Demuestre que si $z$ es raíz del polinomio $p$ con coeficientes reales, entonces $\overline{z}$ también lo es.
\end{exercise}

\chapter{Álgebra lineal: espacios vectoriales y valores propios}
\label{ap:algebra-lineal}
\index{Espacio vectorial}\index{Norma}\index{Producto interno}\index{Valor propio}\index{Teorema espectral}\index{Exponencial de matrices}

Este apéndice recoge los conceptos básicos del álgebra lineal finito-dimensional que se utilizan a lo largo del texto: espacios vectoriales, normas, productos internos, valores propios y la exponencial de matrices. Su lectura permite que el libro sea autocontenido en todo lo referente a herramientas matriciales.

\section{Espacios vectoriales y normas}

\begin{definition}
\label{def:espacio-vectorial-ap}
Un \emph{espacio vectorial} sobre $\mathbb{R}$ es un conjunto $V$ equipado con dos operaciones ---suma $+\colon V\times V\to V$ y producto por escalares $\cdot\colon\mathbb{R}\times V\to V$--- que satisfacen las propiedades usuales de asociatividad, conmutatividad, elemento neutro, inversos aditivos, y distributividad del producto escalar respecto de la suma en $V$ y en $\mathbb{R}$.
\end{definition}

Los ejemplos más importantes para este libro son $\mathbb{R}^n$, el espacio $C([a,b])$ de funciones continuas, y los espacios de matrices $M_{m\times n}(\mathbb{R})$.

\begin{definition}
\label{def:norma-ap}
Una \emph{norma} sobre un espacio vectorial $V$ es una función $\|\cdot\|\colon V\to\mathbb{R}$ tal que:
\begin{enumerate}
  \item[(i)] $\|x\|\geq 0$, con igualdad si y solo si $x=0$;
  \item[(ii)] $\|\lambda x\|=|\lambda|\,\|x\|$ para todo $\lambda\in\mathbb{R}$, $x\in V$;
  \item[(iii)] $\|x+y\|\leq\|x\|+\|y\|$ (desigualdad triangular).
\end{enumerate}
\end{definition}

\begin{example}[Normas en $\mathbb{R}^n$]
Las normas más usuales son:
\[
  \|x\|_1=\sum_{i=1}^{n}|x_i|,\qquad
  \|x\|_2=\Bigl(\sum_{i=1}^{n}x_i^2\Bigr)^{1/2},\qquad
  \|x\|_\infty=\max_{1\leq i\leq n}|x_i|.
\]
\end{example}

\begin{theorem}[Equivalencia de normas]
\label{teo:equivalencia-normas-ap}
En $\mathbb{R}^n$ (y en cualquier espacio de dimensión finita), todas las normas son equivalentes: si $\|\cdot\|_a$ y $\|\cdot\|_b$ son normas, existen constantes $c,C>0$ tales que $c\|x\|_a\leq\|x\|_b\leq C\|x\|_a$ para todo $x$.
\end{theorem}

\section{Producto interno}

\begin{definition}
\label{def:producto-interno-ap}
Un \emph{producto interno} sobre un espacio vectorial real $V$ es una aplicación $\langle\cdot,\cdot\rangle\colon V\times V\to\mathbb{R}$ bilineal, simétrica y definida positiva. Es decir:
\begin{enumerate}
  \item[(i)] $\langle x,x\rangle\geq 0$ y $\langle x,x\rangle=0\Leftrightarrow x=0$;
  \item[(ii)] $\langle x,y\rangle=\langle y,x\rangle$;
  \item[(iii)] $\langle x,\lambda y+\mu z\rangle=\lambda\langle x,y\rangle+\mu\langle x,z\rangle$.
\end{enumerate}
La \emph{norma inducida} es $\|x\|=\sqrt{\langle x,x\rangle}$.
\end{definition}

\begin{theorem}[Cauchy--Schwarz]
\label{teo:cauchy-schwarz-ap}
Para cualesquiera $x,y\in V$,
\[
  |\langle x,y\rangle|\leq\|x\|\,\|y\|.
\]
La igualdad se da si y solo si $x$ e $y$ son linealmente dependientes.
\end{theorem}
\begin{proof}
Si $y=0$ el resultado es trivial. Para $y\neq 0$, sea $\lambda=\langle x,y\rangle/\langle y,y\rangle$. Entonces
\[
  0\leq\langle x-\lambda y,x-\lambda y\rangle=\|x\|^2-\frac{\langle x,y\rangle^2}{\|y\|^2},
\]
de donde se deduce el teorema.
\end{proof}

\begin{corollary}
La norma inducida por un producto interno satisface la desigualdad triangular.
\end{corollary}

\section{Valores propios y vectores propios}

\begin{definition}
\label{def:valor-propio-ap}
Sea $A\in M_n(\mathbb{R})$. Un escalar $\lambda\in\mathbb{C}$ es un \emph{valor propio} de $A$ si existe $v\neq 0$ tal que $Av=\lambda v$. El vector $v$ se denomina \emph{vector propio} asociado a $\lambda$. El polinomio $p_A(\lambda)=\det(A-\lambda I)$ se llama \emph{polinomio característico} de $A$; sus raíces son exactamente los valores propios.
\end{definition}

\begin{theorem}
Los valores propios de una matriz simétrica real $A=A^{\top}$ son todos reales. Más aún, existe una base ortonormal de $\mathbb{R}^n$ formada por vectores propios de $A$.
\end{theorem}
\begin{proof}[Bosquejo]
Si $Av=\lambda v$ con $v\neq 0$, tomando conjugados y usando $A=\overline{A}$ se obtiene $\lambda=\overline{\lambda}$. La existencia de la base ortonormal se prueba por inducción sobre $n$.
\end{proof}

\begin{theorem}[Descomposición espectral]
\label{teo:espectral-ap}
Si $A$ es simétrica real, entonces $A=QDQ^{\top}$, donde $Q$ es una matriz ortogonal ($Q^{\top}Q=I$) y $D$ es diagonal con los valores propios de $A$ en la diagonal.
\end{theorem}

\begin{definition}
Se dice que $A$ es \emph{definida positiva} si $\langle Ax,x\rangle>0$ para todo $x\neq 0$. Equivalentemente, todos sus valores propios son estrictamente positivos.
\end{definition}

\section{Exponencial de matrices}

\begin{definition}
\label{def:exponencial-matriz-ap}
Para $A\in M_n(\mathbb{R})$, la \emph{exponencial de matriz} se define mediante la serie
\[
  e^{A}=\sum_{k=0}^{\infty}\frac{A^k}{k!},
\]
que converge absolutamente en la norma operador $\|A\|=\sup_{\|x\|=1}\|Ax\|$.
\end{definition}

\begin{theorem}
\label{teo:prop-exponencial-ap}
\begin{enumerate}
  \item[(i)] $\frac{d}{dt}e^{tA}=A\,e^{tA}=e^{tA}A$.
  \item[(ii)] Si $A$ y $B$ conmutan ($AB=BA$), entonces $e^{A+B}=e^{A}e^{B}$.
  \item[(iii)] $e^{A}$ es invertible y $(e^{A})^{-1}=e^{-A}$.
  \item[(iv)] Si $A$ es simétrica, $\|e^{A}\|=e^{\lambda_{\max}}$, donde $\lambda_{\max}$ es el mayor valor propio de $A$.
\end{enumerate}
\end{theorem}

\begin{example}
Si $A$ es diagonalizable, $A=PDP^{-1}$, entonces $e^{A}=Pe^{D}P^{-1}$, donde $e^{D}$ es la matriz diagonal con $e^{\lambda_i}$ en la posición $i$-ésima.
\end{example}

\section{Ejercicios}

\begin{exercise}
Demuestre que si $A$ es ortogonal ($A^{\top}A=I$), todos sus valores propios complejos tienen módulo $1$.
\end{exercise}

\begin{exercise}
Calcule $e^{tA}$ para $A=\begin{pmatrix}0&-1\\1&0\end{pmatrix}$. Verifique que $e^{tA}$ es una matriz de rotación.
\end{exercise}

\begin{exercise}
Sea $A$ simétrica. Demuestre que $A$ es definida positiva si y solo si todos sus menores principales son positivos (criterio de Sylvester).
\end{exercise}

\begin{exercise}
Demuestre que toda norma en $\mathbb{R}^n$ es equivalente a la norma euclidiana. Sugerencia: use compacidad de la esfera unidad.
\end{exercise}

\begin{exercise}
Sea $A$ una matriz $n\times n$ con $n$ valores propios distintos. Demuestre que $A$ es diagonalizable.
\end{exercise}

\begin{exercise}
Use la descomposición espectral para probar que si $A$ es simétrica, entonces $\|A\|_2=\max_i|\lambda_i|$, donde $\|\cdot\|_2$ es la norma operador inducida por la norma euclidiana.
\end{exercise}


% --- Índice de palabras ---
\cleardoublepage
\phantomsection
\addcontentsline{toc}{chapter}{Índice de palabras}
\printindex

\backmatter

% --- Referencias (bibliografía) ---
\cleardoublepage
\phantomsection
\nocite{*}
\printbibliography[heading=bibintoc, title=Referencias]

\end{document}
